
\documentclass[11pt]{scrartcl}
\usepackage[sexy]{evan} % comot, baca kalo interested
\usepackage{graphicx} % Required for inserting images

\usepackage[indonesian]{babel}
\usepackage[latin1]{inputenc}
\usepackage{amsmath,amsfonts,amssymb,graphicx,fancyhdr,lastpage, multirow, multicol, hyperref}
\usepackage{enumitem}
\usepackage{enumerate}
\usepackage[table]{xcolor}
\usepackage{tikz}
\usepackage{setspace}
\usepackage{tcolorbox}
\usepackage{pifont}
\linespread{1.3}
\usepackage[left=1.50cm, right=1.50cm, top=1.50cm, bottom=1.50cm]{geometry}
\usepackage{pgfplots,tikz}
\pgfplotsset{compat=1.15}
\usetikzlibrary{arrows}
\usepackage{mathrsfs}
\usepackage{dashbox}


\author{Raymond Christopher Tanto}
\title{\textbf{Big Beautiful IMO Shortlist} }
\subtitle{Version 2.1}
\date{Last updated: 31 July 2026, 12:51 GMT}

\begin{document}
\author{Raymond Christopher Tanto}
\maketitle

\author{ \hyperref[Back]{\textbf{Click here to go back to directory}} }
\date{\textbf{Raymond Christopher Tanto}}

\section*{Preamble}
I believe that doing IMO Shortlist problems is a moral good for all who dream of the IMO.\\
When I was a beginner, doing most of the level 1 and 2 problems of the IMO Shortlist was really\\ effective to start my grind. They encompass nearly all ideas needed to start solving IMO-level problems.\\
I made this document to make it more accessible to do ISL problems, because it is in one file.\\
The beautiful LaTeX template used in this document is evan.sty by Evan Chen, who will undoubtedly go down as the greatest mathematical olympiad coach of all time for his significant contribution to the people. \\
Yes, this is a word play on the Big Beautiful Bill.
\section*{Directory} \label{Back} 
Pages 2 and 3 are directories to each ISL number, classified by subject.
\section*{Legend}
\begin{itemize}[leftmargin=*]
    \item Click on $\href{https://artofproblemsolving.com/community/c3223_imo_shortlist}{\textbf{Pink}}$ to go to the AoPS Threads.
    \item Click the \hyperref[Back]{\textbf{"Click here to go back to directory"}} to go back to page 1.
    \item $\boxed{\text{Boxed}}$ problems are selected as IMO Problems.
    \item $\dbox{\text{Dotted Boxed}}$ problems were selected for IMO 2015 day 2, but had to be scrapped due to a leak.
\end{itemize}

\section*{Contact}
As I added the problems manually by hand without any automation, I expect there would be a lot of mistakes in this hundred page document. Please let me know if there are any. Here are my contacts and profile:
\begin{align*}
    \text{Instagram : } &\href{https://www.instagram.com/raymond_loves_math/}{\text{@raymond\_loves\_math}}\\
    \text{Discord: } &\href{discord.com}{\text{raymond.math.07}} \\ 
  %  \text{Telegram: } &\href{telegram.org}{\text{@RaymondChristopherTanto}} \\ 
\text{AoPS : } &\href{https://artofproblemsolving.com/community/user/886353}{\text{GreenTea2593}}\\
%\text{IMO Website: } &\href{https://www.imo-official.org/results/contestant/34572/}{\text{Raymond Christopher Tanto}}\\
\text{Moresults : } &\href{https://moresults.org/contestants/rct-2}{\text{Raymond Christopher Tanto}} \\ 
\end{align*}


\date{\textbf{Big Beautiful IMO Shortlist 2.1}}
\title{Directory}

\begin{center}
\huge
    \textbf{Algebra Problems}
\end{center}
\normalsize
Click on \textbf{\textcolor{RoyalBlue}{blue}} to teleport!
\begin{itemize}
\Large 
\item[\ding{227}] \textbf{\hyperref[A1]{A1 Problems}}
\item[\ding{227}] \textbf{\hyperref[A2]{A2 Problems}}
\item[\ding{227}] \textbf{\hyperref[A3]{A3 Problems}}
\item[\ding{227}] \textbf{\hyperref[A4]{A4 Problems}}
\item[\ding{227}] \textbf{\hyperref[A5]{A5 Problems}}
\item[\ding{227}] \textbf{\hyperref[A6]{A6 Problems}}
\item[\ding{227}] \textbf{\hyperref[A7]{A7 Problems}}
\item[\ding{227}] \textbf{\hyperref[A8]{A8 Problems}}
\item[\ding{227}] \textbf{\hyperref[A9]{A9 Problems}}
\end{itemize}
\begin{center}
\huge \textbf{Combinatorics Problems}
\end{center}
\normalsize
Click on \textbf{\textcolor{RoyalBlue}{blue}} to teleport!
\begin{itemize}
\Large 
\item[\ding{227}] \textbf{\hyperref[C1]{C1 Problems}}
\item[\ding{227}] \textbf{\hyperref[C2]{C2 Problems}}
\item[\ding{227}] \textbf{\hyperref[C3]{C3 Problems}}
\item[\ding{227}] \textbf{\hyperref[C4]{C4 Problems}}
\item[\ding{227}] \textbf{\hyperref[C5]{C5 Problems}}
\item[\ding{227}] \textbf{\hyperref[C6]{C6 Problems}}
\item[\ding{227}] \textbf{\hyperref[C7]{C7 Problems}}
\item[\ding{227}] \textbf{\hyperref[C8]{C8 Problems}}
\item[\ding{227}] \textbf{\hyperref[C9]{C9 Problems}}
\newpage
\end{itemize}

 

\newpage


\title{Directory}

\begin{center}
\huge \textbf{Geometry Problems}
\end{center} 
\normalsize
Click on \textbf{\textcolor{RoyalBlue}{blue}} to teleport!
\begin{itemize}
\Large 
\item[\ding{227}] \textbf{\hyperref[G1]{G1 Problems}}
\item[\ding{227}] \textbf{\hyperref[G2]{G2 Problems}}
\item[\ding{227}] \textbf{\hyperref[G3]{G3 Problems}}
\item[\ding{227}] \textbf{\hyperref[G4]{G4 Problems}}
\item[\ding{227}] \textbf{\hyperref[G5]{G5 Problems}}
\item[\ding{227}] \textbf{\hyperref[G6]{G6 Problems}}
\item[\ding{227}] \textbf{\hyperref[G7]{G7 Problems}}
\item[\ding{227}] \textbf{\hyperref[G8]{G8 Problems}}
\item[\ding{227}] \textbf{\hyperref[G9]{G9 Problems}}
\item[\ding{227}] \textbf{\hyperref[G10]{G10 Problems}}

\end{itemize} 
\begin{center}
\huge \textbf{Number Theory Problems}
\end{center}
\normalsize
Click on \textbf{\textcolor{RoyalBlue}{blue}} to teleport!
\begin{itemize}
\Large 
\item[\ding{227}] \textbf{\hyperref[N1]{N1 Problems}}
\item[\ding{227}] \textbf{\hyperref[N2]{N2 Problems}}
\item[\ding{227}] \textbf{\hyperref[N3]{N3 Problems}}
\item[\ding{227}] \textbf{\hyperref[N4]{N4 Problems}}
\item[\ding{227}] \textbf{\hyperref[N5]{N5 Problems}}
\item[\ding{227}] \textbf{\hyperref[N6]{N6 Problems}}
\item[\ding{227}] \textbf{\hyperref[N7]{N7 Problems}}
\item[\ding{227}] \textbf{\hyperref[N8]{N8 Problems}}
\item[\ding{227}] \textbf{\hyperref[N9]{N9 Problems}}

\end{itemize}
\newpage

\section*{IMO Shortlist Booklet Links}
\noindent Here are the links for the official IMO Shortlist booklets. They contain the original wording of the problem and the official solutions provided by the IMO Problem Selection Committee. 
\title{Shortlist Booklet Link}
\begin{itemize}[leftmargin=*]
 \item[\ding{227}] \href{https://www.imo-official.org/assets/documents/problems/2026/IMO2026SL.pdf}{\text{2026 IMO Shortlist Booklet}} (Available after the conclusion of IMO 2027)
 
    \item[\ding{227}] \href{https://www.imo-official.org/assets/documents/problems/2025/IMO2025SL.pdf}{\text{2025 IMO Shortlist Booklet}}
\item[\ding{227}] \href{https://www.imo-official.org/assets/documents/problems/2024/IMO2024SL.pdf}{\text{2024 IMO Shortlist Booklet}}
\item[\ding{227}] \href{https://www.imo-official.org/assets/documents/problems/2023/IMO2023SL.pdf}{\text{2023 IMO Shortlist Booklet}}
\item[\ding{227}] \href{https://www.imo-official.org/assets/documents/problems/2022/IMO2022SL.pdf}{\text{2022 IMO Shortlist Booklet}}
\item[\ding{227}] \href{https://www.imo-official.org/assets/documents/problems/2021/IMO2021SL.pdf}{\text{2021 IMO Shortlist Booklet}}
\item[\ding{227}] \href{https://www.imo-official.org/assets/documents/problems/2020/IMO2020SL.pdf}{\text{2020 IMO Shortlist Booklet}}
\item[\ding{227}] \href{https://www.imo-official.org/assets/documents/problems/2019/IMO2019SL.pdf}{\text{2019 IMO Shortlist Booklet}}
\item[\ding{227}] \href{https://www.imo-official.org/assets/documents/problems/2018/IMO2018SL.pdf}{\text{2018 IMO Shortlist Booklet}}
\item[\ding{227}] \href{https://www.imo-official.org/assets/documents/problems/2017/IMO2017SL.pdf}{\text{2017 IMO Shortlist Booklet}}
\item[\ding{227}] \href{https://www.imo-official.org/assets/documents/problems/2016/IMO2016SL.pdf}{\text{2016 IMO Shortlist Booklet}}
\item[\ding{227}] \href{https://www.imo-official.org/assets/documents/problems/2015/IMO2015SL.pdf}{\text{2015 IMO Shortlist Booklet}}
\item[\ding{227}] \href{https://www.imo-official.org/assets/documents/problems/2014/IMO2014SL.pdf}{\text{2014 IMO Shortlist Booklet}}
\item[\ding{227}] \href{https://www.imo-official.org/assets/documents/problems/2013/IMO2013SL.pdf}{\text{2013 IMO Shortlist Booklet}}
\item[\ding{227}] \href{https://www.imo-official.org/assets/documents/problems/2012/IMO2012SL.pdf}{\text{2012 IMO Shortlist Booklet}}
\item[\ding{227}] \href{https://www.imo-official.org/assets/documents/problems/2011/IMO2011SL.pdf}{\text{2011 IMO Shortlist Booklet}}
\item[\ding{227}] \href{https://www.imo-official.org/assets/documents/problems/2010/IMO2010SL.pdf}{\text{2010 IMO Shortlist Booklet}}
\item[\ding{227}] \href{https://www.imo-official.org/assets/documents/problems/2009/IMO2009SL.pdf}{\text{2009 IMO Shortlist Booklet}}
\item[\ding{227}] \href{https://www.imo-official.org/assets/documents/problems/2008/IMO2008SL.pdf}{\text{2008 IMO Shortlist Booklet}}
\item[\ding{227}] \href{https://www.imo-official.org/assets/documents/problems/2007/IMO2007SL.pdf}{\text{2007 IMO Shortlist Booklet}}
\item[\ding{227}] \href{https://www.imo-official.org/assets/documents/problems/2006/IMO2006SL.pdf}{\text{2006 IMO Shortlist Booklet}}


\end{itemize}

\newpage


\section*{AoPS Collection Links}
\noindent Here are the links for the AoPS collections. They also have pre-2006 problems.
\title{AoPS Collection Link} 

\begin{itemize}[leftmargin=*]
\item[\ding{227}] \href{https://artofproblemsolving.com/community/c3223_imo_shortlist}{\text{All Years IMO Shortlist AoPS Collection}}
 
    \item[\ding{227}] \href{https://artofproblemsolving.com/community/c4971703_2025_imo_shortlist}{\text{2025 IMO Shortlist AoPS Collection}}
    
\item[\ding{227}] \href{https://artofproblemsolving.com/community/c3376654_2024_isl}{\text{2024 IMO Shortlist AoPS Collection}}

\item[\ding{227}] \href{https://artofproblemsolving.com/community/c3922196_2023_isl}{\text{2023 IMO Shortlist AoPS Collection}}

\item[\ding{227}] \href{https://artofproblemsolving.com/community/c3381907_2022_imo_shortlist}{\text{2022 IMO Shortlist AoPS Collection}}

\item[\ding{227}] \href{https://artofproblemsolving.com/community/c3080392_2021_imo_shortlist}{\text{2021 IMO Shortlist AoPS Collection}}

\item[\ding{227}] \href{https://artofproblemsolving.com/community/c2409205_2020_imo_shortlist}{\text{2020 IMO Shortlist AoPS Collection}}

\item[\ding{227}] \href{https://artofproblemsolving.com/community/c1308102_2019_imo_shortlist}{\text{2019 IMO Shortlist AoPS Collection}}

\item[\ding{227}] \href{https://artofproblemsolving.com/community/c915701_2018_imo_shortlist}{\text{2018 IMO Shortlist AoPS Collection}}

\item[\ding{227}] \href{https://artofproblemsolving.com/community/c686986_2017_imo_shortlist}{\text{2017 IMO Shortlist AoPS Collection}}

\item[\ding{227}] \href{https://artofproblemsolving.com/community/c482986_2016_imo_shortlist}{\text{2016 IMO Shortlist AoPS Collection}}

\item[\ding{227}] \href{https://artofproblemsolving.com/community/c111148_2015_imo_shortlist}{\text{2015 IMO Shortlist AoPS Collection}}

\item[\ding{227}] \href{https://artofproblemsolving.com/community/c107000_2014_imo_shortlist}{\text{2014 IMO Shortlist AoPS Collection}}

\item[\ding{227}] \href{https://artofproblemsolving.com/community/c3964_2013_imo_shortlist}{\text{2013 IMO Shortlist AoPS Collection}}

\item[\ding{227}] \href{https://artofproblemsolving.com/community/c3963_2012_imo_shortlist}{\text{2012 IMO Shortlist AoPS Collection}}

\item[\ding{227}] \href{https://artofproblemsolving.com/community/c3962_2011_imo_shortlist}{\text{2011 IMO Shortlist AoPS Collection}}

\item[\ding{227}] \href{https://artofproblemsolving.com/community/c3961_2010_imo_shortlist}{\text{2010 IMO Shortlist AoPS Collection}}

\item[\ding{227}] \href{https://artofproblemsolving.com/community/c3960_2009_imo_shortlist}{\text{2009 IMO Shortlist AoPS Collection}}

\item[\ding{227}] \href{https://artofproblemsolving.com/community/c3959_2008_imo_shortlist}{\text{2008 IMO Shortlist AoPS Collection}}

\item[\ding{227}] \href{https://artofproblemsolving.com/community/c3958_2007_imo_shortlist}{\text{2007 IMO Shortlist AoPS Collection}}

\item[\ding{227}] \href{https://artofproblemsolving.com/community/c3957_2006_imo_shortlist}{\text{2006 IMO Shortlist AoPS Collection}}



\end{itemize}

\newpage
\title{Algebra IMO Shortlist Problems}
\normalsize
\section*{A1 Problems} \label{A1} 
\begin{itemize}[leftmargin=*]


\item \href{https://www.imo-official.org/assets/documents/problems/2026/IMO2026SL.pdf}{\textbf{2026 A1}} TO BE REVEALED!

\item \href{https://artofproblemsolving.com/community/c6h3866863p38676959}{\textbf{2025 A1}} \textit{Quadratic solitaire} is a single-player game. To start the game, the player chooses two distinct nonzero integers $a$ and $b$ and writes the equation $x^2+ax+b=0$ on a blackboard. On each turn, if the equation currently written on the blackboard has two distinct nonzero integer solutions $x=u$ and $x=v$, then the player erases the equation and replaces it by one of the equations $x^2+ux+v=0$ and $x^2+vx+u=0$ of their choosing. Otherwise, the game ends. 

Determine all initial choices of $a$ and $b$ such that quadratic solitaire can be played forever.
\item \boxed{\href{https://artofproblemsolving.com/community/c6h3358923p31205921}{{\textbf{2024 A1}}}} Determine all real numbers $\alpha$ such that, for every positive integer $n,$ the integer
$$\lfloor\alpha\rfloor +\lfloor 2\alpha\rfloor +\dots +\lfloor n\alpha\rfloor$$is a multiple of $n.$ (Note that $\lfloor z\rfloor$ denotes the greatest integer less than or equal to $z$.)
\item \href{https://artofproblemsolving.com/community/c6h3359728p31218381}{\textbf{2023 A1}} Professor Oak is feeding his $100$ Pokemon. Each Pokemon has a bowl whose capacity is a positive real number of kilograms. These capacities are known to Professor Oak. The total capacity of all the bowls is $100$ kilograms. Professor Oak distributes $100$ kilograms of food in such a way that each Pokemon receives a non-negative integer number of kilograms of food (which may be larger than the capacity of the bowl). The dissatisfaction level of a Pokemon who received $N$ kilograms of food and whose bowl has a capacity of $C$ kilograms is equal to $\lvert N-C\rvert$.

Find the smallest real number $D$ such that, regardless of the capacities of the bowls, Professor Oak can distribute food in a way that the sum of the dissatisfaction levels over all the $100$ Pokemon is at most $D$.

\item  \href{https://artofproblemsolving.com/community/c6h3107323p28104243}{\textbf{2022 A1}} Let $(a_n)_{n\geqslant  1}$ be a sequence of positive real numbers with the property that
$$(a_{n+1})^2 + a_na_{n+2} \leqslant a_n + a_{n+2}$$for all positive integers $n$. Show that $a_{2022}\leqslant 1$. 
\item \href{https://artofproblemsolving.com/community/c6h2882570p25627665}{\textbf{2021 A1}}	Let $n$ be a positive integer. Given is a subset $A$ of $\{0,1,\dots,5^n\}$ with $4n+2$ elements. Prove that there exist three elements $a<b<c$ from $A$ such that $c+2a>3b$.
\item \href{https://artofproblemsolving.com/community/c6h2625913p22698446}{\textbf{2020 A1}} \textit{Version 1.} Let $n$ be a positive integer, and set $N=2^{n}$. Determine the smallest real number $a_{n}$ such that, for all real $x$,
\[
\sqrt[N]{\frac{x^{2 N}+1}{2}} \leqslant a_{n}(x-1)^{2}+x .
\] 
\textit{Version 2.} For every positive integer $N$, determine the smallest real number $b_{N}$ such that, for all real $x$,
\[
\sqrt[N]{\frac{x^{2 N}+1}{2}} \leqslant b_{N}(x-1)^{2}+x .
\]

\item \boxed{\href{https://artofproblemsolving.com/community/c6h1876068p12744859}{{\textbf{2019 A1}}}}	Let $\mathbb{Z}$ be the set of integers. Determine all functions $f: \mathbb{Z} \rightarrow \mathbb{Z}$ such that, for all integers $a$ and $b$,$$f(2a)+2f(b)=f(f(a+b)).$$
\item \href{https://artofproblemsolving.com/community/c6h1876757p12752810}{\textbf{2018 A1}} Let $\mathbb{Q}_{>0}$ denote the set of all positive rational numbers. Determine all functions $f:\mathbb{Q}_{>0}\to \mathbb{Q}_{>0}$ satisfying$$f(x^2f(y)^2)=f(x)^2f(y)$$for all $x,y\in\mathbb{Q}_{>0}$.

\item \href{https://artofproblemsolving.com/community/c6h1671272p10632289}{\textbf{2017 A1}} Let $a_1,a_2,\dots a_n,k$, and $M$ be positive integers such that
$$\frac{1}{a_1}+\frac{1}{a_2}+\dots+\frac{1}{a_n}=k\quad\text{and}\quad a_1a_2\dots a_n=M.$$If $M>1$, prove that the polynomial
$$P(x)=M(x+1)^k-(x+a_1)(x+a_2)\dots (x+a_n)$$has no positive roots.
\item \href{https://artofproblemsolving.com/community/c6h1480690p8639254}{\textbf{2016 A1}} 	Let $a$, $b$, $c$ be positive real numbers such that $\min(ab,bc,ca) \geqslant 1$. Prove that$$\sqrt[3]{(a^2+1)(b^2+1)(c^2+1)} \leqslant \left(\frac{a+b+c}{3}\right)^2 + 1.$$
\item \dbox{\href{https://artofproblemsolving.com/community/c6h1268809p6621766}{\textbf{2015 A1}}} Suppose that a sequence $a_1,a_2,\dots$ of positive real numbers satisfies\[a_{k+1}\geqslant  \frac{ka_k}{a_k^2+(k-1)}\]for every positive integer $k$. Prove that $a_1+a_2+\dots+a_n\geqslant  n$ for every $n\geqslant  2$.
\item \boxed{\href{https://artofproblemsolving.com/community/c6h596930p3542095}{{\textbf{2014 A1}}}} Let $a_0 < a_1 < a_2 \dots$ be an infinite sequence of positive integers. Prove that there exists a unique integer $n\geqslant  1$ such that \[a_n < \frac{a_0+a_1+a_2+\dots+a_n}{n} \leqslant a_{n+1}.\]
\item \href{https://artofproblemsolving.com/community/c6h597118p3543340}{\textbf{2013 A1}} Let $n$ be a positive integer and let $a_1, \dots, a_{n-1} $ be arbitrary real numbers. Define the sequences $u_0, \dots, u_n $ and $v_0, \dots, v_n $ inductively by $u_0 = u_1  = v_0 = v_1 = 1$, and $u_{k+1} = u_k + a_k u_{k-1}$, $v_{k+1} = v_k + a_{n-k} v_{k-1}$ for $k=1, \dots, n-1.$

Prove that $u_n = v_n.$
\item \boxed{\href{https://artofproblemsolving.com/community/c6h488500p2737336}{{\textbf{2012 A1}}}} Find all functions $f:\mathbb Z\rightarrow \mathbb Z$ such that, for all integers $a,b,c$ that satisfy $a+b+c=0$, the following equality holds:
\[f(a)^2+f(b)^2+f(c)^2=2f(a)f(b)+2f(b)f(c)+2f(c)f(a).\]
\item \boxed{\href{https://artofproblemsolving.com/community/c6h418793p2363530}{{\textbf{2011 A1}}}} 	Given any set $A = \{a_1, a_2, a_3, a_4\}$ of four distinct positive integers, we denote the sum $a_1 +a_2 +a_3 +a_4$ by $s_A$. Let $n_A$ denote the number of pairs $(i, j)$ with $1 \leqslant  i < j \leqslant 4$ for which $a_i +a_j$ divides $s_A$. Find all sets $A$ of four distinct positive integers which achieve the largest possible value of $n_A$.
\item \boxed{\href{https://artofproblemsolving.com/community/c6h356075p1935849}{{\textbf{2010 A1}}}} Find all function $f:\mathbb{R}\rightarrow\mathbb{R}$ such that for all $x,y\in\mathbb{R}$ the following equality holds\[
f(\left\lfloor x\right\rfloor y)=f(x)\left\lfloor f(y)\right\rfloor. \]

\item \href{https://artofproblemsolving.com/community/c6h355778p1932911}{\textbf{2009 A1}} Find the largest possible integer $k$, such that the following statement is true:
Let $2009$ arbitrary non-degenerated triangles be given. In every triangle the three sides are coloured, such that one is blue, one is red and one is white. Now, for every colour separately, let us sort the lengths of the sides. We obtain
\[ \left. \begin{array}{rcl}
 & b_1 \leqslant b_2\leqslant \dots\leqslant b_{2009} & \textrm{the lengths of the blue sides }\\
 & r_1 \leqslant r_2\leqslant \dots\leqslant r_{2009} & \textrm{the lengths of the red sides }\\
 \textrm{and } & w_1 \leqslant w_2\leqslant \dots\leqslant w_{2009} & \textrm{the lengths of the white sides }\\
 \end{array}\right.\]
Then there exist $k$ indices $j$ such that we can form a non-degenerated triangle with side lengths $b_j$, $r_j$, $w_j$.
\item \boxed{\href{https://artofproblemsolving.com/community/c6h215430p1191683}{\textbf{2008 A1}}} Find all functions $ f: (0, \infty) \mapsto (0, \infty)$ (so $f$ is a function from the positive real numbers) such that
\[ \frac {\left( f(w) \right)^2 + \left( f(x) \right)^2}{f(y^2) + f(z^2) } = \frac {w^2 + x^2}{y^2 + z^2}
\]
for all positive real numbers $ w,x,y,z,$ satisfying $ wx = yz.$
\item \boxed{\href{https://artofproblemsolving.com/community/c6h159835p893741}{\textbf{2007 A1}}} Real numbers $ a_{1}$, $ a_{2}$, $ \dots$, $ a_{n}$ are given. For each $ i$, $ (1 \leqslant i \leqslant n )$, define
\[ d_{i} = \max \{ a_{j}\mid 1 \leqslant j \leqslant i \} - \min \{ a_{j}\mid i \leqslant j \leqslant n \}
\]
and let $ d = \max \{d_{i}\mid 1 \leqslant i \leqslant n \}$.
\begin{enumerate}[(a)]
    \item Prove that, for any real numbers $ x_{1}\leqslant x_{2}\leqslant \dots \leqslant x_{n}$,
\[ \max \{ |x_{i} - a_{i}| \mid 1 \leqslant i \leqslant n \}\geqslant  \frac {d}{2}. \quad \quad (*)
\]
    \item Show that there are real numbers $ x_{1}\leqslant x_{2}\leqslant \dots \leqslant x_{n}$ such that the equality holds in (*).
\end{enumerate}


\item \href{https://artofproblemsolving.com/community/c6h130814p741371}{\textbf{2006 A1}} sequence of real numbers $ a_{0},\ a_{1},\ a_{2},\dots$ is defined by the formula
\[ a_{i + 1} = \left\lfloor a_{i}\right\rfloor\cdot \left\langle a_{i}\right\rangle\qquad\text{for}\quad i\geqslant  0;
\]here $a_0$ is an arbitrary real number, $\lfloor a_i\rfloor$ denotes the greatest integer not exceeding $a_i$, and $\left\langle a_i\right\rangle=a_i-\lfloor a_i\rfloor$. Prove that $a_i=a_{i+2}$ for $i$ sufficiently large.
\end{itemize}
\newpage
\section*{A2 Problems} \label{A2}
 \begin{itemize}[leftmargin=*]


\item \href{https://www.imo-official.org/assets/documents/problems/2026/IMO2026SL.pdf}{\textbf{2026 A2}} TO BE REVEALED!

\item \href{https://artofproblemsolving.com/community/c6h3866858p38676946}{\textbf{2025 A2}} The \textit{sunshine cost} of a sequence $a_1,a_2,\ldots,a_{100}$ of integers is the largest possible value of
\[|(a_1+a_2+\cdots+a_i)-a_j|\]as $i$ and $j$ vary over all integers $1,2,\ldots,100$.

Determine the smallest possible sunshine cost over all sequences $a_1, a_2, \ldots, a_{100}$ of pairwise distinct integers.
\item \href{https://artofproblemsolving.com/community/c6h3610446p35340919}{\textbf{2024 A2}} Let $n$ be a positive integer. Find the minimum possible value of
\[
S = 2^0 x_0^2 + 2^1 x_1^2 + \dots + 2^n x_n^2,
\]where $x_0, x_1, \dots, x_n$ are nonnegative integers such that $x_0 + x_1 + \dots + x_n = n$.
\item \href{https://artofproblemsolving.com/community/c6h3359723p31218373}{\textbf{2023 A2}} Let $\mathbb{R}$ be the set of real numbers. Let $f:\mathbb{R}\rightarrow\mathbb{R}$ be a function such that\[f(x+y)f(x-y)\geqslant  f(x)^2-f(y)^2\]for every $x,y\in\mathbb{R}$. Assume that the inequality is strict for some $x_0,y_0\in\mathbb{R}$.

Prove that either $f(x)\geqslant  0$ for every $x\in\mathbb{R}$ or $f(x)\leqslant 0$ for every $x\in\mathbb{R}$.
\item \href{https://artofproblemsolving.com/community/c6h3107342p28104303}{\textbf{2022 A2}} Let $k\geqslant 2$ be an integer. Find the smallest integer $n \geqslant k+1$ with the property that there exists a set of $n$ distinct real numbers such that each of its elements can be written as a sum of $k$ other distinct elements of the set. 
\item \href{https://artofproblemsolving.com/community/c6h2882567p25627655}{\textbf{2021 A2}} For every integer $n \geqslant 1$ consider the $n \times n$ table with entry $\left\lfloor \dfrac{ij}{n+1} \right\rfloor$ at the intersection of row $i$ and column $j$, for every $i=1,\dots,n$ and $j=1,\dots,n$. 

Determine all integers $n \geqslant  1$ for which the sum of the $n^2$ entries in the table is equal to $\frac{1}{4}n^2(n-1)$.
\item \href{https://artofproblemsolving.com/community/c6h2625905p22698387}{\textbf{2020 A2}} Let $\mathcal{A}$ denote the set of all polynomials in three variables $x, y, z$ with integer coefficients. Let $\mathcal{B}$ denote the subset of $\mathcal{A}$ formed by all polynomials which can be expressed as
\begin{align*}
(x + y + z)P(x, y, z) + (xy + yz + zx)Q(x, y, z) + xyzR(x, y, z)
\end{align*}with $P, Q, R \in \mathcal{A}$. Find the smallest non-negative integer $n$ such that $x^i y^j z^k \in \mathcal{B}$ for all non-negative integers $i, j, k$ satisfying $i + j + k \geqslant  n$.
\item \href{https://artofproblemsolving.com/community/c6h2279014p17828913}{\textbf{2019 A2}} Let $u_1, u_2, \dots, u_{2019}$ be real numbers satisfying\[u_{1}+u_{2}+\dots+u_{2019}=0 \quad \text { and } \quad u_{1}^{2}+u_{2}^{2}+\dots+u_{2019}^{2}=1.\]Let $a=\min \left(u_{1}, u_{2}, \dots, u_{2019}\right)$ and $b=\max \left(u_{1}, u_{2}, \dots, u_{2019}\right)$. Prove that
\[
a b \leqslant-\frac{1}{2019}.
\]
\item \boxed{\href{https://artofproblemsolving.com/community/c6h1670581p10626524}{\textbf{2018 A2}}} Find all integers $n \geqslant  3$ for which there exist real numbers $a_1, a_2, \dots a_{n + 2}$ satisfying $a_{n + 1} = a_1$, $a_{n + 2} = a_2$ and
$$a_ia_{i + 1} + 1 = a_{i + 2},$$for $i = 1, 2, \dots, n$.

\item \href{https://artofproblemsolving.com/community/c6h1671307p10632422}{\textbf{2017 A2} }	Let $q$ be a real number. Gugu has a napkin with ten distinct real numbers written on it, and he writes the following three lines of real numbers on the blackboard:
\begin{itemize}[leftmargin=*]
    \item[$\bullet$]  In the first line, Gugu writes down every number of the form $a-b$, where $a$ and $b$ are two (not necessarily distinct) numbers on his napkin.
\item[$\bullet$]In the second line, Gugu writes down every number of the form $qab$, where $a$ and $b$ are
two (not necessarily distinct) numbers from the first line.
\item[$\bullet$] In the third line, Gugu writes down every number of the form $a^2+b^2-c^2-d^2$, where $a, b, c, d$ are four (not necessarily distinct) numbers from the first line.
\end{itemize}
Determine all values of $q$ such that, regardless of the numbers on Gugu's napkin, every number in the second line is also a number in the third line.
\item \href{https://artofproblemsolving.com/community/c6h1480705p8639281}{\textbf{2016 A2}} Find the smallest constant $C > 0$ for which the following statement holds: among any five positive real numbers $a_1,a_2,a_3,a_4,a_5$ (not necessarily distinct), one can always choose distinct subscripts $i,j,k,l$ such that
\[ \left| \frac{a_i}{a_j} - \frac {a_k}{a_l} \right| \leqslant C. \]
\item \href{https://artofproblemsolving.com/community/c6h1268817p6621849}{\textbf{2015 A2}} Determine all functions $f:\mathbb{Z}\rightarrow\mathbb{Z}$ with the property that\[f(x-f(y))=f(f(x))-f(y)-1\]holds for all $x,y\in\mathbb{Z}$.
\item \href{https://artofproblemsolving.com/community/c6h1113206p5083585}{\textbf{2014 A2}} Define the function $f:(0,1)\to (0,1)$ by\[\displaystyle f(x) = \left\{ \begin{array}{lr} x+\frac 12 & \text{if}\ \  x < \frac 12\\ x^2 & \text{if}\ \  x \geqslant \frac 12 \end{array} \right.\]Let $a$ and $b$ be two real numbers such that $0 < a < b < 1$. We define the sequences $a_n$ and $b_n$ by $a_0 = a, b_0 = b$, and $a_n = f( a_{n -1})$, $b_n = f (b_{n -1} )$ for $n > 0$. Show that there exists a positive integer $n$ such that\[(a_n - a_{n-1})(b_n-b_{n-1})<0.\]
\item \href{https://artofproblemsolving.com/community/c6h597119p3543341}{\textbf{2013 A2}} Prove that in any set of $2000$ distinct real numbers there exist two pairs $a>b$ and $c>d$ with $a \neq c$ or $b \neq d $, such that\[ \left| \frac{a-b}{c-d} - 1 \right|< \frac{1}{100000}. \]
\newpage

\item \href{https://artofproblemsolving.com/community/c6h546162p3160551}{\textbf{2012 A2}} Let $\mathbb{Z}$ and $\mathbb{Q}$ be the sets of integers and rationals respectively.
\begin{enumerate}
\item[(a)] Does there exist a partition of $\mathbb{Z}$ into three non-empty subsets $A,B,C$ such that the sets 

$A+B, B+C, C+A$ are disjoint?
\item[(b)] Does there exist a partition of $\mathbb{Q}$ into three non-empty subsets $A,B,C$ such that the sets 

$A+B, B+C, C+A$ are disjoint?
\end{enumerate}
Here $X+Y$ denotes the set $\{ x+y : x \in X, y \in Y \}$, for $X,Y \subseteq \mathbb{Z}$ and for $X,Y \subseteq \mathbb{Q}$.
\item \href{https://artofproblemsolving.com/community/c6h488534p2737640}{\textbf{2011 A2}} Determine all sequences $(x_1,x_2,\dots,x_{2011})$ of positive integers, such that for every positive integer $n$ there exists an integer $a$ with\[\sum^{2011}_{j=1} j  x^n_j = a^{n+1} + 1.\]
\item \href{https://artofproblemsolving.com/community/c6h418679p2362276}{\textbf{2010 A2}} Let the real numbers $a,b,c,d$ satisfy the relations $a+b+c+d=6$ and $a^2+b^2+c^2+d^2=12.$ Prove that
\[36 \leqslant 4 \left(a^3+b^3+c^3+d^3\right) - \left(a^4+b^4+c^4+d^4 \right) \leqslant 48.\]
\item \href{https://artofproblemsolving.com/community/c6h355781p1932917}{\textbf{2009 A2}} Let $a$, $b$, $c$ be positive real numbers such that $\dfrac{1}{a} + \dfrac{1}{b} + \dfrac{1}{c} = a+b+c$. Prove that:
\[\frac{1}{(2a+b+c)^2}+\frac{1}{(a+2b+c)^2}+\frac{1}{(a+b+2c)^2}\leqslant \frac{3}{16}.\]
\item \boxed{\href{https://artofproblemsolving.com/community/c6h215222p1190551}{\textbf{2008 A2}}}
\begin{enumerate}[(a)]
    \item Prove that
\[\frac {x^{2}}{\left(x - 1\right)^{2}} + \frac {y^{2}}{\left(y - 1\right)^{2}} + \frac {z^{2}}{\left(z - 1\right)^{2}} \geqslant  1\]for all real numbers $x$, $y$, $z$, each different from $1$, and satisfying $xyz=1$.
    \item Prove that equality holds above for infinitely many triples of rational numbers $x$, $y$, $z$, each different from $1$, and satisfying $xyz=1$.
\end{enumerate}


 
\item \href{https://artofproblemsolving.com/community/c6h214702p1187155}{\textbf{2007 A2}} Consider those functions $ f: \mathbb{N} \mapsto \mathbb{N}$ which satisfy the condition
\[ f(m + n) \geqslant  f(m) + f(f(n)) - 1
\]
for all $ m,n \in \mathbb{N}.$ Find all possible values of $ f(2007).$
\item \href{https://artofproblemsolving.com/community/c6h145078p821388}{\textbf{2006 A2}} The sequence of real numbers $a_0,a_1,a_2,\dots$ is defined recursively by\[a_0=-1,\qquad\sum_{k=0}^n\dfrac{a_{n-k}}{k+1}=0\quad\text{for}\quad n\geqslant  1.\]Show that $ a_{n} > 0$ for all $ n\geqslant  1$.
\end{itemize}
\newpage
\section*{A3 Problems} \label{A3}  
 \begin{itemize}[leftmargin=*]

\item \href{https://www.imo-official.org/assets/documents/problems/2026/IMO2026SL.pdf}{\textbf{2026 A3}} TO BE REVEALED!

\item \boxed{\href{https://artofproblemsolving.com/community/c6h3610486p35341177}{\textbf{2025 A3}}} Alice and Bazza are playing the inekoalaty game, a two-player game whose rules depend on a positive real number $\lambda$ which is known to both players. On the $n^\text{th}$ turn of the game (starting with $n=1$) the following happens:
\begin{itemize}[leftmargin=*]
\item[$\bullet$] If $n$ is odd, Alice chooses a nonnegative real number $x_n$ such that
\[
    x_1 + x_2 + \dots + x_n \leqslant \lambda n.
  \]
\item[$\bullet$] If $n$ is even, Bazza chooses a nonnegative real number $x_n$ such that
\[
    x_1^2 + x_2^2 + \dots + x_n^2 \leqslant n.
  \]
\end{itemize}
If a player cannot choose a suitable $x_n$, the game ends and the other player wins. If the game goes on forever, neither player wins. All chosen numbers are known to both players.

Determine all values of $\lambda$ for which Alice has a winning strategy and all those for which Bazza has a winning strategy.
\item \href{https://artofproblemsolving.com/community/c6h3610463p35340954}{\textbf{2024 A3}} Decide whether for every sequence $(a_n)$ of positive real numbers,
$$\frac{3^{a_1}+3^{a_2}+\dots+3^{a_n}}{(2^{a_1}+2^{a_2}+\dots+2^{a_n})^2} < \frac{1}{2024}$$
is true for at least one positive integer $n$.
\item \boxed{\href{https://artofproblemsolving.com/community/c6h3107339p28104298}{\textbf{2023 A3}}} Let $x_1,x_2,\dots,x_{2023}$ be pairwise different positive real numbers such that
\[a_n=\sqrt{(x_1+x_2+\dots+x_n)\left(\frac{1}{x_1}+\frac{1}{x_2}+\dots+\frac{1}{x_n}\right)}\]is an integer for every $n=1,2,\dots,2023.$ Prove that $a_{2023} \geqslant  3034$.
\item \boxed{\href{https://artofproblemsolving.com/community/c6h2883212p25635138}{\textbf{2022 A3}}} Let $\mathbb{R}^+$ denote the set of positive real numbers. Find all functions $f: \mathbb{R}^+ \to \mathbb{R}^+$ such that for each $x \in \mathbb{R}^+$, there is exactly one $y \in \mathbb{R}^+$ satisfying$$xf(y)+yf(x) \leqslant 2.$$ 
\item \href{https://artofproblemsolving.com/community/c6h2882539p25627503}{\textbf{2021 A3}} For each integer $n\geqslant 1,$ compute the smallest possible value of  $\left\lfloor\dfrac{a_1}{1}\right\rfloor+\left\lfloor\dfrac{a_2}{2}\right\rfloor+\dots+\left\lfloor\dfrac{a_n}{n}\right\rfloor$ over all permutations $(a_1,\dots,a_n)$ of $(1,\dots,n)$.
\item \href{https://artofproblemsolving.com/community/c6h2625839p22697887}{\textbf{2020 A3}} Suppose that $a,b,c,d$ are positive real numbers satisfying $(a+c)(b+d)=ac+bd$. Find the smallest possible value of
$$\frac{a}{b}+\frac{b}{c}+\frac{c}{d}+\frac{d}{a}.$$
\newpage
\item \href{https://artofproblemsolving.com/community/c6h2279061p17829510}{\textbf{2019 A3}} Let $n \geqslant  3$ be a positive integer and let $\left(a_{1}, a_{2}, \dots, a_{n}\right)$ be a strictly increasing sequence of $n$ positive real numbers with sum equal to 2. Let $X$ be a subset of $\{1,2, \dots, n\}$ such that the value of
\[
\left|1-\sum_{i \in X} a_{i}\right|
\]is minimised. Prove that there exists a strictly increasing sequence of $n$ positive real numbers $\left(b_{1}, b_{2}, \dots, b_{n}\right)$ with sum equal to 2 such that
\[
\sum_{i \in X} b_{i}=1.
\]
\item \href{https://artofproblemsolving.com/community/c6h1876779p12752861}{\textbf{2018 A3}} 	Given any set $S$ of positive integers, show that at least one of the following two assertions holds:
\begin{enumerate}[(1)]
    \item There exist distinct finite subsets $F$ and $G$ of $S$ such that $\sum_{x\in F}1/x=\sum_{x\in G}1/x$;
    \item There exists a positive rational number $r<1$ such that $\sum_{x\in F}1/x\neq r$ for all finite subsets $F$ of $S$.
\end{enumerate}
\item \href{https://artofproblemsolving.com/community/c6h1671322p10632599}{\textbf{2017 A3}} Let $S$ be a finite set, and let $\mathcal{A}$ be the set of all functions from $S$ to $S$. Let $f$ be an element of $\mathcal{A}$, and let $T=f(S)$ be the image of $S$ under $f$. Suppose that $f\circ g\circ f\ne g\circ f\circ g$ for every $g$ in $\mathcal{A}$ with $g\ne f$. Show that $f(T)=T$.
\item \href{https://artofproblemsolving.com/community/c6h1480716p8639312}{\textbf{2016 A3}} Find all positive integers $n$ such that the following statement holds: Suppose real numbers $a_1$, $a_2$, $\dots$, $a_n$, $b_1$, $b_2$, $\dots$, $b_n$ satisfy $|a_k|+|b_k|=1$ for all $k=1,\dots,n$. Then there exists $\varepsilon_1$, $\varepsilon_2$, $\dots$, $\varepsilon_n$, each of which is either $-1$ or $1$, such that
\[ \left| \sum_{i=1}^n \varepsilon_i a_i \right| + \left| \sum_{i=1}^n \varepsilon_i b_i \right| \leqslant 1. \]
\item \href{https://artofproblemsolving.com/community/c6h1268841p6622146}{\textbf{2015 A3}} 	Let $n$ be a fixed positive integer. Find the maximum possible value of\[ \sum_{1 \leqslant r < s \leqslant 2n} (s-r-n)x_rx_s, \]where $-1 \leqslant x_i \leqslant 1$ for all $i = 1, \dots , 2n$.
\item \href{https://artofproblemsolving.com/community/c6h1113180p5083537}{\textbf{2014 A3}} For a sequence $x_1,x_2,\dots,x_n$ of real numbers, we define its $\textit{price}$ as\[\max_{1\leqslant i\leqslant n}|x_1+\dots +x_i|.\]Given $n$ real numbers, Dave and George want to arrange them into a sequence with a low price. Diligent Dave checks all possible ways and finds the minimum possible price $D$. Greedy George, on the other hand, chooses $x_1$ such that $|x_1 |$ is as small as possible; among the remaining numbers, he chooses $x_2$ such that $|x_1 + x_2 |$ is as small as possible, and so on. Thus, in the $i$-th step he chooses $x_i$ among the remaining numbers so as to minimise the value of $|x_1 + x_2 + \dots + x_i |$. In each step, if several numbers provide the same value, George chooses one at random. Finally he gets a sequence with price $G$. 

Find the least possible constant $c$ such that for every positive integer $n$, for every collection of $n$ real numbers, and for every possible sequence that George might obtain, the resulting values satisfy the inequality $G\leqslant cD$.
\newpage
\item \boxed{\href{https://artofproblemsolving.com/community/c6h1181547p5720286}{\textbf{2013 A3}}} Let $\mathbb Q_{>0}$ be the set of all positive rational numbers. Let $f:\mathbb Q_{>0}\to\mathbb R$ be a function satisfying the following three conditions: 
\begin{enumerate}[(i)]
    \item for all $x,y\in\mathbb Q_{>0}$, we have $f(x)f(y)\geqslant  f(xy)$; 
    \item for all $x,y\in\mathbb Q_{>0}$, we have $f(x+y)\geqslant  f(x)+f(y)$; 
    \item there exists a rational number $a>1$ such that $f(a)=a$.
Prove that $f(x)=x$ for all $x\in\mathbb Q_{>0}$.
\end{enumerate}
\item \boxed{\href{https://artofproblemsolving.com/community/c6h488342p2736375}{\textbf{2012 A3}}}  Let $n\geqslant 3$ be an integer, and let $a_2,a_3,\dots ,a_n$ be positive real numbers such that $a_{2}a_{3}\dots a_{n}=1$. Prove that
\[(1 + a_2)^2 (1 + a_3)^3 \dotsm (1 + a_n)^n > n^n.\]
\item \href{https://artofproblemsolving.com/community/c6h488535p2737643}{\textbf{2011 A3}} Determine all pairs $(f,g)$ of functions from the set of real numbers to itself that satisfy\[g(f(x+y)) = f(x) + (2x + y)g(y)\]for all real numbers $x$ and $y$.
\item \href{https://artofproblemsolving.com/community/c6h418680p2362280}{\textbf{2010 A3}} Let $x_1, \dots , x_{100}$ be nonnegative real numbers such that $x_i + x_{i+1} + x_{i+2} \leqslant 1$ for all $i = 1, \dots , 100$ (we put $x_{101 } = x_1, x_{102} = x_2).$ Find the maximal possible value of the sum $S = \displaystyle{\sum^{100}_{i=1} x_i x_{i+2}.}$
\item \boxed{\href{https://artofproblemsolving.com/community/c6h289055p1562848}{\textbf{2009 A3}} } Determine all functions $ f$ from the set of positive integers to the set of positive integers such that, for all positive integers $ a$ and $ b$, there exists a non-degenerate triangle with sides of lengths
\[ a, f(b) \text{ and } f(b + f(a) - 1).\]
(A triangle is non-degenerate if its vertices are not collinear.)
\item \href{https://artofproblemsolving.com/community/c6h287852p1555894}{\textbf{2008 A3}} 	Let $ S\subseteq\mathbb{R}$ be a set of real numbers. We say that a pair $ (f, g)$ of functions from $ S$ into $ S$ is a \textit{Spanish Couple} on $ S$, if they satisfy the following conditions: 
\begin{enumerate}[(i)]
    \item Both functions are strictly increasing, i.e. $ f(x) < f(y)$ and $ g(x) < g(y)$ for all $ x$, $ y\in S$ with $ x < y$;
    \item The inequality $ f\left(g\left(g\left(x\right)\right)\right) < g\left(f\left(x\right)\right)$ holds for all $ x\in S$.
\end{enumerate}



Decide whether there exists a Spanish Couple
on the set $ S = \mathbb{N}$ of positive integers;
on the set $ S = \{a - \frac {1}{b}: a, b\in\mathbb{N}\}$
\item \href{https://artofproblemsolving.com/community/c6h214703p1187159}{\textbf{2007 A3}} Let $ n$ be a positive integer, and let $ x$ and $ y$ be a positive real number such that $ x^n + y^n = 1.$ Prove that
\[ \left(\sum^n_{k = 1} \frac {1 + x^{2k}}{1 + x^{4k}} \right) \cdot \left( \sum^n_{k = 1} \frac {1 + y^{2k}}{1 + y^{4k}} \right) < \frac {1}{(1 - x) \cdot (1 - y)}.
\]
\item \href{https://artofproblemsolving.com/community/c6h155689p874971}{\textbf{2006 A3}} The sequence $c_{0}, c_{1},\dots , c_{n}, \dots $ is defined by $c_{0}= 1, c_{1}= 0$, and $c_{n+2}= c_{n+1}+c_{n}$ for $n \geqslant  0$. Consider the set $S$ of ordered pairs $(x, y)$ for which there is a finite set $J$ of positive integers such that $x=\textstyle\sum_{j \in J}{c_{j}}$, $y=\textstyle\sum_{j \in J}{c_{j-1}}$. Prove that there exist real numbers $\alpha$, $\beta$, and $M$ with the following property: An ordered pair of nonnegative integers $(x, y)$ satisfies the inequality\[m < \alpha x+\beta y < M\]if and only if $(x, y) \in S$.
\end{itemize}
\newpage
\section*{A4 Problems} \label{A4}
 \begin{itemize}[leftmargin=*]

\item \boxed{\href{https://artofproblemsolving.com/community/c6h3866898p38677166}{\textbf{2026 A4}}} Let $\mathbb{R}_{>0}$ be the set of positive real numbers. Determine all functions $f :\mathbb{R}_{>0}\to \mathbb{R}_{>0}$ such that
\[\sqrt{\frac{x^2 + f(y)^2}2}\geqslant \frac{f(x) + y}2 \geqslant \sqrt{xf(y)}\]
for every $x ,y\in\mathbb{R}_{>0}$ .

\item \href{https://artofproblemsolving.com/community/c6h3866883p38677019}{\textbf{2025 A4}} Let $\mathbb{Z}_{\geqslant 0}$ be the set of all nonnegative integers. Let $f:\mathbb{Z}_{\geqslant 0}\rightarrow\mathbb{Z}_{\geqslant 0}$ be an unbounded function such that, if $m$ and $n$ are nonnegative integers satisfying
\[ f(m+n)=\max\{f(0),f(1),\dots,f(m+n)\}, \]then
\[ f(m+n)=f(m)+f(n). \]Prove that there exist positive integers $A, B, C$ and $D$ such that for all nonnegative integers $n$,
\[ f(An+B)=Cn+D. \]
\textit{We say that $f$ is unbounded if for each nonnegative integer $N$, there exists some nonnegative integer $n$ such that $f(n)\geqslant N$.}
\item \href{https://artofproblemsolving.com/community/c6h3610435p35340902}{\textbf{2024 A4}} Let \(\mathbb{Z}_{>0}\) be the set of all positive integers. Determine all subsets \(\mathcal{S}\) of \(\{2^{0},2^{1},2^{2},\dots\}\) for which there exists a function \(f\colon\mathbb{Z}_{>0}\to\mathbb{Z}_{>0}\) such that
\[\mathcal{S}=\{f(a+b)-f(a)-f(b)\mid a,b\in\mathbb{Z}_{>0}\}.\]
\item \href{https://artofproblemsolving.com/community/c6h3359742p31218446}{\textbf{2023 A4}} Let $\mathbb R_{>0}$ be the set of positive real numbers. Determine all functions $f \colon \mathbb R_{>0} \to \mathbb R_{>0}$ such that\[x \big(f(x) + f(y)\big) \geqslant  \big(f(f(x)) + y\big) f(y)\]for every $x, y \in \mathbb R_{>0}$.
\item \href{https://artofproblemsolving.com/community/c6h3107365p28104481}{\textbf{2022 A4}} Let $n \geqslant  3$ be an integer, and let $x_1,x_2,\dots,x_n$ be real numbers in the interval $[0,1]$. Let $s=x_1+x_2+\dots+x_n$, and assume that $s \geqslant  3$. Prove that there exist integers $i$ and $j$ with $1 \leqslant i<j \leqslant n$ such that\[2^{j-i}x_ix_j>2^{s-3}.\]
\item \boxed{\href{https://artofproblemsolving.com/community/c6h2625850p22697952}{\textbf{2021 A4}}} 	Show that the inequality\[\sum_{i=1}^n \sum_{j=1}^n \sqrt{|x_i-x_j|}\leqslant \sum_{i=1}^n \sum_{j=1}^n \sqrt{|x_i+x_j|}\]holds for all real numbers $x_1,\dots x_n.$
\item \boxed{\href{https://artofproblemsolving.com/community/c6h2278647p17821569}{\textbf{2020 A4}}} The real numbers $a, b, c, d$ are such that $a\geqslant  b\geqslant  c\geqslant  d>0$ and $a+b+c+d=1$. Prove that
\[(a+2b+3c+4d)a^ab^bc^cd^d<1.\]
\item \href{https://artofproblemsolving.com/community/c6h2279021p17829013}{\textbf{2019 A4}} Let $n\geqslant  2$ be a positive integer and $a_1,a_2, \dots ,a_n$ be real numbers such that\[a_1+a_2+\dots+a_n=0.\]Define the set $A$ by
\[A=\left\{(i, j)\,|\,1 \leqslant i<j \leqslant n,\left|a_{i}-a_{j}\right| \geqslant  1\right\}\]Prove that, if $A$ is not empty, then
\[\sum_{(i, j) \in A} a_{i} a_{j}<0.\]
\item \href{https://artofproblemsolving.com/community/c6h1876762p12752819}{\textbf{2018 A4}} Let $a_0,a_1,a_2,\dots $ be a sequence of real numbers such that $a_0=0, a_1=1,$ and for every $n\geqslant  2$ there exists $1 \leqslant k \leqslant n$ satisfying\[ a_n=\frac{a_{n-1}+\dots + a_{n-k}}{k}. \]Find the maximum possible value of $a_{2018}-a_{2017}$.
\item \href{https://artofproblemsolving.com/community/c6h1671279p10632313}{\textbf{2017 A4}} A sequence of real numbers $a_1,a_2,\dots$ satisfies the relation
$$a_n=-\max_{i+j=n}(a_i+a_j)\qquad\text{for all}\quad n>2017.$$Prove that the sequence is bounded, i.e., there is a constant $M$ such that $|a_n|\leqslant M$ for all positive integers $n$.
\item \href{https://artofproblemsolving.com/community/c6h1480721p8639325}{\textbf{2016 A4}} Find all functions $f:(0,\infty)\rightarrow (0,\infty)$ such that for any $x,y\in (0,\infty)$,$$xf(x^2)f(f(y)) + f(yf(x)) = f(xy) \left(f(f(x^2)) + f(f(y^2))\right).$$
\item \boxed{\href{https://artofproblemsolving.com/community/c6h1113162p5083463}{\textbf{2015 A4}}} Let $\mathbb R$ be the set of real numbers. Determine all functions $f:\mathbb R\to\mathbb R$ that satisfy the equation\[f(x+f(x+y))+f(xy)=x+f(x+y)+yf(x)\]for all real numbers $x$ and $y$.
\item \href{https://artofproblemsolving.com/community/c6h1113177p5083531}{\textbf{2014 A4}} Determine all functions $f: \mathbb{Z}\to\mathbb{Z}$ satisfying\[f\big(f(m)+n\big)+f(m)=f(n)+f(3m)+2014\]for all integers $m$ and $n$.
\item \href{https://artofproblemsolving.com/community/c6h597122p3543362}{\textbf{2013 A4}} 	Let $n$ be a positive integer, and consider a sequence $a_1 , a_2 , \dots , a_n $ of positive integers. Extend it periodically to an infinite sequence $a_1 , a_2 , \dots $ by defining $a_{n+i} = a_i $ for all $i \geqslant 1$. If\[a_1 \leqslant a_2 \leqslant \dots \leqslant a_n \leqslant a_1 +n  \]and\[a_{a_i } \leqslant n+i-1 \quad\text{for}\quad i=1,2,\dots, n, \]prove that\[a_1 + \dots +a_n \leqslant n^2. \]
\item \href{https://artofproblemsolving.com/community/c6h546164p3160553}{\textbf{2012 A4}} 	Let $f$ and $g$ be two nonzero polynomials with integer coefficients and $\deg f>\deg g$. Suppose that for infinitely many primes $p$ the polynomial $pf+g$ has a rational root. Prove that $f$ has a rational root.
\item \href{https://artofproblemsolving.com/community/c6h488536p2737644}{\textbf{2011 A4}} Determine all pairs $(f,g)$ of functions from the set of positive integers to itself that satisfy\[f^{g(n)+1}(n) + g^{f(n)}(n) = f(n+1) - g(n+1) + 1\]for every positive integer $n$. Here, $f^k(n)$ means $\underbrace{f(f(\dots f)}_{k}(n) \dots ))$.
\item \href{https://artofproblemsolving.com/community/c6h418681p2362283}{\textbf{2010 A4}} A sequence $x_1, x_2, \dots$ is defined by $x_1 = 1$ and $x_{2k}=-x_k, x_{2k-1} = (-1)^{k+1}x_k$ for all $k \geqslant  1.$ Prove that $x_1 + x_2 + \dots + x_n \geqslant  0$ for all $n\geqslant 1$.
\item \href{https://artofproblemsolving.com/community/c6h355782p1932920}{\textbf{2009 A4}} 	Let $a$, $b$, $c$ be positive real numbers such that $ab+bc+ca\leqslant 3abc$. Prove that
\[\sqrt{\frac{a^2+b^2}{a+b}}+\sqrt{\frac{b^2+c^2}{b+c}}+\sqrt{\frac{c^2+a^2}{c+a}}+3\leqslant \sqrt{2}\left(\sqrt{a+b}+\sqrt{b+c}+\sqrt{c+a}\right)\] 
\item \href{https://artofproblemsolving.com/community/c6h287853p1555896}{\textbf{2008 A4}} 	For an integer $ m$, denote by $ t(m)$ the unique number in $ \{1, 2, 3\}$ such that $ m + t(m)$ is a multiple of $ 3$. A function $ f: \mathbb{Z}\to\mathbb{Z}$ satisfies $ f( - 1) = 0$, $ f(0) = 1$, $ f(1) = - 1$ and $ f\left(2^{n} + m\right) = f\left(2^n - t(m)\right) - f(m)$ for all integers $ m$, $ n\geqslant 0$ with $ 2^n > m$. Prove that $ f(3p)\geqslant 0$ holds for all integers $ p\geqslant 0$.
\item \href{https://artofproblemsolving.com/community/c6h211465p1165901}{\textbf{2007 A4}} Find all functions $ f: \mathbb{R}^{ + }\to\mathbb{R}^{ + }$ satisfying $ f\left(x + f\left(y\right)\right) = f\left(x + y\right) + f\left(y\right)$ for all pairs of positive reals $ x$ and $ y$. Here, $ \mathbb{R}^{ + }$ denotes the set of all positive reals. 
\item \href{https://artofproblemsolving.com/community/c6h155691p874975}{\textbf{2006 A4}} Prove the inequality:
\[\sum_{i < j}{\frac {a_{i}a_{j}}{a_{i} + a_{j}}}\leqslant \frac {n}{2(a_{1} + a_{2} +\dots + a_{n})}\cdot \sum_{i < j}{a_{i}a_{j}}\]
for positive reals $ a_{1},a_{2},\dots,a_{n}$.
\end{itemize}


 
\newpage
\section*{A5 Problems}  \label{A5}
 \begin{itemize}[leftmargin=*]

 \item \href{https://www.imo-official.org/assets/documents/problems/2026/IMO2026SL.pdf}{\textbf{2026 A5}} TO BE REVEALED!
 
 \item \href{https://artofproblemsolving.com/community/c6h3866881p38677013}{\textbf{2025 A5}} Determine all integers $n \geqslant 2$ such that the polynomial
 $$P_n(x,y,z)=x^n+y^n+z^n-xyz^{n-2}-xy^{n-2}z-x^{n-2}yz$$
 can be written as $P_n(x,y,z)=Q_n(x,y,z)\cdot R_n(x,y,z)$, where $Q_n(x,y,z)$ and $R_n(x,y,z)$ are non-constant polynomials with integer coefficients.
\item \href{https://artofproblemsolving.com/community/c6h3610458p35340939}{\textbf{2024 A5}} Find all periodic sequence $a_1,a_2,\dots$ of real numbers such that the following conditions hold for all $n\geqslant  1$:$$a_{n+2}+a_{n}^2=a_n+a_{n+1}^2\quad\text{and}\quad |a_{n+1}-a_n|\leqslant 1.$$
\item \href{https://artofproblemsolving.com/community/c6h3359722p31218372}{\textbf{2023 A5}} Let $a_1,a_2,\dots,a_{2023}$ be positive integers such that
\begin{itemize}[leftmargin=*]
    \item[$\bullet$]
$a_1,a_2,\dots,a_{2023}$ is a permutation of $1,2,\dots,2023$, and
\item[$\bullet$] $|a_1-a_2|,|a_2-a_3|,\dots,|a_{2022}-a_{2023}|$ is a permutation of $1,2,\dots,2022$.
\end{itemize}
Prove that $\max(a_1,a_{2023})\geqslant 507$.
    \item \href{https://artofproblemsolving.com/community/c6h3107362p28104461}{\textbf{2022 A5}} Find all positive integers $n \geqslant  2$ for which there exist $n$ real numbers $a_1<\dots<a_n$ and a real number $r>0$ such that the $\tfrac{1}{2}n(n-1)$ differences $a_j-a_i$ for $1 \leqslant i<j \leqslant n$ are equal, in some order, to the numbers $r^1,r^2,\dots,r^{\frac{1}{2}n(n-1)}$.
\item \href{https://artofproblemsolving.com/community/c6h2882545p25627517}{\textbf{2021 A5}} Let $n\geqslant  2$ be an integer and let $a_1, a_2, \dots, a_n$ be positive real numbers with sum $1$. Prove that$$\sum_{k=1}^n \frac{a_k}{1-a_k}(a_1+a_2+\dots+a_{k-1})^2 < \frac{1}{3}.$$
\item \href{https://artofproblemsolving.com/community/c6h2625872p22698165}{\textbf{2020 A5}} 	A magician intends to perform the following trick. She announces a positive integer $n$, along with $2n$ real numbers $x_1 < \dots < x_{2n}$, to the audience. A member of the audience then secretly chooses a polynomial $P(x)$ of degree $n$ with real coefficients, computes the $2n$ values $P(x_1), \dots , P(x_{2n})$, and writes down these $2n$ values on the blackboard in non-decreasing order. After that the magician announces the secret polynomial to the audience. Can the magician find a strategy to perform such a trick?
\item \href{https://artofproblemsolving.com/community/c6h2279006p17828803}{\textbf{2019 A5}} Let $x_1, x_2, \dots, x_n$ be different real numbers. Prove that
\[\sum_{1 \leqslant i \leqslant n} \prod_{j \neq i} \frac{1-x_{i} x_{j}}{x_{i}-x_{j}}=\left\{\begin{array}{ll}
0, & \text { if } n \text { is even; } \\
1, & \text { if } n \text { is odd. }
\end{array}\right.\]
\item \href{https://artofproblemsolving.com/community/c6h1876752p12752798}{\textbf{2018 A5}} Determine all functions $f:(0,\infty)\to\mathbb{R}$ satisfying$$\left(x+\frac{1}{x}\right)f(y)=f(xy)+f\left(\frac{y}{x}\right)$$for all $x,y>0$.
\newpage
\item \href{https://artofproblemsolving.com/community/c6h1671280p10632315}{\textbf{2017 A5}} An integer $n \geqslant  3$ is given. We call an $n$-tuple of real numbers $(x_1, x_2, \dots, x_n)$ Shiny if for each permutation $y_1, y_2, \dots, y_n$ of these numbers, we have
$$\sum \limits_{i=1}^{n-1} y_i y_{i+1} = y_1y_2 + y_2y_3 + y_3y_4 + \dots + y_{n-1}y_n \geqslant  -1.$$Find the largest constant $K = K(n)$ such that
$$\sum \limits_{1 \leqslant i < j \leqslant n} x_i x_j \geqslant  K$$holds for every Shiny $n$-tuple $(x_1, x_2, \dots, x_n)$.
\item \href{https://artofproblemsolving.com/community/c6h1480728p8639342}{\textbf{2016 A5}} 	Consider fractions $\frac{a}{b}$ where $a$ and $b$ are positive integers. 
\begin{enumerate}[(a)]
    \item Prove that for every positive integer $n$, there exists such a fraction $\frac{a}{b}$ such that $\sqrt{n} \leqslant \frac{a}{b} \leqslant \sqrt{n+1}$ and $b \leqslant \sqrt{n}+1$. 
    \item Show that there are infinitely many positive integers $n$ such that no such fraction $\frac{a}{b}$ satisfies $\sqrt{n} \leqslant \frac{a}{b} \leqslant \sqrt{n+1}$ and $b \leqslant \sqrt{n}$.
\end{enumerate}
\item \href{https://artofproblemsolving.com/community/c6h1268849p6622192}{\textbf{2015 A5}} Let $2\mathbb{Z} + 1$ denote the set of odd integers. Find all functions $f:\mathbb{Z} \mapsto 2\mathbb{Z} + 1$ satisfying\[ f(x + f(x) + y) + f(x - f(x) - y) = f(x+y) + f(x-y) \]for every $x, y \in \mathbb{Z}$. 
\item \href{https://artofproblemsolving.com/community/c6h1113181p5083541}{\textbf{2014 A5}} 	Consider all polynomials $P(x)$ with real coefficients that have the following property: for any two real numbers $x$ and $y$ one has\[|y^2-P(x)|\leqslant 2|x|\quad\text{if and only if}\quad |x^2-P(y)|\leqslant 2|y|.\]Determine all possible values of $P(0)$. 
\item \href{https://artofproblemsolving.com/community/c6h597123p3543371}{\textbf{2013 A5}} Let $\mathbb{Z}_{\geqslant 0}$ be the set of all nonnegative integers. Find all the functions $f: \mathbb{Z}_{\geqslant 0} \rightarrow \mathbb{Z}_{\geqslant 0} $ satisfying the relation
\[ f(f(f(n))) = f(n+1 ) +1 \]
for all $ n\in \mathbb{Z}_{\geqslant 0}$.
\item \href{https://artofproblemsolving.com/community/c6h546165p3160554}{\textbf{2012 A5}} Find all functions $f:\mathbb{R} \rightarrow \mathbb{R}$ that satisfy the conditions
\[f(1+xy)-f(x+y)=f(x)f(y) \quad \text{for all } x,y \in \mathbb{R},\]
and $f(-1) \neq 0$.
\newpage
\item \href{https://artofproblemsolving.com/community/c6h488537p2737645}{\textbf{2011 A5}} Prove that for every positive integer $n,$ the set $\{2,3,4,\dots,3n+1\}$ can be partitioned into $n$ triples in such a way that the numbers from each triple are the lengths of the sides of some obtuse triangle. 
\item \href{https://artofproblemsolving.com/community/c6h418682p2362286}{\textbf{2010 A5}} Denote by $\mathbb{Q}^+$ the set of all positive rational numbers. Determine all functions $f : \mathbb{Q}^+ \mapsto \mathbb{Q}^+$ which satisfy the following equation for all $x, y \in \mathbb{Q}^+:$\[f\left( f(x)^2y \right) = x^3 f(xy).\] 
\item \href{https://artofproblemsolving.com/community/c6h355779p1932913}{\textbf{2009 A5}} Let $f$ be any function that maps the set of real numbers into the set of real numbers. Prove that there exist real numbers $x$ and $y$ such that\[f\left(x-f(y)\right)>yf(x)+x\] 	
\item \href{https://artofproblemsolving.com/community/c6h287855p1555899}{\textbf{2008 A5}} Let $ a$, $ b$, $ c$, $ d$ be positive real numbers such that $ abcd = 1$ and $ a + b + c + d > \dfrac{a}{b} + \dfrac{b}{c} + \dfrac{c}{d} + \dfrac{d}{a}$. Prove that
\[ a + b + c + d < \dfrac{b}{a} + \dfrac{c}{b} + \dfrac{d}{c} + \dfrac{a}{d}\] 
\item \href{https://artofproblemsolving.com/community/c6h214704p1187165}{\textbf{2007 A5}} Let $ c > 2,$ and let $ a(1), a(2), \dots$ be a sequence of nonnegative real numbers such that
\[ a(m + n) \leqslant 2 \cdot a(m) + 2 \cdot a(n) \text{ for all } m,n \geqslant  1,
\]
and $ a\left(2^k \right) \leqslant \dfrac {1}{(k + 1)^c} \text{ for all } k \geqslant  0.$ Prove that the sequence $ a(n)$ is bounded. 
\item \href{https://artofproblemsolving.com/community/c6h130812p741368}{\textbf{2006 A5}} 	If $a,b,c$ are the sides of a triangle, prove that
\[\frac{\sqrt{b+c-a}}{\sqrt{b}+\sqrt{c}-\sqrt{a}}+\frac{\sqrt{c+a-b}}{\sqrt{c}+\sqrt{a}-\sqrt{b}}+\frac{\sqrt{a+b-c}}{\sqrt{a}+\sqrt{b}-\sqrt{c}}\leqslant 3 \]
\end{itemize}


 
\newpage
\section*{A6 Problems}  \label{A6}
 \begin{itemize}[leftmargin=*]

\item \href{https://www.imo-official.org/assets/documents/problems/2026/IMO2026SL.pdf}{\textbf{2026 A6}} TO BE REVEALED!

 \item \href{https://artofproblemsolving.com/community/c6h3866851p38676935}{\textbf{2025 A6}}  Let $S$ be a set of positive integers, possibly infinite, such that no positive integer greater than $1$ divides all elements of $S$. Determine all non-periodic infinite sequences $a_1,a_2,a_3,\dots$ of positive integers such that, for all positive integers $n$,
\begin{itemize}[leftmargin=*] 
\item[$\bullet$] $a_n\leqslant |a_{n+\ell}-\ell|$ for all $\ell$ in $S$, and
\item[$\bullet$] $a_n=|a_{n+\ell}-\ell|$ for at least one $\ell$ in $S$.
\end{itemize}
\textit{We say that an infinite sequence $a_1,a_2,a_3,\dots$ is periodic if there exists a positive integer $t$ such that $a_n=a_{n+t}$ for all positive integer $n$.}
\item \href{https://artofproblemsolving.com/community/c6h3610454p35340929}{\textbf{2024 A6}} Let $a_0, a_1, a_2, \dots$ be an infinite strictly increasing sequence of positive integers such that for each $n \geqslant  1$ we have

$$
a_n \in \left\{ \frac{a_{n-1} + a_{n+1}}{2},\ \sqrt{a_{n-1} \cdot a_{n+1}} \right\}.
$$Let $b_1, b_2, \dots$ be an infinite sequence of letters defined as

$$
b_n = 
\begin{cases}
A & \text{if } a_n = \frac{1}{2}(a_{n-1} + a_{n+1}) \\
G & \text{otherwise}.
\end{cases}
$$Prove that there exist positive integers $n_0$ and $d$ such that for all $n \geqslant  n_0$ we have $b_{n+d} = b_n$.
\item \boxed{\href{https://artofproblemsolving.com/community/c6h3106754p28097579}{\textbf{2023 A6}}} For each integer $k\geqslant  2$, determine all infinite sequences of positive integers $a_1$, $a_2$, $\dots$ for which there exists a polynomial $P$ of the form\[ P(x)=x^k+c_{k-1}x^{k-1}+\dots + c_1 x+c_0, \]where $c_0$, $c_1$, $\dots$ , $c_{k-1}$ are non-negative integers, such that\[ P(a_n)=a_{n+1}a_{n+2}\dots a_{n+k} \]for every integer $n\geqslant  1$.
    \item  \href{https://artofproblemsolving.com/community/c6h3107327p28104254}{\textbf{2022 A6}} Let $\mathbb R$ be the set of real numbers. We denote by $\mathcal F$ the set of all functions $f\colon\mathbb R\to\mathbb R$ such that
$$f(x + f(y)) = f(x) + f(y)$$for every $x,y\in\mathbb R$ Find all rational numbers $q$ such that for every function $f\in\mathcal F$, there exists some $z\in\mathbb R$ satisfying $f(z)=qz$.
\item \boxed{\href{https://artofproblemsolving.com/community/c6h2625864p22698082}{\textbf{2021 A6}}} Let $m\geqslant 2$ be an integer, $A$ a finite set of integers (not necessarily positive) and $B_1,B_2,\dots,B_m$ subsets of $A$. Suppose that, for every $k=1,2,\dots,m$, the sum of the elements of $B_k$ is $m^k$. 

Prove that $A$ contains at least $\dfrac{m}{2}$ elements.
\item \href{https://artofproblemsolving.com/community/c6h2625860p22698062}{\textbf{2020 A6}} Find all functions $f : \mathbb{Z}\rightarrow \mathbb{Z}$ satisfying
\[f^{a^{2} + b^{2}}(a+b) = af(a) +bf(b)\]for all integers $a$ and $b$.

\item \href{https://artofproblemsolving.com/community/c6h2278982p17828583}{\textbf{2019 A6}} 	A polynomial $P(x, y, z)$ in three variables with real coefficients satisfies the identities
$$P(x, y, z)=P(x, y, xy-z)=P(x, zx-y, z)=P(yz-x, y, z).$$
Prove that there exists a polynomial $F(t)$ in one variable such that
$$P(x,y,z)=F(x^2+y^2+z^2-xyz).$$
\item \href{https://artofproblemsolving.com/community/c6h1876783p12752880}{\textbf{2018 A6}} 	Let $m,n\geqslant  2$ be integers. Let $f(x_1,\dots, x_n)$ be a polynomial with real coefficients such that$$f(x_1,\dots, x_n)=\left\lfloor \frac{x_1+\dots + x_n}{m} \right\rfloor\text{ for every } x_1,\dots, x_n\in \{0,1,\dots, m-1\}.$$Prove that the total degree of $f$ is at least $n$.
\item \boxed{\href{https://artofproblemsolving.com/community/c6h1480146p8633190}{\textbf{2017 A6}}} Let $\mathbb{R}$ be the set of real numbers. Determine all functions $f: \mathbb{R} \rightarrow \mathbb{R}$ such that, for any real numbers $x$ and $y$,\[ f(f(x)f(y)) + f(x+y) = f(xy). \]
\item \boxed{\href{https://artofproblemsolving.com/community/c6h1270996p6642565}{\textbf{2016 A6}}} The equation
$$(x-1)(x-2)\dots(x-2016)=(x-1)(x-2)\dots (x-2016)$$is written on the board, with $2016$ linear factors on each side. What is the least possible value of $k$ for which it is possible to erase exactly $k$ of these $4032$ linear factors so that at least one factor remains on each side and the resulting equation has no real solutions?
\item \href{https://artofproblemsolving.com/community/c6h1268901p6622740}{\textbf{2015 A6}} Let $n$ be a fixed integer with $n \geqslant 2$. We say that two polynomials $P$ and $Q$ with real coefficients are block-similar if for each $i \in \{1, 2, \dots, n\}$ the sequences
\begin{eqnarray*}
P(2015i), P(2015i - 1), \dots, P(2015i - 2014) & \text{and}\\
Q(2015i), Q(2015i - 1), \dots, Q(2015i - 2014)
\end{eqnarray*}
are permutations of each other. 
\begin{enumerate}[(a)]
    \item Prove that there exist distinct block-similar polynomials of degree $n + 1$. 
    \item Prove that there do not exist distinct block-similar polynomials of degree $n$. 
\end{enumerate}
\item \href{https://artofproblemsolving.com/community/c6h1113182p5083542}{\textbf{2014 A6}} Find all functions $f : \mathbb{Z} \to\mathbb{ Z}$ such that
\[ n^2+4f(n)=f(f(n))^2 \]
for all $n\in \mathbb{Z}$. 
\item \href{https://artofproblemsolving.com/community/c6h597124p3543375}{\textbf{2013 A6}} 	Let $m \neq 0 $ be an integer. Find all polynomials $P(x) $ with real coefficients such that
\[ (x^3 - mx^2 +1 ) P(x+1)  + (x^3+mx^2+1) P(x-1) =2(x^3 - mx +1 ) P(x) \]
for all real number $x$.
\item \href{https://artofproblemsolving.com/community/c6h546166p3160555}{\textbf{2012 A6}} Let $f: \mathbb{N} \rightarrow \mathbb{N}$ be a function, and let $f^m$ be $f$ applied $m$ times. Suppose that for every $n \in \mathbb{N}$ there exists a $k \in \mathbb{N}$ such that $f^{2k}(n)=n+k$, and let $k_n$ be the smallest such $k$. 

Prove that the sequence $k_1,k_2,\dots $ is unbounded.
\item \boxed{\href{https://artofproblemsolving.com/community/c6h418798p2363539}{\textbf{2011 A6}}}	Let $f : \mathbb R \to \mathbb R$ be a real-valued function defined on the set of real numbers that satisfies
\[f(x + y) \leqslant yf(x) + f(f(x))\]
for all real numbers $x$ and $y$. Prove that $f(x) = 0$ for all $x \leqslant 0$. 
\item \href{https://artofproblemsolving.com/community/c6h418683p2362289}{\textbf{2010 A6}} Suppose that $f$ and $g$ are two functions defined on the set of positive integers and taking positive integer values. Suppose also that the equations $f(g(n)) = f(n) + 1$ and $g(f(n)) = g(n) + 1$ hold for all positive integers. Prove that $f(n) = g(n)$ for all positive integer $n.$ 
\item \boxed{\href{https://artofproblemsolving.com/community/c6h288840p1561573}{\textbf{2009 A6}}} Suppose that $ s_1,s_2,s_3, \dots$ is a strictly increasing sequence of positive integers such that the sub-sequences\[s_{s_1},\, s_{s_2},\, s_{s_3},\, \dots\qquad\text{and}\qquad s_{s_1+1},\, s_{s_2+1},\, s_{s_3+1},\, \dots\]are both arithmetic progressions. Prove that the sequence $ s_1, s_2, s_3, \dots$ is itself an arithmetic progression. 	
\item \href{https://artofproblemsolving.com/community/c6h287857p1555902}{\textbf{2008 A6}} 	Let $ f: \mathbb{R}\to\mathbb{N}$ be a function which satisfies $ f\left(x + \dfrac{1}{f(y)}\right) = f\left(y + \dfrac{1}{f(x)}\right)$ for all $ x$, $ y\in\mathbb{R}$. Prove that there is a positive integer which is not a value of $ f$. 
\item \href{https://artofproblemsolving.com/community/c6h214705p1187171}{\textbf{2007 A6}} Let $ a_1, a_2, \dots, a_{100}$ be nonnegative real numbers such that $ a^2_1 + a^2_2 + \dots + a^2_{100} = 1.$ Prove that
\[ a^2_1 \cdot a_2 + a^2_2 \cdot a_3 + \dots + a^2_{100} \cdot a_1 < \frac {12}{25}.
\] 
\item \boxed{\href{https://artofproblemsolving.com/community/c6h101294p571945}{\textbf{2006 A6}}} Determine the least real number $M$ such that the inequality\[|ab(a^{2}-b^{2})+bc(b^{2}-c^{2})+ca(c^{2}-a^{2})| \leqslant M(a^{2}+b^{2}+c^{2})^{2}\]holds for all real numbers $a$, $b$ and $c$.
\end{itemize}


 
\newpage
\section*{A7 Problems}  \label{A7}
 \begin{itemize}[leftmargin=*]

 \item \href{https://www.imo-official.org/assets/documents/problems/2026/IMO2026SL.pdf}{\textbf{2026 A7}} TO BE REVEALED!
 
 \item \href{https://artofproblemsolving.com/community/c6h3866874p38676992}{\textbf{2025 A7}}  For each integer $k\geqslant 3$, prove that there exists a unique $k$-tuplet$(x_1,x_2,\dots,x_k)$ of positive real numbers such that $x_1\geqslant x_2\geqslant\cdots \geqslant x_k$, and$$\left\lfloor nx_1+\frac{1}{2}\right\rfloor+\left\lfloor nx_2+\frac{1}{2}\right\rfloor+\cdots+\left\lfloor nx_k+\frac{1}{2}\right\rfloor=n$$for every integer $n$.

\item \boxed{\href{https://artofproblemsolving.com/community/c6h3359771p31218720}{\textbf{2024 A7}}} Let $\mathbb{Q}$ be the set of rational numbers. A function $f: \mathbb{Q} \to \mathbb{Q}$ is called \textit{aquaesulian} if the following property holds: for every $x,y \in \mathbb{Q}$,
\[ f(x+f(y)) = f(x) + y \quad \text{or} \quad f(f(x)+y) = x + f(y). \]Show that there exists an integer $c$ such that for any aquaesulian function $f$ there are at most $c$ different rational numbers of the form $f(r) + f(-r)$ for some rational number $r$, and find the smallest possible value of $c$.
\item \href{https://artofproblemsolving.com/community/c6h3359726p31218377}{\textbf{2023 A7}} Let $N$ be a positive integer. Prove that there exist three permutations $a_1,\dots,a_N$, $b_1,\dots,b_N$, and $c_1,\dots,c_N$ of $1,\dots,N$ such that\[\left|\sqrt{a_k}+\sqrt{b_k}+\sqrt{c_k}-2\sqrt{N}\right|<2023\]for every $k=1,2,\dots,N$.

    \item \href{https://artofproblemsolving.com/community/c6h3107335p28104281}{\textbf{2022 A7}} 	For a positive integer $n$ we denote by $s(n)$ the sum of the digits of $n$. Let $P(x)=x^n+a_{n-1}x^{n-1}+\dots+a_1x+a_0$ be a polynomial, where $n \geqslant  2$ and $a_i$ is a positive integer for all $0 \leqslant i \leqslant n-1$. 
    
    Could it be the case that, for all positive integers $k$, $s(k)$ and $s(P(k))$ have the same parity?
    
\item \href{https://artofproblemsolving.com/community/c6h2882578p25627708}{\textbf{2021 A7}} Let $n\geqslant  1$ be an integer, and let $x_0,x_1,\dots,x_{n+1}$ be $n+2$ non-negative real numbers that satisfy $x_ix_{i+1}-x_{i-1}^2\geqslant  1$ for all $i=1,2,\dots,n.$ Show that\[x_0+x_1+\dots+x_n+x_{n+1}>\bigg(\frac{2n}{3}\bigg)^{3/2}.\]
\item \href{https://artofproblemsolving.com/community/c6h2625876p22698187}{\textbf{2020 A7}} Let $n$ and $k$ be positive integers. Prove that for $a_1, \dots, a_n \in [1,2^k]$ one has
\[ \sum_{i = 1}^n \frac{a_i}{\sqrt{a_1^2 + \dots + a_i^2}} \leqslant 4 \sqrt{kn}. \]
\newpage

\item \href{https://artofproblemsolving.com/community/c6h2278978p17828524}{\textbf{2019 A7}} Let $\mathbb Z$ be the set of integers. We consider functions $f :\mathbb Z\to\mathbb Z$ satisfying
\[f\left(f(x+y)+y\right)=f\left(f(x)+y\right)\]for all integers $x$ and $y$. For such a function, we say that an integer $v$ is f-rare if the set
\[X_v=\{x\in\mathbb Z:f(x)=v\}\]is finite and nonempty. 
\begin{enumerate}[(a)]
    \item Prove that there exists such a function $f$ for which there is an $f$-rare integer. 
    \item Prove that no such function $f$ can have more than one $f$-rare integer.
\end{enumerate}
\item \href{https://artofproblemsolving.com/community/c6h1876747p12752777}{\textbf{2018 A7}} Find the maximal value of
\[S = \sqrt[3]{\frac{a}{b+7}} + \sqrt[3]{\frac{b}{c+7}} + \sqrt[3]{\frac{c}{d+7}} + \sqrt[3]{\frac{d}{a+7}},\]where $a$, $b$, $c$, $d$ are nonnegative real numbers which satisfy $a+b+c+d = 100$.
\item \href{https://artofproblemsolving.com/community/c6h1671327p10632648}{\textbf{2017 A7}} Let $a_0,a_1,a_2,\dots$ be a sequence of integers and $b_0,b_1,b_2,\dots$ be a sequence of positive integers such that $a_0=0,a_1=1$, and
\[
a_{n+1} =
        \begin{cases}
            a_nb_n+a_{n-1} & \text{if $b_{n-1}=1$} \\
            a_nb_n-a_{n-1} & \text{if $b_{n-1}>1$}
        \end{cases}\qquad\text{for }n=1,2,\dots.
\]for $n=1,2,\dots$.

Prove that at least one of the two numbers $a_{2017}$ and $a_{2018}$ must be greater than or equal to $2017$.
\item \href{https://artofproblemsolving.com/community/c6h1480730p8639356}{\textbf{2016 A7}} 	Find all functions $f:\mathbb{R}\rightarrow\mathbb{R}$ such that $f(0)\neq 0$ and for all $x,y\in\mathbb{R}$,
\[ f(x+y)^2 = 2f(x)f(y) + \max \left\{ f(x^2+y^2), f(x^2)+f(y^2) \right\}. \]
\item \href{https://artofproblemsolving.com/community/c6h546167p3160557}{\textbf{2012 A7}} We say that a function $f:\mathbb{R}^k \rightarrow \mathbb{R}$ is a metapolynomial if, for some positive integers $m$ and $n$, it can be represented in the form
\[f(x_1,\dots , x_k )=\max_{i=1,\dots , m} \min_{j=1,\dots , n}P_{i,j}(x_1,\dots , x_k),\]
where $P_{i,j}$ are multivariate polynomials. 

Prove that the product of two metapolynomials is also a metapolynomial.
\item \href{https://artofproblemsolving.com/community/c6h488538p2737646}{\textbf{2011 A7}} Let $a,b$ and $c$ be positive real numbers satisfying $\min(a+b,b+c,c+a) > \sqrt{2}$ and $a^2+b^2+c^2=3.$ Prove that

\[\frac{a}{(b+c-a)^2} + \frac{b}{(c+a-b)^2} + \frac{c}{(a+b-c)^2} \geqslant  \frac{3}{(abc)^2}.\] 
\newpage

\item \boxed{\href{https://artofproblemsolving.com/community/c6h356197p1936918}{\textbf{2010 A7}}}  Let $a_1, a_2, a_3, \dots$ be a sequence of positive real numbers, and $s$ be a positive integer, such that
\[a_n = \max \{ a_k + a_{n-k} \mid 1 \leqslant k \leqslant n-1 \} \ \textrm{ for all } \ n > s.\]
Prove there exist positive integers $\ell \leqslant s$ and $N$, such that
\[a_n = a_{\ell} + a_{n - \ell} \ \textrm{ for all } \ n \geqslant  N.\]

\item \href{https://artofproblemsolving.com/community/c6h355780p1932915}{\textbf{2009 A7}} Find all functions $f$ from the set of real numbers into the set of real numbers which satisfy for all $x$, $y$ the identity\[ f\left(xf(x+y)\right) = f\left(yf(x)\right) +x^2.\] 	

\item \href{https://artofproblemsolving.com/community/c6h287856p1555900}{\textbf{2008 A7}} Prove that for any four positive real numbers $ a$, $ b$, $ c$, $ d$ the inequality
\[ \frac {(a - b)(a - c)}{a + b + c} + \frac {(b - c)(b - d)}{b + c + d} + \frac {(c - d)(c - a)}{c + d + a} + \frac {(d - a)(d - b)}{d + a + b}\geqslant 0\]
holds. Determine all cases of equality. 

\item \boxed{\href{https://artofproblemsolving.com/community/c6h159989p894658}{\textbf{2007 A7}}} Let $ n$ be a positive integer. Consider
\[ S = \left\{ (x,y,z) \mid x,y,z \in \{ 0, 1, \dots, n\}, x + y + z > 0 \right \} \]
as a set of $ (n + 1)^{3} - 1$ points in the three-dimensional space.  Determine the smallest possible number of planes, the union of which contains $ S$ but does not include $(0,0,0)$. 

\end{itemize}
\newpage
\section*{A8 Problems}  \label{A8}
 \begin{itemize}[leftmargin=*]

 \item \href{https://www.imo-official.org/assets/documents/problems/2026/IMO2026SL.pdf}{\textbf{2026 A8}} TO BE REVEALED!
 
 \item \href{https://artofproblemsolving.com/community/c6h3866871p38676988}{\textbf{2025 A8}} Tim and Tam play a game. To start the game, Tim writes some nonzero real numbers not necessarily distinct, on a blackboard. In each round, the following happens.
\begin{itemize}[leftmargin=*]
    \item[$\bullet$]
First, Tam chooses a polynomial $P_k(x)=a_kx^k+a_{k-1}x^{k-1}+\cdots+a_1x+a_0$ whose coefficients $a_k,a_{k-1},\dots,a_1,a_0$ are all of the numbers currently written on the blackboard in some order. 

(For example, if the numbers on the blackboard, are $4,-3,4$, then $k=2$ and Tam may choose the polynomial $4x^2-3x+4$ or $-3x^2+4x+4$ but not $4x-3$ or $-3x^2-3x+4$.)
\item[$\bullet$] Then, if the equation $P_k(x)=0$ has no real solutions, the game is stopped. Otherwise, Tim chooses a real number $r$ such that $P_k(r)=0$ and writes it on the blackboard, so there is now one more number on the blackboard.

\end{itemize}
Determine whether Tim can ensure that the game is never stopped, no matter what Tam does.
\item \href{https://artofproblemsolving.com/community/c6h3610460p35340944}{\textbf{2024 A8}} Let $p \ne q$ be coprime positive integers. Determine all infinite sequences $a_1$, $a_2$, $\dots$ of positive integers such that the following conditions hold for all $n \geqslant 1$:
\begin{align*}
    \max(a_n, a_{n+1}, \dots, a_{n+p}) - \min(a_n, a_{n+1}, \dots, a_{n+p}) &= p \\
    \max(a_n, a_{n+1}, \dots, a_{n+q}) - \min(a_n, a_{n+1}, \dots, a_{n+q}) &= q   
\end{align*}
    \item \href{https://artofproblemsolving.com/community/c6h3107360p28104441}{\textbf{2022 A8}} For a positive integer $n$, an $n$-sequence is a sequence $(a_0,\dots,a_n)$ of non-negative integers satisfying the following condition: if $i$ and $j$ are non-negative integers with $i+j \leqslant n$, then $a_i+a_j \leqslant n$ and $a_{a_i+a_j}=a_{i+j}$. Let $f(n)$ be the number of $n$-sequences. 
    
    Prove that there exist positive real numbers $c_1$, $c_2$, and $\lambda$ such that\[c_1\lambda^n<f(n)<c_2\lambda^n\]for all positive integers $n$.
    
\item \href{https://artofproblemsolving.com/community/c6h2882529p25627449}{\textbf{2021 A8}} Determine all functions $f: \mathbb{R} \rightarrow \mathbb{R}$ that satisfy$$(f(a)-f(b))(f(b)-f(c))(f(c)-f(a)) = f(ab^2+bc^2+ca^2) - f(a^2b+b^2c+c^2a)$$for all real numbers $a$, $b$, $c$.

\item \href{https://artofproblemsolving.com/community/c6h2625874p22698175}{\textbf{2020 A8}} Let $\mathbb{R}^+$ be the set of positive real numbers. Determine all functions $f:\mathbb{R}^+$ $\rightarrow$ $\mathbb{R}^+$ such that for all positive real numbers $x$ and $y$,
$$f(x+f(xy))+y=f(x)f(y)+1.$$

\item \href{https://artofproblemsolving.com/community/c6h1671299p10632376}{\textbf{2017 A8}} A function $f:\mathbb{R} \to \mathbb{R}$ has the following property:
$$\text{For every } x,y \in \mathbb{R} \text{ such that }(f(x)+y)(f(y)+x) > 0, \text{ we have } f(x)+y = f(y)+x.$$Prove that $f(x)+y \leqslant f(y)+x$ whenever $x>y$.

\item \href{https://artofproblemsolving.com/community/c6h1480693p8639259}{\textbf{2016 A8}} 	Find the largest real constant $a$ such that for all $n \geqslant  1$ and for all real numbers $x_0, x_1, \dots , x_n$ satisfying $0 = x_0 < x_1 < x_2 < \dots < x_n$ we have
\[\frac{1}{x_1-x_0} + \frac{1}{x_2-x_1} + \dots + \frac{1}{x_n-x_{n-1}} \geqslant  a \left( \frac{2}{x_1} + \frac{3}{x_2} + \dots + \frac{n+1}{x_n} \right).\]

\item \href{https://artofproblemsolving.com/community/c6h418684p2362291}{\textbf{2010 A8}} Given six positive numbers $a,b,c,d,e,f$ such that $a < b < c < d < e < f.$ Let $a+c+e=S$ and $b+d+f=T.$ Prove that
\[2ST > \sqrt{3(S+T)\left(S(bd + bf + df) + T(ac + ae + ce) \right)}.\]
\end{itemize}

\newpage

\section*{A9 Problems} \label{A9}
\begin{itemize}[leftmargin=*]
    \item \href{https://www.imo-official.org/assets/documents/problems/2026/IMO2026SL.pdf}{\textbf{2026 A9}} TO BE REVEALED!
\end{itemize}

\newpage
\title{Combinatorics IMO Shortlist}

\section*{C1 Problems} \label{C1}
\begin{itemize}[leftmargin=*]


\item \boxed{\href{https://www.imo-official.org/assets/documents/problems/2026/IMO2026SL.pdf}{\textbf{2026 C1}}} Shan-Yu and Mulan are playing a game. Let $\theta$ be an angle with $0^\circ<\theta<180^\circ$ known to both players. Initially, Shan-Yu makes a paper triangle $\mathcal{T}$ with measurements of his choice. Then, they repeatedly perform the following steps:
\begin{itemize}[leftmargin=*]
\item[$\bullet$] If $\mathcal{T}$ has at least one angle measuring exactly $\theta$, then the game stops and Mulan wins.
\item[$\bullet$] Otherwise, Mulan chooses a point $P$ on the perimeter of $\mathcal{T}$, different from its three vertices. She then makes a straight cut from $P$ to the opposite vertex of $\mathcal{T}$, splitting it into two triangles.
\item[$\bullet$] Shan-Yu discards one of the two triangles. The remaining triangle becomes the new $\mathcal{T}$.
\end{itemize}
For which real values of $\theta$ can Mulan guarantee her victory in finitely many steps, no matter how Shan-Yu plays?

\item \boxed{\href{https://artofproblemsolving.com/community/c6h3609787p35332003}{\textbf{2025 C1}}} A line in the plane is called $sunny$ if it is \textbf{not} parallel to any of the $x$-axis, the $y$-axis, or the line $x+y=0$.

Let $n\geqslant  3$ be a given integer. Determine all nonnegative integers $k$ such that there exist $n$ distinct lines in the plane satisfying both of the following:
\begin{itemize}[leftmargin=*]
\item[$\bullet$]  for all positive integers $a$ and $b$ with $a+b\leqslant n+1$, the point $(a,b)$ lies on at least one of the lines; and
\item[$\bullet$] exactly $k$ of the $n$ lines are sunny.
\end{itemize}
\item \href{https://artofproblemsolving.com/community/c6h3610442p35340913}{\textbf{2024 C1}} Let $n$ be a positive integer. A class of $n$ students run $n$ races, in each of which they are ranked with no draws. A student is eligible for a rating $(a,\,b)$ for positive integers $a$ and $b$ if they come in the top $b$ places in at least $a$ of the races. Their final score is the maximum possible value of $a-b$ across all ratings for which they are eligible.

Find the maximum possible sum of all the scores of the $n$ students.
\item \href{https://artofproblemsolving.com/community/c6h3359749p31218491}{\textbf{2023 C1}} Let $m$ and $n$ be positive integers greater than $1$. In each unit square of an $m\times n$ grid lies a coin with its tail side up. A move consists of the following steps.
\begin{enumerate}[1.]
\item select a $2\times 2$ square in the grid;
\item flip the coins in the top-left and bottom-right unit squares;
\item flip the coin in either the top-right or bottom-left unit square.
 
\end{enumerate}
Determine all pairs $(m,n)$ for which it is possible that every coin shows head-side up after a finite number of moves.

    \item \href{https://artofproblemsolving.com/community/c6h3107334p28104276}{\textbf{2022 C1}} 	A $\pm 1$-sequence is a sequence of $2022$ numbers $a_1, \dots, a_{2022},$ each equal to either $+1$ or $-1$. Determine the largest $C$ so that, for any $\pm 1$-sequence, there exists an integer $k$ and indices $1 \leqslant t_1 < \dots < t_k \leqslant 2022$ so that $t_{i+1} - t_i \leqslant 2$ for all $i$, and$$\left| \sum_{i = 1}^{k} a_{t_i} \right| \geqslant C.$$
    
    \item \href{https://artofproblemsolving.com/community/c6h2882547p25627520}{\textbf{2021 C1}} Let $S$ be an infinite set of positive integers, such that there exist four pairwise distinct $a,b,c,d \in S$ with $\gcd(a,b) \neq \gcd(c,d)$. 
    
    Prove that there exist three pairwise distinct $x,y,z \in S$ such that $\gcd(x,y)=\gcd(y,z) \neq \gcd(z,x)$.
    
    \item \href{https://artofproblemsolving.com/community/c6h2625893p22698309}{\textbf{2020 C1}} Let $n$ be a positive integer. Find the number of permutations $a_1, a_2,\dots a_n$ of the
sequence $1, 2, \dots , n$ satisfying
$$a_1 \leqslant 2a_2\leqslant 3a_3 \leqslant \dots \leqslant na_n.$$

    \item \href{https://artofproblemsolving.com/community/c6h2279044p17829263}{\textbf{2019 C1}} The infinite sequence $a_0,a _1, a_2, \dots$ of (not necessarily distinct) integers has the following properties: $0\leqslant a_i \leqslant i$ for all integers $i\geqslant 0$, and\[\binom{k}{a_0} + \binom{k}{a_1} + \dots + \binom{k}{a_k} = 2^k\]for all integers $k\geqslant 0$. Prove that all integers $N\geqslant 0$ occur in the sequence (that is, for all $N\geqslant 0$, there exists $i\geqslant 0$ with $a_i=N$).

    \item \href{https://artofproblemsolving.com/community/c6h1876759p12752812}{\textbf{2018 C1}} 	Let $n\geqslant  3$ be an integer. Prove that there exists a set $S$ of $2n$ positive integers satisfying the following property: For every $m=2,3,\dots,n$ the set $S$ can be partitioned into two subsets with equal sums of elements, with one of subsets of cardinality $m$.
    
    \item \href{https://artofproblemsolving.com/community/c6h1671298p10632375}{\textbf{2017 C1}} A rectangle $\mathcal{R}$ with odd integer side lengths is divided into small rectangles with integer side lengths. Prove that there is at least one among the small rectangles whose distances from the four sides of $\mathcal{R}$ are either all odd or all even.
    
    \item \href{https://artofproblemsolving.com/community/c6h1480687p8639246}{\textbf{2016 C1}} The leader of an IMO team chooses positive integers $n$ and $k$ with $n>k$, and announces them to the deputy leader and a contestant. The leader then secretly tells the deputy leader an $n$-digit binary string, and the deputy leader writes down all $n$-digit binary strings which differ from the leader's in exactly $k$ positions. (For example, if $n = 3$ and $k = 1$, and if the leader chooses $101$, the deputy leader would write down $001, 111$ and $100$.) The contestant is allowed to look at the strings written by the deputy leader and guess the leader's string.
    
    What is the minimum number of guesses (in terms of $n$ and $k$) needed to guarantee the correct answer?
    \item \href{https://artofproblemsolving.com/community/c6h1268873p6622370}{\textbf{2015 C1}} In Lineland there are $n\geqslant 1$ towns, arranged along a road running from left to right. Each town has a left bulldozer (put to the left of the town and facing left) and a right bulldozer (put to the right of the town and facing right). The sizes of the $2n$ bulldozers are distinct. Every time when a left and right bulldozer confront each other, the larger bulldozer pushes the smaller one off the road. On the other hand, bulldozers are quite unprotected at their rears; so, if a bulldozer reaches the rear-end of another one, the first one pushes the second one off the road, regardless of their sizes. 
    
    Let $A$ and $B$ be two towns, with $B$ to the right of $A$. We say that town $A$ can sweep town $B$ away if the right bulldozer of $A$ can move over to $B$ pushing off all bulldozers it meets. Similarly town $B$ can sweep town $A$ away if the left bulldozer of $B$ can move over to $A$ pushing off all bulldozers of all towns on its way. 
    
    Prove that there is exactly one town that cannot be swept away by any other one.
    \item \href{https://artofproblemsolving.com/community/c6h1113183p5083543}{\textbf{2014 C1}} Let $n$ points be given inside a rectangle $R$ such that no two of them lie on a line parallel to one of the sides of $R$. The rectangle $R$ is to be dissected into smaller rectangles with sides parallel to the sides of $R$ in such a way that none of these rectangles contains any of the given points in its interior. 
    
    Prove that we have to dissect $R$ into at least $n + 1$ smaller rectangles.

    \newpage
    
    \item \href{https://artofproblemsolving.com/community/c6h597127p3543383}{\textbf{2013 C1}} Let $n$ be an positive integer. Find the smallest integer $k$ with the following property; Given any real numbers $a_1 , \dots , a_d $ such that $a_1 + a_2 + \dots + a_d = n$ and $0 \leqslant a_i \leqslant 1$ for $i=1,2,\dots ,d$, it is possible to partition these numbers into $k$ groups (some of which may be empty) such that the sum of the numbers in each group is at most $1$.
    \item \href{https://artofproblemsolving.com/community/c6h546168p3160559}{\textbf{2012 C1}} Several positive integers are written in a row. Iteratively, Alice chooses two adjacent numbers $x$ and $y$ such that $x>y$ and $x$ is to the left of $y$, and replaces the pair $(x,y)$ by either $(y+1,x)$ or $(x-1,x)$. Prove that she can perform only finitely many such iterations.
    \item \boxed{\href{https://artofproblemsolving.com/community/c6h418978p2365036}{\textbf{2011 C1}}} Let $n > 0$ be an integer. We are given a balance and $n$ weights of weight $2^0, 2^1, \dots, 2^{n-1}$. We are to place each of the $n$ weights on the balance, one after another, in such a way that the right pan is never heavier than the left pan. At each step we choose one of the weights that has not yet been placed on the balance, and place it on either the left pan or the right pan, until all of the weights have been placed. Determine the number of ways in which this can be done.
    \item \href{https://artofproblemsolving.com/community/c6h418685p2362292}{\textbf{2010 C1}} 	In a concert, $20$ singers will perform. For each singer, there is a (possibly empty) set of other singers such that he wishes to perform later than all the singers from that set. 
    
    Can it happen that there are exactly $2010$ orders of the singers such that all their wishes are satisfied?
    \item \href{https://artofproblemsolving.com/community/c6h355783p1932923}{\textbf{2009 C1}} Consider $2009$ cards, each having one gold side and one black side, lying on parallel on a long table. Initially all cards show their gold sides. Two player, standing by the same long side of the table, play a game with alternating moves. Each move consists of choosing a block of $50$ consecutive cards, the leftmost of which is showing gold, and turning them all over, so those which showed gold now show black and vice versa. The last player who can make a legal move wins. 
    \begin{enumerate}[(a)]
        \item Does the game necessarily end?
        \item Does there exist a winning strategy for the starting player?
    \end{enumerate}
    \item \href{https://artofproblemsolving.com/community/c6h287859p1555905}{\textbf{2008 C1}} 	In the plane we consider rectangles whose sides are parallel to the coordinate axes and have positive length. Such a rectangle will be called a box. Two boxes intersect if they have a common point in their interior or on their boundary. 
    
    Find the largest $ n$ for which there exist $ n$ boxes $ B_1$, $ \dots$, $ B_n$ such that $ B_i$ and $ B_j$ intersect if and only if $ i\not\equiv j\pm 1\pmod n$.
    \item \href{https://artofproblemsolving.com/community/c6h214707p1187174}{\textbf{2007 C1}} 	Let $ n > 1$ be an integer. Find all sequences $ a_1, a_2, \dots , a_{n^2 + n}$ satisfying the following conditions:
    \begin{enumerate}[(a)]
        \item $a_i \in \left\{0,1\right\} \text{ for all } 1 \leqslant i \leqslant n^2 + n;$
        \item $a_{i + 1} + a_{i + 2} + \dots + a_{i + n} < a_{i + n + 1} + a_{i + n + 2} + \dots + a_{i + 2n} \text{ for all } 0 \leqslant i \leqslant n^2 - n.$
    \end{enumerate}
    \item \href{https://artofproblemsolving.com/community/c6h155692p874978}{\textbf{2006 C1}} We have $ n \geqslant  2$ lamps $ L_{1}, \dots ,L_{n}$ in a row, each of them being either on or off. Every second we simultaneously modify the state of each lamp as follows: if the lamp $ L_{i}$ and its neighbours (only one neighbour for $ i = 1$ or $ i = n$, two neighbours for other $ i$) are in the same state, then $ L_{i}$ is switched off; otherwise, $ L_{i}$ is switched on. Initially all the lamps are off except the leftmost one which is on. 
    \begin{enumerate}[(a)]
        \item Prove that there are infinitely many integers $n$ for which all the lamps will eventually be off.
        \item Prove that there are infinitely many integers $n$ for which the lamps will never be all off.
    \end{enumerate}
\end{itemize}
\newpage

\section*{C2 Problems} \label{C2}
\begin{itemize}[leftmargin=*]


\item \href{https://www.imo-official.org/assets/documents/problems/2026/IMO2026SL.pdf}{\textbf{2026 C2}} TO BE REVEALED!

\item \href{https://artofproblemsolving.com/community/c6h3866876p38676996}{\textbf{2025 C2}}
There is a row of $n$ paddocks, labelled from left to right with the integers $1$ to $n$, where $n\geqslant 3$. Skippy the Kangaroo grazes in the paddock according to the following rule:

If Skippy is grazing in paddock $k$, then she makes a sequence of $k$ hops from a paddock to an adjacent paddock. The first of the $k$ hops is always to the right, unless Skippy is in paddock $n$, in which case it is to the left. Each of the following $k-1$ hops is in the same direction as the previous hop, unless Skippy is in paddock $1$ or $n$. Skippy grazes in the paddock that she is in after the $k^{\text{th}}$ hop.

For example, if $n=8$ and Skippy is grazing in paddock $3$, the next $4$ paddocks she will graze in are $6$, $4$, $8$, and $2$, in this order. 

Skippy continues grazing in this way indefinitely. A paddock is $overgrazed$ if Skippy eventually grazes in it, regardless of which paddock she starts in. Determine all $n\geqslant 3$ for which there exists exactly one overgrazed paddock.
\item \href{https://artofproblemsolving.com/community/c6h3610436p35340903}{\textbf{2024 C2}} Let $n$ be a positive integer. The integers $1,\,2,\,3,\,\dots,\,n^2$ are to be written in the cells of an $n\times n$ board such that each integer is written in exactly one cell and each cell contains exactly one integer. For every integer $d$ with $d\mid n$, the $d$-division of the board is the division of the board into $(n/d)^2$ nonoverlapping sub-boards, each of size $d\times d$, such that each cell is contained in exactly one $d\times d$ sub-board.

We say that $n$ is a cool number if the integers can be written on the $n\times n$ board such that, for each integer $d$ with $d\mid n$ and $1<d<n$, in the $d$-division of the board, the sum of the integers written in each $d\times d$ sub-board is not a multiple of $d$.

Determine all even cool numbers.
\item \href{https://artofproblemsolving.com/community/c6h3359755p31218537}{\textbf{2023 C2}} Determine the maximal length $L$ of a sequence $a_1,\dots,a_L$ of positive integers satisfying both the following properties:
\begin{itemize}[leftmargin=*]
\item[$\bullet$] every term in the sequence is less than or equal to $2^{2023}$, and
\item[$\bullet$] there does not exist a consecutive subsequence $a_i,a_{i+1},\dots,a_j$ (where $1\leqslant i\leqslant j\leqslant L$) with a choice of signs $s_i,s_{i+1},\dots,s_j\in\{1,-1\}$ for which\[s_ia_i+s_{i+1}a_{i+1}+\dots+s_ja_j=0.\]
\end{itemize}
    \item \boxed{\href{https://artofproblemsolving.com/community/c6h2883211p25635135}{\textbf{2022 C2}}} The Bank of Oslo issues two types of coin: aluminum (denoted A) and bronze (denoted B). Marianne has $n$ aluminum coins and $n$ bronze coins arranged in a row in some arbitrary initial order. A chain is any subsequence of consecutive coins of the same type. Given a fixed positive integer $k \leqslant 2n$, Gilberty repeatedly performs the following operation: he identifies the longest chain containing the $k^{th}$ coin from the left and moves all coins in that chain to the left end of the row. For example, if $n=4$ and $k=4$, the process starting from the ordering $AABBBABA$ would be \[AABBBABA \to BBBAAABA \to AAABBBBA \to BBBBAAAA \to \dots\] 
    Find all pairs $(n,k)$ with $1 \leqslant k \leqslant 2n$ such that for every initial ordering, at some moment during the process, the leftmost $n$ coins will all be of the same type.
    \item \href{https://artofproblemsolving.com/community/c6h2882538p25627501}{\textbf{2021 C2}}  Let $n\geqslant 3$ be an integer. An integer $m\geqslant n+1$ is called $n-colourful$ if,given infinitely many marbles in each of $n$ colours $C_1,C_2,\dots,C_n$, it is possible to place $m$ of them around circle so that in any group of $n+1$ consecutive marbles there is at least one marble of colour $C_i$ for each $i=1,2,\dots,n$.

    Prove that there are only finitely many positive integers which are not $n-colourful$. Find the largest among them.
    
    \item \href{https://artofproblemsolving.com/community/c6h2625840p22697894}{\textbf{2020 C2}}  In a regular 100-gon, 41 vertices are colored black and the remaining 59 vertices are colored white. Prove that there exist 24 convex quadrilaterals $Q_{1}, \dots, Q_{24}$ whose corners are vertices of the 100-gon, so that
    \begin{enumerate}[(a)]
        \item the quadrilaterals $Q_{1}, \dots, Q_{24}$ are pairwise disjoint, and
        \item every quadrilateral $Q_{i}$ has three corners of one color and one corner of the other color.
    \end{enumerate}
    \item \href{https://artofproblemsolving.com/community/c6h2279034p17829157}{\textbf{2019 C2}} You are given a set of $n$ blocks, each weighing at least $1$; their total weight is $2n$. Prove that for every real number $r$ with $0 \leqslant r \leqslant 2n-2$ you can choose a subset of the blocks whose total weight is at least $r$ but at most $r + 2$.
    \item \boxed{\href{https://artofproblemsolving.com/community/c6h1671290p10632348}{\textbf{2018 C2}}} A site is any point $(x, y)$ in the plane such that $x$ and $y$ are both positive integers less than or equal to 20. Initially, each of the 400 sites is unoccupied. Amy and Ben take turns placing stones with Amy going first. On her turn, Amy places a new red stone on an unoccupied site such that the distance between any two sites occupied by red stones is not equal to $\sqrt{5}$. On his turn, Ben places a new blue stone on any unoccupied site. (A site occupied by a blue stone is allowed to be at any distance from any other occupied site.) They stop as soon as a player cannot place a stone. 
    
    Find the greatest $K$ such that Amy can ensure that she places at least $K$ red stones, no matter how Ben places his blue stones.
    \item \href{https://artofproblemsolving.com/community/c6h1671304p10632393}{\textbf{2017 C2}} Let $n$ be a positive integer. Define a chameleon to be any sequence of $3n$ letters, with exactly $n$ occurrences of each of the letters $a, b,$ and $c$. Define a swap to be the transposition of two adjacent letters in a chameleon. 
    
    Prove that for any chameleon $X$ , there exists a chameleon $Y$ such that $X$ cannot be changed to $Y$ using fewer than $\dfrac{3n^2}{2}$ swaps.
    \item \href{https://artofproblemsolving.com/community/c6h1480692p8639256}{\textbf{2016 C2}} Find all positive integers $n$ for which all positive divisors of $n$ can be put into the cells of a rectangular table under the following constraints:
    \begin{enumerate}[(a)]
        \item each cell contains a distinct divisor;
        \item the sums of all rows are equal; and
        \item the sums of all columns are equal.
    \end{enumerate}
    \item \boxed{\href{https://artofproblemsolving.com/community/c6h1112753p5079689}{\textbf{2015 C2}}} We say that a finite set $\mathcal{S}$ of points in the plane is balanced if, for any two different points $A$ and $B$ in $\mathcal{S}$, there is a point $C$ in $\mathcal{S}$ such that $AC=BC$. We say that $\mathcal{S}$ is centre-free if for any three different points $A$, $B$ and $C$ in $\mathcal{S}$, there is no points $P$ in $\mathcal{S}$ such that $PA=PB=PC$. 
    \begin{enumerate}[(a)]
        \item Show that for all integers $n\geqslant 3$, there exists a balanced set consisting of $n$ points.
        \item Determine all integers $n\geqslant 3$ for which there exists a balanced centre-free set consisting of $n$ points.
    \end{enumerate}
    \item \href{https://artofproblemsolving.com/community/c6h1113184p5083546}{\textbf{2014 C2}} We have $2^m$ sheets of paper, with the number $1$ written on each of them. We perform the following operation. In every step we choose two distinct sheets; if the numbers on the two sheets are $a$ and $b$, then we erase these numbers and write the number $a + b$ on both sheets. 
    
    Prove that after $m2^{m -1}$ steps, the sum of the numbers on all the sheets is at least $4^m$.
    \item \boxed{\href{https://artofproblemsolving.com/community/c6h1181527p5720110}{\textbf{2013 C2}}} A configuration of $4027$ points in the plane is called Colombian if it consists of $2013$ red points and $2014$ blue points, and no three of the points of the configuration are collinear. By drawing some lines, the plane is divided into several regions. An arrangement of lines is good for a Colombian configuration if the following two conditions are satisfied: 
    \begin{enumerate}[(a)]
        \item No line passes through any point of the configuration.
        \item No region contains points of both colors.
    \end{enumerate} 
    Find the least value of $k$ such that for any Colombian configuration of $4027$ points, there is a good arrangement of $k$ lines.
    \item \href{https://artofproblemsolving.com/community/c6h546169p3160560}{\textbf{2012 C2}} Let $n \geqslant  1$ be an integer. What is the maximum number of disjoint pairs of elements of the set $\{ 1,2,\dots , n \}$ such that the sums of the different pairs are different integers not exceeding $n$?
    \item \href{https://artofproblemsolving.com/community/c6h488626p2738008}{\textbf{2011 C2}} Suppose that $1000$ students are standing in a circle. Prove that there exists an integer $k$ with $100 \leqslant k \leqslant 300$ such that in this circle there exists a contiguous group of $2k$ students, for which the first half contains the same number of girls as the second half.
    \item \href{https://artofproblemsolving.com/community/c6h418686p2362296}{\textbf{2010 C2}} On some planet, there are $2^N$ countries $(N \geqslant  4).$ Each country has a flag $N$ units wide and one unit high composed of $N$ fields of size $1 \times 1,$ each field being either yellow or blue. No two countries have the same flag. We say that a set of $N$ flags is diverse if these flags can be arranged into an $N \times N$ square so that all $N$ fields on its main diagonal will have the same color. 
    
    Determine the smallest positive integer $M$ such that among any $M$ distinct flags, there exist $N$ flags forming a diverse set.
    \item \href{https://artofproblemsolving.com/community/c6h355784p1932924}{\textbf{2009 C2}} For any integer $n\geqslant  2$, let $N(n)$ be the maximal number of triples $(a_i, b_i, c_i)$, $i=1, \dots, N(n)$, consisting of nonnegative integers $a_i$, $b_i$ and $c_i$ such that the following two conditions are satisfied:
    \begin{enumerate}[(a)]
        \item $a_i+b_i+c_i=n$ for all $i=1, \dots, N(n)$,
        \item If $i\neq j$ then $a_i\neq a_j$, $b_i\neq b_j$ and $c_i\neq c_j$
    \end{enumerate} 
    Determine $N(n)$ for all $n\geqslant  2$.
    
    \item \href{https://artofproblemsolving.com/community/c6h271833p1472110}{\textbf{2008 C2}} Let $n \in \mathbb N$ and $A_n$ set of all permutations $(a_1, \dots, a_n)$ of the set $\{1, 2, \dots , n\}$ for which
\[k|2(a_1 + \dots+ a_k), \text{ for all } 1 \leqslant k \leqslant n.\]
Find the number of elements of the set $A_n$.
    \item \href{https://artofproblemsolving.com/community/c6h195050p1071295}{\textbf{2007 C2}} A rectangle $ D$ is partitioned in several ($ \geqslant 2$) rectangles with sides parallel to those of $ D$. Given that any line parallel to one of the sides of $ D$, and having common points with the interior of $ D$, also has common interior points with the interior of at least one rectangle of the partition; prove that there is at least one rectangle of the partition having no common points with $ D$'s boundary.
    \newpage
    \item \boxed{\href{https://artofproblemsolving.com/community/c6h101306p571973}{\textbf{2006 C2}}} Let $P$ be a regular $2006$-gon. A diagonal is called good if its endpoints divide the boundary of $P$ into two parts, each composed of an odd number of sides of $P$. The sides of $P$ are also called good. Suppose $P$ has been dissected into triangles by $2003$ diagonals, no two of which have a common point in the interior of $P$.
    
    Find the maximum number of isosceles triangles having two good sides that could appear in such a configuration.
\end{itemize}
\newpage

\section*{C3 Problems} \label{C3}
\begin{itemize}[leftmargin=*]


\item \href{https://www.imo-official.org/assets/documents/problems/2026/IMO2026SL.pdf}{\textbf{2026 C3}} TO BE REVEALED!

\item \href{https://artofproblemsolving.com/community/c6h3866891p38677080}{\textbf{2025 C3}} There are $2025$ white balls and $2025$ black balls arranged in a row. A $balanced$ segment is a contiguous nonempty segment of balls that contains the same number of white balls as black balls. A $balanced$ $removal$ is the operation of selecting a balanced segment $S$, removing all the balls in $S$, and shifting left the balls that were to the right of $S$ so that the remaining balls form a contiguous row.
    
Determine the smallest positive integer $N$ satisfying the following property: for every configuration of balls and for every integer $k$ satisfying $0 \leqslant k \leqslant 2025$, there exists a sequence of at most $N$ balanced removals after which there are exactly $2k$ remaining balls.
\item \href{https://artofproblemsolving.com/community/c6h3610441p35340911}{\textbf{2024 C3}} Let $n$ be a positive integer. There are $2n$ knights sitting at a round table. They consist of $n$ pairs of partners, each pair of which wishes to shake hands. A pair can shake hands only when next to each other. Every minute, one pair of adjacent knights swaps places. 

Find the minimum number of exchanges of adjacent knights such that, regardless of the initial arrangement, every knight can meet her partner and shake hands at some time.
\item \boxed{\href{https://artofproblemsolving.com/community/c6h3107350p28104367}{\textbf{2023 C3}} } Let $n$ be a positive integer. A Japanese triangle consists of $1 + 2 + \dots + n$ circles arranged in an equilateral triangular shape such that for each $i = 1$, $2$, $\dots$, $n$, the $i^{th}$ row contains exactly $i$ circles, exactly one of which is coloured red. A ninja path in a Japanese triangle is a sequence of $n$ circles obtained by starting in the top row, then repeatedly going from a circle to one of the two circles immediately below it and finishing in the bottom row. Here is an example of a Japanese triangle with $n = 6$, along with a ninja path in that triangle containing two red circles.
\begin{center}
\begin{asy}
// credit to vEnhance for the diagram (which was better than my original asy):
size(4cm);
  pair X = dir(240); pair Y = dir(0);
  path c = scale(0.5)*unitcircle;
  int[] t = {0,0,2,2,3,0};
  for (int i=0; i<=5; ++i) {
    for (int j=0; j<=i; ++j) {
      filldraw(shift(i*X+j*Y)*c, (t[i]==j) ? lightred : white);
      draw(shift(i*X+j*Y)*c);
    }
  }
  draw((0,0)--(X+Y)--(2*X+Y)--(3*X+2*Y)--(4*X+2*Y)--(5*X+2*Y),linewidth(1.5));
  path q = (3,-3sqrt(3))--(-3,-3sqrt(3));
  draw(q,Arrows(TeXHead, 1));
  label("$n = 6$", q, S);
label("$n = 6$", q, S);
\end{asy}
\end{center}
In terms of $n$, find the greatest $k$ such that in each Japanese triangle there is a ninja path containing at least $k$ red circles.
\item \href{https://artofproblemsolving.com/community/c6h3107325p28104250}{\textbf{2022 C3}} In each square of a garden shaped like a $2022 \times 2022$ board, there is initially a tree of height $0$. A gardener and a lumberjack alternate turns playing the following game, with the gardener taking the first turn:
\begin{enumerate}[(a)]
\item  The gardener chooses a square in the garden. Each tree on that square and all the surrounding squares of which there are at most eight) then becomes one unit taller.
\item  The lumberjack then chooses four different squares on the board. Each tree of positive height on those squares then becomes one unit shorter.
\end{enumerate}
We say that a tree is majestic if its height is at least $10^6$. Determine the largest $K$ such that the gardener can ensure there are eventually $K$ majestic trees on the board, no matter how the lumberjack plays. 
\item \boxed{\href{https://artofproblemsolving.com/community/c6h2625845p22697921}{\textbf{2021 C3}}} Two squirrels, Bushy and Jumpy, have collected 2021 walnuts for the winter. Jumpy numbers the walnuts from 1 through 2021, and digs 2021 little holes in a circular pattern in the ground around their favourite tree. The next morning Jumpy notices that Bushy had placed one walnut into each hole, but had paid no attention to the numbering. Unhappy, Jumpy decides to reorder the walnuts by performing a sequence of 2021 moves. In the $k$-th move, Jumpy swaps the positions of the two walnuts adjacent to walnut $k$.

Prove that there exists a value of $k$ such that, on the $k$-th move, Jumpy swaps some walnuts $a$ and $b$ such that $a<k<b$. 
\item \boxed{\href{https://artofproblemsolving.com/community/c6h2278648p17821585}{\textbf{2020 C3}}} There is an integer $n > 1$. There are $n^2$ stations on a slope of a mountain, all at different altitudes. Each of two cable car companies, $A$ and $B$, operates $k$ cable cars; each cable car provides a transfer from one of the stations to a higher one (with no intermediate stops). The $k$ cable cars of $A$ have $k$ different starting points and $k$ different finishing points, and a cable car which starts higher also finishes higher. The same conditions hold for $B$. We say that two stations are linked by a company if one can start from the lower station and reach the higher one by using one or more cars of that company (no other movements between stations are allowed). 

Determine the smallest positive integer $k$ for which one can guarantee that there are two stations that are linked by both companies. 
\item \boxed{\href{https://artofproblemsolving.com/community/c6h1876772p12752847}{\textbf{2019 C3}}}  The Bank of Bath issues coins with an $H$ on one side and a $T$ on the other. Harry has $n$ of these coins arranged in a line from left to right. He repeatedly performs the following operation: if there are exactly $k>0$ coins showing $H$, then he turns over the $k$th coin from the left; otherwise, all coins show $T$ and he stops. For example, if $n=3$ the process starting with the configuration $THT$ would be $THT \to HHT  \to HTT \to TTT$, which stops after three operations.\\
(a) Show that, for each initial configuration, Harry stops after a finite number of operations.\\
(b) For each initial configuration $C$, let $L(C)$ be the number of operations before Harry stops. For example, $L(THT) = 3$ and $L(TTT) = 0$.
Determine the average value of $L(C)$ over all $2^n$ possible initial configurations $C$. 
\item \href{https://artofproblemsolving.com/community/c6h1876784p12752893}{\textbf{2018 C3}} 	Let $n$ be a given positive integer. Sisyphus performs a sequence of turns on a board consisting of $n + 1$ squares in a row, numbered $0$ to $n$ from left to right. Initially, $n$ stones are put into square $0$, and the other squares are empty. At every turn, Sisyphus chooses any nonempty square, say with $k$ stones, takes one of these stones and moves it to the right by at most $k$ squares (the stone should say within the board). Sisyphus' aim is to move all $n$ stones to square $n$.
Prove that Sisyphus cannot reach the aim in less than
\[ \left \lceil \frac{n}{1} \right \rceil + \left \lceil \frac{n}{2} \right \rceil + \left \lceil \frac{n}{3} \right \rceil + \dots + \left \lceil \frac{n}{n} \right \rceil. \] 
\item \href{https://artofproblemsolving.com/community/c6h1671318p10632509}{\textbf{2017 C3}} Sir Alex plays the following game on a row of 9 cells. Initially, all cells are empty. In each move, Sir Alex is allowed to perform exactly one of the following two operations:
Choose any number of the form $2^j$, where $j$ is a non-negative integer, and put it into an empty cell.
Choose two (not necessarily adjacent) cells with the same number in them; denote that number by $2^j$. Replace the number in one of the cells with $2^{j+1}$ and erase the number in the other cell.
At the end of the game, one cell contains $2^n$, where $n$ is a given positive integer, while the other cells are empty. 

Determine the maximum number of moves that Sir Alex could have made, in terms of $n$. 

\newpage

\item \href{https://artofproblemsolving.com/community/c6h1480694p8639260}{\textbf{2016 C3}} Let $n$ be a positive integer relatively prime to $6$. We paint the vertices of a regular $n$-gon with three colours so that there is an odd number of vertices of each colour. 

Show that there exists an isosceles triangle whose three vertices are of different colours. 

\item \href{https://artofproblemsolving.com/community/c6h1268855p6622233}{\textbf{2015 C3}} For a finite set $A$ of positive integers, a partition of $A$ into two disjoint nonempty subsets $A_1$ and $A_2$ is $\textit{good}$ if the least common multiple of the elements in $A_1$ is equal to the greatest common divisor of the elements in $A_2$. Determine the minimum value of $n$ such that there exists a set of $n$ positive integers with exactly $2015$ good partitions. 

\item \boxed{\href{https://artofproblemsolving.com/community/c6h596929p3542094}{\textbf{2014 C3}}} 	Let $n \geqslant 2$ be an integer. Consider an $n \times n$ chessboard consisting of $n^2$ unit squares. A configuration of $n$ rooks on this board is peaceful if every row and every column contains exactly one rook. Find the greatest positive integer $k$ such that, for each peaceful configuration of $n$ rooks, there is a $k \times k$ square which does not contain a rook on any of its $k^2$ unit squares. 
\item \href{https://artofproblemsolving.com/community/c6h597130p3543398}{\textbf{2013 C3}} A crazy physicist discovered a new kind of particle wich he called an imon, after some of them mysteriously appeared in his lab. Some pairs of imons in the lab can be entangled, and each imon can participate in many entanglement relations. The physicist has found a way to perform the following two kinds of operations with these particles, one operation at a time.
\begin{enumerate}[(a)]
    \item  If some imon is entangled with an odd number of other imons in the lab, then the physicist can destroy it.
    \item  At any moment, he may double the whole family of imons in the lab by creating a copy $I'$ of each imon $I$. During this procedure, the two copies $I'$ and $J'$ become entangled if and only if the original imons $I$ and $J$ are entangled, and each copy $I'$ becomes entangled with its original imon $I$; no other entanglements occur or disappear at this moment.
\end{enumerate}
Prove that the physicist may apply a sequence of such operations resulting in a family of imons, no two of which are entangled. 
\item \href{https://artofproblemsolving.com/community/c6h546170p3160563}{\textbf{2012 C3}} 	In a $999 \times 999$ square table some cells are white and the remaining ones are red. Let $T$ be the number of triples $(C_1,C_2,C_3)$ of cells, the first two in the same row and the last two in the same column, with $C_1,C_3$ white and $C_2$ red. Find the maximum value $T$ can attain. 
\item \boxed{\href{https://artofproblemsolving.com/community/c6h418796p2363537}{\textbf{2011 C3}}} Let $\mathcal{S}$ be a finite set of at least two points in the plane. Assume that no three points of $\mathcal S$ are collinear. A windmill is a process that starts with a line $\ell$ going through a single point $P \in \mathcal S$. The line rotates clockwise about the pivot $P$ until the first time that the line meets some other point belonging to $\mathcal S$. This point, $Q$, takes over as the new pivot, and the line now rotates clockwise about $Q$, until it next meets a point of $\mathcal S$. This process continues indefinitely.
Show that we can choose a point $P$ in $\mathcal S$ and a line $\ell$ going through $P$ such that the resulting windmill uses each point of $\mathcal S$ as a pivot infinitely many times. 
\item \href{https://artofproblemsolving.com/community/c6h418687p2362298}{\textbf{2010 C3}} 2500 chess kings have to be placed on a $100 \times 100$ chessboard so that
\begin{enumerate}[(a)]
    \item no king can capture any other one (i.e. no two kings are placed in two squares sharing a common vertex);
    \item each row and each column contains exactly 25 kings.
\end{enumerate}
Find the number of such arrangements. 

\textbf{Note:} Two arrangements differing by rotation or symmetry are supposed to be different.
\newpage
\item \href{https://artofproblemsolving.com/community/c6h355915p1934456}{\textbf{2009 C3}} Let $n$ be a positive integer. Given a sequence $\varepsilon_1$, $\dots$, $\varepsilon_{n - 1}$ with $\varepsilon_i = 0$ or $\varepsilon_i = 1$ for each $i = 1$, $\dots$, $n - 1$, the sequences $a_0$, $\dots$, $a_n$ and $b_0$, $\dots$, $b_n$ are constructed by the following rules:\[a_0 = b_0 = 1, \quad a_1 = b_1 = 7,\]\[\begin{array}{lll}
	a_{i+1} = 
	\begin{cases}
		2a_{i-1} + 3a_i, \\
		3a_{i-1} + a_i, 
	\end{cases} & 
        \begin{array}{l} 
                \text{if } \varepsilon_i = 0, \\  
                \text{if } \varepsilon_i = 1, \end{array} 
         & \text{for each } i = 1, \dots, n - 1, \\[15pt]
        b_{i+1}= 
        \begin{cases}
		2b_{i-1} + 3b_i, \\
		3b_{i-1} + b_i, 
	\end{cases} & 
        \begin{array}{l} 
                \text{if } \varepsilon_{n-i} = 0, \\  
                \text{if } \varepsilon_{n-i} = 1, \end{array} 
         & \text{for each } i = 1, \dots, n - 1.
	\end{array}\]Prove that $a_n = b_n$. 
\item \href{https://artofproblemsolving.com/community/c6h287860p1555907}{\textbf{2008 C3}} In the coordinate plane consider the set $ S$ of all points with integer coordinates. For a positive integer $ k$, two distinct points $A$, $ B\in S$ will be called $ k$-friends if there is a point $ C\in S$ such that the area of the triangle $ ABC$ is equal to $ k$. A set $ T\subset S$ will be called $ k$-clique if every two points in $ T$ are $ k$-friends. Find the least positive integer $ k$ for which there exits a $ k$-clique with more than 200 elements. 
\item \href{https://artofproblemsolving.com/community/c6h214709p1187179}{\textbf{2007 C3}} Find all positive integers $ n$ for which the numbers in the set $ S = \{1,2, \dots,n \}$ can be colored red and blue, with the following condition being satisfied: The set $ S \times S \times S$ contains exactly $ 2007$ ordered triples $ \left(x, y, z\right)$ such that:
\begin{enumerate}[(a)]
    \item the numbers $ x$, $ y$, $ z$ are of the same color,
and
    \item the number $ x + y + z$ is divisible by $ n$. 
\end{enumerate}


\item \href{https://artofproblemsolving.com/community/c6h155694p874985}{\textbf{2006 C3}} Let $ S$ be a finite set of points in the plane such that no three of them are on a line. For each convex polygon $ P$ whose vertices are in $ S$, let $ a(P)$ be the number of vertices of $ P$, and let $ b(P)$ be the number of points of $ S$ which are outside $ P$. A line segment, a point, and the empty set are considered as convex polygons of $ 2$, $ 1$, and $ 0$ vertices respectively. Prove that for every real number $ x$\[\sum_{P}{x^{a(P)}(1 - x)^{b(P)}} = 1,\]where the sum is taken over all convex polygons with vertices in $ S$. 
\end{itemize}
\newpage
\section*{C4 Problems}  \label{C4}
\begin{itemize}[leftmargin=*]


\item \href{https://www.imo-official.org/assets/documents/problems/2026/IMO2026SL.pdf}{\textbf{2026 C4}} TO BE REVEALED!

\item \href{https://artofproblemsolving.com/community/c6h3866872p38676989}{\textbf{2025 C4}} Bluey is placing lollies on a $1000\times1000$ grid. Initially, every cell of the grid is empty. At each step, Bluey chooses a row or column in which the total number of lollies is a multiple of $3$, chooses an empty cell in this row or column, and places a lolly in it. If there is no such row or column, Bluey stops.

Determine the maximum number of lollies that Bluey can place.
\item \boxed{\href{https://artofproblemsolving.com/community/c6h3359777p31218774}{\textbf{2024 C4}}} Turbo the snail plays a game on a board with $2024$ rows and $2023$ columns. There are hidden monsters in $2022$ of the cells. Initially, Turbo does not know where any of the monsters are, but he knows that there is exactly one monster in each row except the first row and the last row, and that each column contains at most one monster.

Turbo makes a series of attempts to go from the first row to the last row. On each attempt, he chooses to start on any cell in the first row, then repeatedly moves to an adjacent cell sharing a common side. (He is allowed to return to a previously visited cell.) If he reaches a cell with a monster, his attempt ends and he is transported back to the first row to start a new attempt. The monsters do not move, and Turbo remembers whether or not each cell he has visited contains a monster. If he reaches any cell in the last row, his attempt ends and the game is over.

Determine the minimum value of $n$ for which Turbo has a strategy that guarantees reaching the last row on the $n$-th attempt or earlier, regardless of the locations of the monsters.
\item \href{https://artofproblemsolving.com/community/c6h3359724p31218375}{\textbf{2023 C4}} Let $n\geqslant  2$ be a positive integer. Paul has a $1\times n^2$ rectangular strip consisting of $n^2$ unit squares, where the $i^{\text{th}}$ square is labelled with $i$ for all $1\leqslant i\leqslant n^2$. He wishes to cut the strip into several pieces, where each piece consists of a number of consecutive unit squares, and then translate (without rotating or flipping) the pieces to obtain an $n\times n$ square satisfying the following property: if the unit square in the $i^{\text{th}}$ row and $j^{\text{th}}$ column is labelled with $a_{ij}$, then $a_{ij}-(i+j-1)$ is divisible by $n$.

Determine the smallest number of pieces Paul needs to make in order to accomplish this.
\item \href{https://artofproblemsolving.com/community/c6h3107331p28104263}{\textbf{2022 C4}} Let $n > 3$ be a positive integer. Suppose that $n$ children are arranged in a circle, and $n$ coins are distributed between them (some children may have no coins). At every step, a child with at least 2 coins may give 1 coin to each of their immediate neighbors on the right and left. Determine all initial distributions of the coins from which it is possible that, after a finite number of steps, each child has exactly one coin. 
\item \href{https://artofproblemsolving.com/community/c6h2882594p25627941}{\textbf{2021 C4}} The kingdom of Anisotropy consists of $n$ cities. For every two cities there exists exactly one direct one-way road between them. We say that a path from $X$ to $Y$ is a sequence of roads such that one can move from $X$ to $Y$ along this sequence without returning to an already visited city. A collection of paths is called diverse if no road belongs to two or more paths in the collection.
Let $A$ and $B$ be two distinct cities in Anisotropy. Let $N_{AB}$ denote the maximal number of paths in a diverse collection of paths from $A$ to $B$. Similarly, let $N_{BA}$ denote the maximal number of paths in a diverse collection of paths from $B$ to $A$. Prove that the equality $N_{AB} = N_{BA}$ holds if and only if the number of roads going out from $A$ is the same as the number of roads going out from $B$. 

\newpage

\item \href{https://artofproblemsolving.com/community/c6h2625914p22698447}{\textbf{2020 C4}} 	The Fibonacci numbers $F_0, F_1, F_2, \dots $ are defined inductively by $F_0=0, F_1=1$, and $F_{n+1}=F_n+F_{n-1}$ for $n \geqslant 1$. Given an integer $n \geqslant 2$, determine the smallest size of a set $S$ of integers such that for every $k=2, 3, \dots , n$ there exist some $x, y \in S$ such that $x-y=F_k$. 

\item \href{https://artofproblemsolving.com/community/c6h2279023p17829075}{\textbf{2019 C4}} On a flat plane in Camelot, King Arthur builds a labyrinth $\mathfrak{L}$ consisting of $n$ walls, each of which is an infinite straight line. No two walls are parallel, and no three walls have a common point. Merlin then paints one side of each wall entirely red and the other side entirely blue.

At the intersection of two walls there are four corners: two diagonally opposite corners where a red side and a blue side meet, one corner where two red sides meet, and one corner where two blue sides meet. At each such intersection, there is a two-way door connecting the two diagonally opposite corners at which sides of different colours meet.

After Merlin paints the walls, Morgana then places some knights in the labyrinth. The knights can walk through doors, but cannot walk through walls.

Let $k(\mathfrak{L})$ be the largest number $k$ such that, no matter how Merlin paints the labyrinth $\mathfrak{L},$ Morgana can always place at least $k$ knights such that no two of them can ever meet. For each $n,$ what are all possible values for $k(\mathfrak{L}),$ where $\mathfrak{L}$ is a labyrinth with $n$ walls? 
\item \boxed{\href{https://artofproblemsolving.com/community/c6h1670582p10626557}{\textbf{2018 C4}}} 	An anti-Pascal triangle is an equilateral triangular array of numbers such that, except for the numbers in the bottom row, each number is the absolute value of the difference of the two numbers immediately below it. For example, the following is an anti-Pascal triangle with four rows which contains every integer from $1$ to $10$.
\[\begin{array}{
c@{\hspace{4pt}}c@{\hspace{4pt}}
c@{\hspace{4pt}}c@{\hspace{2pt}}c@{\hspace{2pt}}c@{\hspace{4pt}}c
} \vspace{4pt}
 & & & 4 & & &  \\\vspace{4pt}
 & & 2 & & 6 & &  \\\vspace{4pt}
 & 5 & & 7 & & 1 & \\\vspace{4pt}
 8 & & 3 & & 10 & & 9 \\\vspace{4pt}
\end{array}\]Does there exist an anti-Pascal triangle with $2018$ rows which contains every integer from $1$ to $1 + 2 + 3 + \dots + 2018$? 
\item \boxed{\href{https://artofproblemsolving.com/community/c6h1480685p8639240}{\textbf{2017 C4}}} 	An integer $N \geqslant 2$ is given. A collection of $N(N + 1)$ soccer players, no two of whom are of the same height, stand in a row. Sir Alex wants to remove $N(N - 1)$ players from this row leaving a new row of $2N$ players in which the following $N$ conditions hold:
\begin{enumerate}[(1)]
    \item no one stands between the two tallest players,
    \item no one stands between the third and fourth tallest players,\\
    \vdots
    \item[(N)]  no one stands between the two shortest players.
\end{enumerate}

Show that this is always possible. 

\newpage

\item \boxed{\href{https://artofproblemsolving.com/community/c6h1270471p6637677}{\textbf{2016 C4}}}  Find all integers $n$ for which each cell of $n \times n$ table can be filled with one of the letters $I,M$ and $O$ in such a way that:
\begin{enumerate}[(a)]
    \item  in each row and each column, one third of the entries are $I$, one third are $M$ and one third are $O$; and 
\item  in any diagonal, if the number of entries on the diagonal is a multiple of three, then one third of the entries are $I$, one third are $M$ and one third are $O$.
\end{enumerate}
\textbf{Note.} The rows and columns of an $n \times n$ table are each labelled $1$ to $n$ in a natural order. Thus each cell corresponds to a pair of positive integer $(i,j)$ with $1 \leqslant i,j \leqslant n$. For $n>1$, the table has $4n-2$ diagonals of two types. A diagonal of first type consists all cells $(i,j)$ for which $i+j$ is a constant, and the diagonal of this second type consists all cells $(i,j)$ for which $i-j$ is constant. 
\item \href{https://artofproblemsolving.com/community/c6h1268903p6622756}{\textbf{2015 C4}} Let $n$ be a positive integer. Two players $A$ and $B$ play a game in which they take turns choosing positive integers $k \leqslant n$. The rules of the game are:
\begin{enumerate}[(a)]
    \item A player cannot choose a number that has been chosen by either player on any previous turn.
\item A player cannot choose a number consecutive to any of those the player has already chosen on any previous turn.
\item The game is a draw if all numbers have been chosen; otherwise the player who cannot choose a number anymore loses the game.
\end{enumerate}
The player $A$ takes the first turn. Determine the outcome of the game, assuming that both players use optimal strategies. 
\item \href{https://artofproblemsolving.com/community/c6h1113185p5083550}{\textbf{2014 C4}} Construct a tetromino by attaching two $2 \times 1$ dominoes along their longer sides such that the midpoint of the longer side of one domino is a corner of the other domino. This construction yields two kinds of tetrominoes with opposite orientations. Let us call them $S$- and $Z$-tetrominoes, respectively.
Assume that a lattice polygon $P$ can be tiled with $S$-tetrominoes. Prove that no matter how we tile $P$ using only $S$- and $Z$-tetrominoes, we always use an even number of $Z$-tetrominoes. 
\item \href{https://artofproblemsolving.com/community/c6h597131p3543403}{\textbf{2013 C4}} Let $n$ be a positive integer, and let $A$ be a subset of $\{ 1,\dots ,n\}$. An $A$-partition of $n$ into $k$ parts is a representation of n as a sum $n = a_1 + \dots + a_k$, where the parts $a_1 , \dots , a_k $ belong to $A$ and are not necessarily distinct. The number of different parts in such a partition is the number of (distinct) elements in the set $\{ a_1 , a_2 , \dots , a_k \} $.
We say that an $A$-partition of $n$ into $k$ parts is optimal if there is no $A$-partition of $n$ into $r$ parts with $r<k$. Prove that any optimal $A$-partition of $n$ contains at most $\sqrt[3]{6n}$ different parts. 
\item \href{https://artofproblemsolving.com/community/c6h546171p3160567}{\textbf{2012 C4}} 	Players $A$ and $B$ play a game with $N \geqslant  2012$ coins and $2012$ boxes arranged around a circle. Initially $A$ distributes the coins among the boxes so that there is at least $1$ coin in each box. Then the two of them make moves in the order $B,A,B,A,\dots $ by the following rules:
\begin{enumerate}[(a)]
    \item On every move of his $B$ passes $1$ coin from every box to an adjacent box.
\item  On every move of hers $A$ chooses several coins that were not involved in $B$'s previous move and are in different boxes. She passes every coin to an adjacent box.
\end{enumerate}

Player $A$'s goal is to ensure at least $1$ coin in each box after every move of hers, regardless of how $B$ plays and how many moves are made. Find the least $N$ that enables her to succeed. 
\item \href{https://artofproblemsolving.com/community/c6h488627p2738014}{\textbf{2011 C4}} Determine the greatest positive integer $k$ that satisfies the following property: The set of positive integers can be partitioned into $k$ subsets $A_1, A_2, \dots, A_k$ such that for all integers $n \geqslant  15$ and all $i \in \{1, 2, \dots, k\}$ there exist two distinct elements of $A_i$ whose sum is $n.$ 
\item \boxed{\href{https://artofproblemsolving.com/community/c6h356196p1936917}{\textbf{2010 C4}}} Each of the six boxes $B_1$, $B_2$, $B_3$, $B_4$, $B_5$, $B_6$ initially contains one coin. The following operations are allowed : \\
Type 1: Choose a non-empty box $B_j$, $1\leqslant j \leqslant 5$, remove one coin from $B_j$ and add two coins to $B_{j+1}$; \\
Type 2: Choose a non-empty box $B_k$, $1\leqslant k \leqslant 4$, remove one coin from $B_k$ and swap the contents (maybe empty) of the boxes $B_{k+1}$ and $B_{k+2}$. \\
Determine if there exists a finite sequence of operations of the allowed types, such that the five boxes $B_1$, $B_2$, $B_3$, $B_4$, $B_5$ become empty, while box $B_6$ contains exactly $2010^{2010^{2010}}$ coins. 
\item \href{https://artofproblemsolving.com/community/c6h355785p1932929}{\textbf{2009 C4}} For an integer $m\geqslant  1$, we consider partitions of a $2^m\times 2^m$ chessboard into rectangles consisting of cells of chessboard, in which each of the $2^m$ cells along one diagonal forms a separate rectangle of side length $1$. Determine the smallest possible sum of rectangle perimeters in such a partition. 
\item \boxed{\href{https://artofproblemsolving.com/community/c6h215429p1191679}{\textbf{2008 C4}}} 	Let $ n$ and $ k$ be positive integers with $ k \geqslant  n$ and $ k - n$ an even number. Let $ 2n$ lamps labelled $ 1$, $ 2$, $\dots$, $ 2n$ be given, each of which can be either on or off. Initially all the lamps are off. We consider sequences of steps: at each step one of the lamps is switched (from on to off or from off to on). 
Let $ N$ be the number of such sequences consisting of $ k$ steps and resulting in the state where lamps $ 1$ through $ n$ are all on, and lamps $ n + 1$ through $ 2n$ are all off. 
Let $ M$ be number of such sequences consisting of $ k$ steps, resulting in the state where lamps $ 1$ through $ n$ are all on, and lamps $ n + 1$ through $ 2n$ are all off, but where none of the lamps $ n + 1$ through $ 2n$ is ever switched on.
Determine $ \frac {N}{M}$. 
\item \href{https://artofproblemsolving.com/community/c6h214669p1186971}{\textbf{2007 C4}} Let $ A_0 = (a_1,\dots,a_n)$ be a finite sequence of real numbers. For each $ k\geqslant  0$, from the sequence $ A_k = (x_1,\dots,x_k)$ we construct a new sequence $ A_{k + 1}$ in the following way.
1. We choose a partition $ \{1,\dots,n\} = I\cup J$, where $ I$ and $ J$ are two disjoint sets, such that the expression
\[ \left|\sum_{i\in I}x_i - \sum_{j\in J}x_j\right|
\]
attains the smallest value. (We allow $ I$ or $ J$ to be empty; in this case the corresponding sum is 0.) If there are several such partitions, one is chosen arbitrarily.
2. We set $ A_{k + 1} = (y_1,\dots,y_n)$ where $ y_i = x_i + 1$ if $ i\in I$, and $ y_i = x_i - 1$ if $ i\in J$.

Prove that for some $ k$, the sequence $ A_k$ contains an element $ x$ such that $\displaystyle{|x|\geqslant \frac n2}$. 
\item \href{https://artofproblemsolving.com/community/c6h155696p874991}{\textbf{2006 C4}} A cake has the form of an $ n$ x $ n$ square composed of $ n^{2}$ unit squares. Strawberries lie on some of the unit squares so that each row or column contains exactly one strawberry; call this arrangement $\mathcal{A}$.

Let $\mathcal{B}$ be another such arrangement. Suppose that every grid rectangle with one vertex at the top left corner of the cake contains no fewer strawberries of arrangement $\mathcal{B}$ than of arrangement $\mathcal{A}$. Prove that arrangement $\mathcal{B}$ can be obtained from $ \mathcal{A}$ by performing a number of switches, defined as follows:

A switch consists in selecting a grid rectangle with only two strawberries, situated at its top right corner and bottom left corner, and moving these two strawberries to the other two corners of that rectangle. 
\end{itemize}
\newpage

\section*{C5 Problems} \label{C5}
\begin{itemize}[leftmargin=*]


\item \href{https://www.imo-official.org/assets/documents/problems/2026/IMO2026SL.pdf}{\textbf{2026 C5}} TO BE REVEALED!

\item \href{https://artofproblemsolving.com/community/c6h3866870p38676985}{\textbf{2025 C5}} Let $n$ be a positive integer, and let $(a_1,a_2,\dots,a_n)$ be any $n-$tuplet of distinct integers. An inversion is a pair $(i,j)$ of integers satisfy $1 \leqslant i < j \leqslant n$ and $a_i > a_j$. An inversion is called odd if $j-i$ is odd. Let $A$ be the number of all inversions, and let $B$ be the number of odd inversions.
\begin{enumerate}[leftmargin=*]
\item[(a)] Prove that $A \leqslant 5B$ for all $n$ and all $n-$tuplets $(a_1,a_2,\dots,a_n)$.
\item[(b)] Prove that there exist an $n$ and an $n-$tuplets $(a_1,a_2,\dots,a_n)$ for which $A> 4.99B$.
\end{enumerate}


\item \href{https://artofproblemsolving.com/community/c6h3610469p35340978}{\textbf{2024 C5}} Let $N$ be a positive integer. Geoff and Ceri play a game in which they start by writing the numbers $1, 2, \dots, N$ on a board. They then take turns to make a move, starting with Geoff. Each move consists of choosing a pair of integers $(k, n)$, where $k \geqslant  0$ and $n$ is one of the integers on the board, and then erasing every integer $s$ on the board such that $2^k \mid n - s$. The game continues until the board is empty. The player who erases the last integer on the board loses.

Determine all values of $N$ for which Geoff can ensure that he wins, no matter how Ceri plays.
\item \href{https://artofproblemsolving.com/community/c6h3359765p31218619}{\textbf{2023 C5}} Elisa has $2023$ treasure chests, all of which are unlocked and empty at first. Each day, Elisa adds a new gem to one of the unlocked chests of her choice, and afterwards, a fairy acts according to the following rules:
\begin{itemize}[leftmargin=*]
\item[$\bullet$] if more than one chests are unlocked, it locks one of them, or
\item[$\bullet$] if there is only one unlocked chest, it unlocks all the chests.
\end{itemize}
Given that this process goes on forever, prove that there is a constant $C$ with the following property: Elisa can ensure that the difference between the numbers of gems in any two chests never exceeds $C$, regardless of how the fairy chooses the chests to unlock.
    \item  \href{https://artofproblemsolving.com/community/c6h3107370p28104581}{\textbf{2022 C5}}	Let $m,n \geqslant  2$ be integers, let $X$ be a set with $n$ elements, and let $X_1,X_2,\dots,X_m$ be pairwise distinct non-empty, not necessary disjoint subset of $X$. A function $f \colon X \to \{1,2,\dots,n+1\}$ is called nice if there exists an index $k$ such that\[\sum_{x \in X_k} f(x)>\sum_{x \in X_i} f(x) \quad \text{for all } i \ne k.\]Prove that the number of nice functions is at least $n^n$.
    \item \href{https://artofproblemsolving.com/community/c6h2882530p25627454}{\textbf{2021 C5}} 	Let $n$ and $k$ be two integers with $n>k\geqslant  1$. There are $2n+1$ students standing in a circle. Each student $S$ has $2k$ neighbors - namely, the $k$ students closest to $S$ on the left, and the $k$ students closest to $S$ on the right.

Suppose that $n+1$ of the students are girls, and the other $n$ are boys. Prove that there is a girl with at least $k$ girls among her neighbors.
    \item \href{https://artofproblemsolving.com/community/c6h2625882p22698214}{\textbf{2020 C5}} Let $p$ be an odd prime, and put $N=\frac{1}{4} (p^3 -p) -1.$ The numbers $1,2, \dots, N$ are painted arbitrarily in two colors, red and blue. For any positive integer $n \leqslant N,$ denote $r(n)$ the fraction of integers $\{ 1,2, \dots, n \}$ that are red.
Prove that there exists a positive integer $a \in \{ 1,2, \dots, p-1\}$ such that $r(n) \neq a/p$ for all $n = 1,2, \dots , N.$
\newpage
    \item \boxed{\href{https://artofproblemsolving.com/community/c6h1876067p12744851}{\textbf{2019 C5}}} A social network has $2019$ users, some pairs of whom are friends. Whenever user $A$ is friends with user $B$, user $B$ is also friends with user $A$. Events of the following kind may happen repeatedly, one at a time:
    
Three users $A$, $B$, and $C$ such that $A$ is friends with both $B$ and $C$, but $B$ and $C$ are not friends, change their friendship statuses such that $B$ and $C$ are now friends, but $A$ is no longer friends with $B$, and no longer friends with $C$. All other friendship statuses are unchanged.

Initially, $1010$ users have $1009$ friends each, and $1009$ users have $1010$ friends each. Prove that there exists a sequence of such events after which each user is friends with at most one other user.
    \item  \href{https://artofproblemsolving.com/community/c6h1876763p12752820}{\textbf{2018 C5}} Let $k$ be a positive integer. The organising commitee of a tennis tournament is to schedule the matches for $2k$ players so that every two players play once, each day exactly one match is played, and each player arrives to the tournament site the day of his first match, and departs the day of his last match. For every day a player is present on the tournament, the committee has to pay $1$ coin to the hotel. The organisers want to design the schedule so as to minimise the total cost of all players' stays. Determine this minimum cost.
    \item \boxed{\href{https://artofproblemsolving.com/community/c6h1480157p8633324}{\textbf{2017 C5}}} A hunter and an invisible rabbit play a game in the Euclidean plane. The rabbit's starting point, $A_0,$ and the hunter's starting point, $B_0$ are the same. After $n-1$ rounds of the game, the rabbit is at point $A_{n-1}$ and the hunter is at point $B_{n-1}.$ In the $n^{\text{th}}$ round of the game, three things occur in order:
    \begin{enumerate}[(a)]
        \item The rabbit moves invisibly to a point $A_n$ such that the distance between $A_{n-1}$ and $A_n$ is exactly $1.$
        \item A tracking device reports a point $P_n$ to the hunter. The only guarantee provided by the tracking device to the hunter is that the distance between $P_n$ and $A_n$ is at most $1.$
        \item The hunter moves visibly to a point $B_n$ such that the distance between $B_{n-1}$ and $B_n$ is exactly $1.$
    \end{enumerate}
Is it always possible, no matter how the rabbit moves, and no matter what points are reported by the tracking device, for the hunter to choose her moves so that after $10^9$ rounds, she can ensure that the distance between her and the rabbit is at most $100?$
    \item \href{https://artofproblemsolving.com/community/c6h1480696p8639265}{\textbf{2016 C5}} 	Let $n \geqslant  3$ be a positive integer. Find the maximum number of diagonals in a regular $n$-gon one can select, so that any two of them do not intersect in the interior or they are perpendicular to each other.
    \item \boxed{\href{https://artofproblemsolving.com/community/c6h1113168p5083494}{\textbf{2015 C5}}} The sequence $a_1,a_2,\dots$ of integers satisfies the conditions:
\begin{enumerate}[(a)]
    \item $1\leqslant a_j\leqslant 2015$ for all $j\geqslant 1$,
    \item $k+a_k\neq \ell+a_\ell$ for all $1\leqslant k<\ell$.
\end{enumerate}
Prove that there exist two positive integers $b$ and $N$ for which\[\left\vert\sum_{j=m+1}^n(a_j-b)\right\vert\leqslant 1007^2\]for all integers $m$ and $n$ such that $n>m\geqslant N$.
\newpage
    \item \boxed{\href{https://artofproblemsolving.com/community/c6h597095p3543151}{\textbf{2014 C5}}} A set of lines in the plane is in general position if no two are parallel and no three pass through the same point. A set of lines in general position cuts the plane into regions, some of which have finite area; we call these its finite regions. Prove that for all sufficiently large $n$, in any set of $n$ lines in general position it is possible to colour at least $\sqrt{n}$ lines blue in such a way that none of its finite regions has a completely blue boundary.

\textbf{Note:} Results with $\sqrt{n}$ replaced by $c\sqrt{n}$ will be awarded points depending on the value of the constant $c$.
    \item \href{https://artofproblemsolving.com/community/c6h597133p3543427}{\textbf{2013 C5}} Let $r$ be a positive integer, and let $a_0 , a_1 , \dots $ be an infinite sequence of real numbers. Assume that for all nonnegative integers $m$ and $s$ there exists a positive integer $n \in [m+1, m+r]$ such that
\[ a_m + a_{m+1} +\dots +a_{m+s} = a_n + a_{n+1} +\dots +a_{n+s} \]
Prove that the sequence is periodic, i.e. there exists some $p \geqslant 1 $ such that $a_{n+p} =a_n $ for all $n \geqslant 0$.
    \item \href{https://artofproblemsolving.com/community/c6h546172p3160571}{\textbf{2012 C5}} The columns and the row of a $3n \times 3n$ square board are numbered $1,2,\dots ,3n$. Every square $(x,y)$ with $1 \leqslant x,y \leqslant 3n$ is colored asparagus, byzantium or citrine according as the modulo $3$ remainder of $x+y$ is $0,1$ or $2$ respectively. One token colored asparagus, byzantium or citrine is placed on each square, so that there are $3n^2$ tokens of each color. Suppose that one can permute the tokens so that each token is moved to a distance of at most $d$ from its original position, each asparagus token replaces a byzantium token, each byzantium token replaces a citrine token, and each citrine token replaces an asparagus token. 
    
    Prove that it is possible to permute the tokens so that each token is moved to a distance of at most $d+2$ from its original position, and each square contains a token with the same color as the square.
    \item \href{https://artofproblemsolving.com/community/c6h488628p2738017}{\textbf{2011 C5}} Let $m$ be a positive integer, and consider a $m\times m$ checkerboard consisting of unit squares. At the centre of some of these unit squares there is an ant. At time $0$, each ant starts moving with speed $1$ parallel to some edge of the checkerboard. When two ants moving in the opposite directions meet, they both turn $90^{\circ}$ clockwise and continue moving with speed $1$. When more than $2$ ants meet, or when two ants moving in perpendicular directions meet, the ants continue moving in the same direction as before they met. When an ant reaches one of the edges of the checkerboard, it falls off and will not re-appear.

Considering all possible starting positions, determine the latest possible moment at which the last ant falls off the checkerboard, or prove that such a moment does not necessarily exist.
    \item \href{https://artofproblemsolving.com/community/c6h418688p2362302}{\textbf{2010 C5}} $n \geqslant  4$ players participated in a tennis tournament. Any two players have played exactly one game, and there was no tie game. We call a company of four players $bad$ if one player was defeated by the other three players, and each of these three players won a game and lost another game among themselves. Suppose that there is no bad company in this tournament. Let $w_i$ and $l_i$ be respectively the number of wins and losses of the $i$-th player. Prove that\[\sum^n_{i=1} \left(w_i - l_i\right)^3 \geqslant  0.\]
    \newpage
    \item \href{https://artofproblemsolving.com/community/c6h355786p1932930}{\textbf{2009 C5}} Five identical empty buckets of $2$-liter capacity stand at the vertices of a regular pentagon. Cinderella and her wicked Stepmother go through a sequence of rounds: At the beginning of every round, the Stepmother takes one liter of water from the nearby river and distributes it arbitrarily over the five buckets. Then Cinderella chooses a pair of neighbouring buckets, empties them to the river and puts them back. Then the next round begins. The Stepmother goal's is to make one of these buckets overflow. Cinderella's goal is to prevent this. Can the wicked Stepmother enforce a bucket overflow?
    \item \href{https://artofproblemsolving.com/community/c6h287861p1555908}{\textbf{2008 C5}} Let $ S = \{x_1, x_2, \dots, x_{k + l}\}$ be a $ (k + l)$-element set of real numbers contained in the interval $ [0, 1]$; $ k$ and $ l$ are positive integers. A $ k$-element subset $ A\subset S$ is called nice if
\[ \left |\frac {1}{k}\sum_{x_i\in A} x_i - \frac {1}{l}\sum_{x_j\in S\setminus A} x_j\right |\leqslant \frac {k + l}{2kl}\]
Prove that the number of nice subsets is at least $ \dfrac{2}{k + l}\dbinom{k + l}{k}$.
    \item \href{https://artofproblemsolving.com/community/c6h195043p1071276}{\textbf{2007 C5}} In the Cartesian coordinate plane define the strips $ S_n = \{(x,y)|n\leqslant x < n + 1\}$ for every integer $ n.$ Assume each strip $ S_n$ is colored either red or blue, and let $ a$ and $ b$ be two distinct positive integers. Prove that there exists a rectangle with side length $ a$ and $ b$ such that its vertices have the same color.
    \item \href{https://artofproblemsolving.com/community/c6h155699p875001}{\textbf{2006 C5}} An $ (n, k) -$ tournament is a contest with $ n$ players held in $ k$ rounds such that:
\begin{enumerate}[(a)]
    \item  Each player plays in each round, and every two players meet at most once.
    \item If player $ A$ meets player $ B$ in round $ i$, player $ C$ meets player $ D$ in round $ i$, and player $ A$ meets player $ C$ in round $ j$, then player $ B$ meets player $ D$ in round $ j$.
\end{enumerate}
Determine all pairs $ (n, k)$ for which there exists an $ (n, k) -$ tournament.
\end{itemize}
\newpage

\section*{C6 Problems} \label{C6}
\begin{itemize}[leftmargin=*]


\item \href{https://www.imo-official.org/assets/documents/problems/2026/IMO2026SL.pdf}{\textbf{2026 C6}} TO BE REVEALED!


\item \href{https://artofproblemsolving.com/community/c6h3866845p38676929}{\textbf{2025 C6}} On a $45\times 45$ square grid there is an echidna. From any cell, the echidna can move to any other cell that shares a side. The echidna makes a sequence of $2024$ moves, after which it has visited each cell exactly once. Prove that it is possible to write the numbers from $1$ to $2025$, one in each cell, in such a way that:
    \begin{itemize}[leftmargin=*]
        \item[$\bullet$]  For any pair of adjacent cells in the same row, the cell containing the larger number was visited earlier.
\item[$\bullet$] For any pair of adjacent cells in the same column, the cell containing the larger number was visited later.

    \end{itemize}
\item \href{https://artofproblemsolving.com/community/c6h3610456p35340931}{\textbf{2024 C6}} Let $n$ and $T$ be positive integers. James has $4n$ marbles with weights $1$, $2$, $\dots$, $4n$. He places them on a balance scale, so that both sides have equal weight. Andrew may move a marble from one side of the scale to the other, so that the absolute difference in weights of the two sides remains at most $T$.

Find, in terms of $n$, the minimum positive integer $T$ such that Andrew may make a sequence of moves such that each marble ends up on the opposite side of the scale, regardless of how James initially placed the marbles.
\item \href{https://artofproblemsolving.com/community/c6h3359747p31218478}{\textbf{2023 C6}} Let $N$ be a positive integer, and consider an $N \times N$ grid. A right-down path is a sequence of grid cells such that each cell is either one cell to the right of or one cell below the previous cell in the sequence. A right-up path is a sequence of grid cells such that each cell is either one cell to the right of or one cell above the previous cell in the sequence.

Prove that the cells of the $N \times N$ grid cannot be partitioned into less than $N$ right-down or right-up paths. For example, the following partition of the $5 \times 5$ grid uses $5$ paths.
\begin{center}
\begin{asy}
size(4cm);
draw((5,-1)--(0,-1)--(0,-2)--(5,-2)--(5,-3)--(0,-3)--(0,-4)--(5,-4),gray+linewidth(0.5)+miterjoin);
draw((1,-5)--(1,0)--(2,0)--(2,-5)--(3,-5)--(3,0)--(4,0)--(4,-5),gray+linewidth(0.5)+miterjoin);
draw((0,0)--(5,0)--(5,-5)--(0,-5)--cycle,black+linewidth(2.5)+miterjoin);
draw((0,-1)--(3,-1)--(3,-2)--(1,-2)--(1,-4)--(4,-4)--(4,-3)--(2,-3)--(2,-2),black+linewidth(2.5)+miterjoin);
draw((3,0)--(3,-1),black+linewidth(2.5)+miterjoin);
draw((1,-4)--(1,-5),black+linewidth(2.5)+miterjoin);
draw((4,-3)--(4,-1)--(5,-1),black+linewidth(2.5)+miterjoin);
\end{asy}
\end{center}
    \item \href{https://artofproblemsolving.com/community/c6h3107333p28104274}{\textbf{2022 C6}} Let $n$ be a positive integer. We start with $n$ piles of pebbles, each initially containing a single pebble. One can perform moves of the following form: choose two piles, take an equal number of pebbles from each pile and form a new pile out of these pebbles. Find (in terms of $n$) the smallest number of nonempty piles that one can obtain by performing a finite sequence of moves of this form.

    \newpage
    
    \item \href{https://artofproblemsolving.com/community/c6h2882577p25627683}{\textbf{2021 C6}} 	A hunter and an invisible rabbit play a game on an infinite square grid. First the hunter fixes a colouring of the cells with finitely many colours. The rabbit then secretly chooses a cell to start in. Every minute, the rabbit reports the colour of its current cell to the hunter, and then secretly moves to an adjacent cell that it has not visited before (two cells are adjacent if they share an edge). The hunter wins if after some finite time either:
    \begin{enumerate}[(a)]
        \item the rabbit cannot move; or
        \item the hunter can determine the cell in which the rabbit started.
    \end{enumerate}
Decide whether there exists a winning strategy for the hunter.
    \item \boxed{\href{https://artofproblemsolving.com/community/c6h2278652p17821656}{\textbf{2020 C6}}} There are $4n$ pebbles of weights $1, 2, 3, \dots, 4n.$ Each pebble is coloured in one of $n$ colours and there are four pebbles of each colour. Show that we can arrange the pebbles into two piles so that the following two conditions are both satisfied:
    \begin{enumerate}[(a)]
        \item The total weights of both piles are the same.
        \item Each pile contains two pebbles of each colour.
    \end{enumerate}
    \item \href{https://artofproblemsolving.com/community/c6h2279018p17828946}{\textbf{2019 C6}} Let $n>1$ be an integer. Suppose we are given $2n$ points in the plane such that no three of them are collinear. The points are to be labelled $A_1, A_2, \dots , A_{2n}$ in some order. We then consider the $2n$ angles $\angle A_1A_2A_3, \angle A_2A_3A_4, \dots , \angle A_{2n-2}A_{2n-1}A_{2n}, \angle A_{2n-1}A_{2n}A_1, \angle A_{2n}A_1A_2$. We measure each angle in the way that gives the smallest positive value (i.e. between $0^{\circ}$ and $180^{\circ}$). Prove that there exists an ordering of the given points such that the resulting $2n$ angles can be separated into two groups with the sum of one group of angles equal to the sum of the other group.	
    \item \href{https://artofproblemsolving.com/community/c6h1876766p12752833}{\textbf{2018 C6}} Let $a$ and $b$ be distinct positive integers. The following infinite process takes place on an initially empty board.
    \begin{enumerate}[(a)]
        \item If there is at least a pair of equal numbers on the board, we choose such a pair and increase one of its components by $a$ and the other by $b$.
        \item If no such pair exists, we write two times the number $0$.
    \end{enumerate}
Prove that, no matter how we make the choices in $(a)$, operation $(b)$ will be performed only finitely many times.
    \item \href{https://artofproblemsolving.com/community/c6h1671284p10632333}{\textbf{2017 C6}}  Let $n > 1$ be a given integer. An $n \times n \times n$ cube is composed of $n^3$ unit cubes. Each unit cube is painted with one colour. For each $n \times n \times 1$ box consisting of $n^2$ unit cubes (in any of the three possible orientations), we consider the set of colours present in that box (each colour is listed only once). This way, we get $3n$ sets of colours, split into three groups according to the orientation.

It happens that for every set in any group, the same set appears in both of the other groups. Determine, in terms of $n$, the maximal possible number of colours that are present.
\newpage

    \item \href{https://artofproblemsolving.com/community/c6h1480703p8639274}{\textbf{2016 C6}} There are $n \geqslant  3$ islands in a city. Initially, the ferry company offers some routes between some pairs of islands so that it is impossible to divide the islands into two groups such that no two islands in different groups are connected by a ferry route.

After each year, the ferry company will close a ferry route between some two islands $X$ and $Y$. At the same time, in order to maintain its service, the company will open new routes according to the following rule: for any island which is connected to a ferry route to exactly one of $X$ and $Y$, a new route between this island and the other of $X$ and $Y$ is added.

Suppose at any moment, if we partition all islands into two nonempty groups in any way, then it is known that the ferry company will close a certain route connecting two islands from the two groups after some years. Prove that after some years there will be an island which is connected to all other islands by ferry routes.
    \item \dbox{\href{https://artofproblemsolving.com/community/c6h1268814p6621831}{\textbf{2015 C6}}} 	Let $S$ be a nonempty set of positive integers. We say that a positive integer $n$ is \textit{clean} if it has a unique representation as a sum of an odd number of distinct elements from $S$. 
    
    Prove that there exist infinitely many positive integers that are not clean.
    
    \item \href{https://artofproblemsolving.com/community/c6h1113186p5083553}{\textbf{2014 C6}} We are given an infinite deck of cards, each with a real number on it. For every real number $x$, there is exactly one card in the deck that has $x$ written on it. Now two players draw disjoint sets $A$ and $B$ of $100$ cards each from this deck. We would like to define a rule that declares one of them a winner. This rule should satisfy the following conditions:
    \begin{enumerate}[(a)]
        \item The winner only depends on the relative order of the $200$ cards: if the cards are laid down in increasing order face down and we are told which card belongs to which player, but not what numbers are written on them, we can still decide the winner.
\item If we write the elements of both sets in increasing order as $A =\{ a_1 , a_2 , \dots, a_{100} \}$ and $B= \{ b_1 , b_2 , \dots , b_{100} \}$, and $a_i > b_i$ for all $i$, then $A$ beats $B$.
\item If three players draw three disjoint sets $A, B, C$ from the deck, $A$ beats $B$ and $B$ beats $C$ then $A$ also beats $C$.
    \end{enumerate}

How many ways are there to define such a rule? Here, we consider two rules as different if there exist two sets $A$ and $B$ such that $A$ beats $B$ according to one rule, but $B$ beats $A$ according to the other.
    \item \href{https://artofproblemsolving.com/community/c6h597136p3543443}{\textbf{2013 C6}} In some country several pairs of cities are connected by direct two-way flights. It is possible to go from any city to any other by a sequence of flights. The distance between two cities is defined to be the least possible numbers of flights required to go from one of them to the other. It is known that for any city there are at most $100$ cities at distance exactly three from it. Prove that there is no city such that more than $2550$ other cities have distance exactly four from it.
    \newpage
    
    \item \boxed{\href{https://artofproblemsolving.com/community/c6h488351p2736406}{\textbf{2012 C6}}} The liar's guessing game is a game played between two players $A$ and $B$. The rules of the game depend on two positive integers $k$ and $n$ which are known to both players.

At the start of the game $A$ chooses integers $x$ and $N$ with $1 \leqslant x \leqslant N.$ Player $A$ keeps $x$ secret, and truthfully tells $N$ to player $B$. Player $B$ now tries to obtain information about $x$ by asking player $A$ questions as follows: each question consists of $B$ specifying an arbitrary set $S$ of positive integers (possibly one specified in some previous question), and asking $A$ whether $x$ belongs to $S$. Player $B$ may ask as many questions as he wishes. After each question, player $A$ must immediately answer it with yes or no, but is allowed to lie as many times as she wants; the only restriction is that, among any $k+1$ consecutive answers, at least one answer must be truthful.

After $B$ has asked as many questions as he wants, he must specify a set $X$ of at most $n$ positive integers. If $x$ belongs to $X$, then $B$ wins; otherwise, he loses. Prove that:
\begin{enumerate}[(a)]
    \item If $n \geqslant 2^k,$ then $B$ can guarantee a win.
    \item For all sufficiently large $k$, there exists an integer $n \geqslant 1.99^k$ such that $B$ cannot guarantee a win.
\end{enumerate}
    \item \href{https://artofproblemsolving.com/community/c6h488630p2738024}{\textbf{2011 C6}} 	Let $n$ be a positive integer, and let $W = \dots x_{-1}x_0x_1x_2 \dots$ be an infinite periodic word, consisting of just letters $a$ and/or $b$. Suppose that the minimal period $N$ of $W$ is greater than $2^n$.

A finite nonempty word $U$ is said to appear in $W$ if there exist indices $k \leqslant \ell$ such that $U=x_k x_{k+1} \dots x_{\ell}$. A finite word $U$ is called ubiquitous if the four words $Ua$, $Ub$, $aU$, and $bU$ all appear in $W$. 

Prove that there are at least $n$ ubiquitous finite nonempty words.

    \item \href{https://artofproblemsolving.com/community/c6h418689p2362306}{\textbf{2010 C6}} Given a positive integer $k$ and other two integers $b > w > 1.$ There are two strings of pearls, a string of $b$ black pearls and a string of $w$ white pearls. The length of a string is the number of pearls on it. One cuts these strings in some steps by the following rules. In each step:
\begin{enumerate}[(a)]
    \item The strings are ordered by their lengths in a non-increasing order. If there are some strings of equal lengths, then the white ones precede the black ones. Then $k$ first ones (if they consist of more than one pearl) are chosen; if there are less than $k$ strings longer than 1, then one chooses all of them.
\item Next, one cuts each chosen string into two parts differing in length by at most one. (For instance, if there are strings of $5, 4, 4, 2$ black pearls, strings of $8, 4, 3$ white pearls and $k = 4,$ then the strings of 8 white, 5 black, 4 white and 4 black pearls are cut into the parts $(4,4), (3,2), (2,2)$ and $(2,2)$ respectively.) The process stops immediately after the step when a first isolated white pearl appears.
\end{enumerate}
Prove that at this stage, there will still exist a string of at least two black pearls.
    \item \href{https://artofproblemsolving.com/community/c6h355787p1932931}{\textbf{2009 C6}} On a $999\times 999$ board a limp rook can move in the following way: From any square it can move to any of its adjacent squares, i.e. a square having a common side with it, and every move must be a turn, i.e. the directions of any two consecutive moves must be perpendicular. A non-intersecting route of the limp rook consists of a sequence of pairwise different squares that the limp rook can visit in that order by an admissible sequence of moves. Such a non-intersecting route is called cyclic, if the limp rook can, after reaching the last square of the route, move directly to the first square of the route and start over.
How many squares does the longest possible cyclic, non-intersecting route of a limp rook visit?

\newpage

    \item \href{https://artofproblemsolving.com/community/c6h287863p1555914}{\textbf{2008 C6}} For $ n\geqslant 2$, let $ S_1$, $ S_2$, $ \dots$, $ S_{2^n}$ be $ 2^n$ subsets of $ A = \{1, 2, 3, \dots, 2^{n + 1}\}$ that satisfy the following property: There do not exist indices $ a$ and $ b$ with $ a < b$ and elements $ x$, $ y$, $ z\in A$ with $ x < y < z$ and $ y$, $ z\in S_a$, and $ x$, $ z\in S_b$. Prove that at least one of the sets $ S_1$, $ S_2$, $ \dots$, $ S_{2^n}$ contains no more than $ 4n$ elements.
    \item \boxed{\href{https://artofproblemsolving.com/community/c6h159838p893746}{\textbf{2007 C6}}} In a mathematical competition some competitors are friends. Friendship is always mutual. Call a group of competitors a clique if each two of them are friends. (In particular, any group of fewer than two competitiors is a clique.) The number of members of a clique is called its size.

Given that, in this competition, the largest size of a clique is even, prove that the competitors can be arranged into two rooms such that the largest size of a clique contained in one room is the same as the largest size of a clique contained in the other room.
    \item \href{https://artofproblemsolving.com/community/c6h155700p875004}{\textbf{2006 C6}} A holey triangle is an upward equilateral triangle of side length $n$ with $n$ upward unit triangular holes cut out. A diamond is a $60^\circ-120^\circ$ unit rhombus.
Prove that a holey triangle $T$ can be tiled with diamonds if and only if the following condition holds: Every upward equilateral triangle of side length $k$ in $T$ contains at most $k$ holes, for $1\leqslant k\leqslant n$.
\end{itemize}
\newpage

\section*{C7 Problems} \label{C7}
\begin{itemize}[leftmargin=*]


\item \boxed{\href{https://www.imo-official.org/assets/documents/problems/2026/IMO2026SL.pdf}{\textbf{2026 C7}}} Let $n$ be a positive integer. Liu Bang and Xiang Yu have a stick of length $1$ and want to divide it between themselves. Liu Bang marks at most n points on the stick, and then Xiang Yu marks at most $n$ points on the stick. The marked points are distinct. Then, the stick is cut at all marked points, creating a number of pieces. Afterwards, they take turns claiming any unclaimed piece of the stick, with Liu Bang going first. Each player's goal is to maximise the total length of their own pieces. 

For each $n$, determine the largest value $c$ such that Liu Bang may guarantee a total length of at least $c$, regardless of Xiang Yu's play.

\item \href{https://artofproblemsolving.com/community/c6h3866867p38676971}{\textbf{2025 C7}} Let $P$ be a regular $100-$gon. An \textit{Aussie triangulation} is formed by dividing $P$ into a set $T$ of non-overlapping triangles such that:
    \begin{itemize}[leftmargin=*]
\item[$\bullet$] there are exactly $2025$ different points, including the $100$ vertices of $P$ that are a vertex of at least one triangle in $T$, and
\item[$\bullet$] no three of these $2025$ points are collinear.
\end{itemize}
A $quokkalateral$ is a convex quadrilateral that consists of two triangles in $T$ sharing a common side. Two quokkalaterals may share a triangle in $T$. 

Determine the minimum number of quokkalateral among all Aussie triangulation.
\item \boxed{\href{https://artofproblemsolving.com/community/c6h3358930p31206050}{\textbf{2024 C7}}} Let $a_1, a_2, a_3, \dots$ be an infinite sequence of positive integers, and let $N$ be a positive integer. Suppose that, for each $n > N$, $a_n$ is equal to the number of times $a_{n-1}$ appears in the list $a_1, a_2, \dots, a_{n-1}$.

Prove that at least one of the sequence $a_1, a_3, a_5, \dots$ and $a_2, a_4, a_6, \dots$ is eventually periodic.

(An infinite sequence $b_1, b_2, b_3, \dots$ is eventually periodic if there exist positive integers $p$ and $M$ such that $b_{m+p} = b_m$ for all $m \geqslant M$.)
\item \href{https://artofproblemsolving.com/community/c6h3359751p31218524}{\textbf{2023 C7}} The Imomi archipelago consists of $n\geqslant  2$ islands. Between each pair of distinct islands is a unique ferry line that runs in both directions, and each ferry line is operated by one of $k$ companies. It is known that if any one of the $k$ companies closes all its ferry lines, then it becomes impossible for a traveller, no matter where the traveller starts at, to visit all the islands exactly once (in particular, not returning to the island the traveller started at).

Determine the maximal possible value of $k$ in terms of $n$.
    \item \href{https://artofproblemsolving.com/community/c6h3107374p28104616}{\textbf{2022 C7}} Lucy starts by writing $s$ integer-valued $2022$-tuples on a blackboard. After doing that, she can take any two (not necessarily distinct) tuples $\mathbf{v}=(v_1,\dots,v_{2022})$ and $\mathbf{w}=(w_1,\dots,w_{2022})$ that she has already written, and apply one of the following operations to obtain a new tuple:
\begin{align*}
\mathbf{v}+\mathbf{w}&=(v_1+w_1,\dots,v_{2022}+w_{2022}) \\
\mathbf{v} \lor \mathbf{w}&=(\max(v_1,w_1),\dots,\max(v_{2022},w_{2022}))
\end{align*}and then write this tuple on the blackboard.

It turns out that, in this way, Lucy can write any integer-valued $2022$-tuple on the blackboard after finitely many steps. What is the smallest possible number $s$ of tuples that she initially wrote?
    \item \href{https://artofproblemsolving.com/community/c6h2882582p25627814}{\textbf{2021 C7}} Consider a checkered $3m\times 3m$ square, where $m$ is an integer greater than $1.$ A frog sits on the lower left corner cell $S$ and wants to get to the upper right corner cell $F.$ The frog can hop from any cell to either the next cell to the right or the next cell upwards.

Some cells can be sticky, and the frog gets trapped once it hops on such a cell. A set $X$ of cells is called blocking if the frog cannot reach $F$ from $S$ when all the cells of $X$ are sticky. A blocking set is minimal if it does not contain a smaller blocking set.
\begin{enumerate}[(a)]
    \item Prove that there exists a minimal blocking set containing at least $3m^2-3m$ cells.
    \item Prove that every minimal blocking set containing at most $3m^2$ cells.
\end{enumerate}
    \item \href{https://artofproblemsolving.com/community/c6h2625884p22698238}{\textbf{2020 C7}} Consider any rectangular table having finitely many rows and columns, with a real number $a(r, c)$ in the cell in row $r$ and column $c$. A pair $(R, C)$, where $R$ is a set of rows and $C$ a set of columns, is called a saddle pair if the following two conditions are satisfied:
    \begin{enumerate}[(a)]
        \item  For each row $r^{\prime}$, there is $r \in R$ such that $a(r, c) \geqslant  a\left(r^{\prime}, c\right)$ for all $c \in C$;
\item  For each column $c^{\prime}$, there is $c \in C$ such that $a(r, c) \leqslant a\left(r, c^{\prime}\right)$ for all $r \in R$.
    \end{enumerate}

A saddle pair $(R, C)$ is called a minimal pair if for each saddle pair $\left(R^{\prime}, C^{\prime}\right)$ with $R^{\prime} \subseteq R$ and $C^{\prime} \subseteq C$, we have $R^{\prime}=R$ and $C^{\prime}=C$. Prove that any two minimal pairs contain the same number of rows.
    \item \href{https://artofproblemsolving.com/community/c6h2278979p17828535}{\textbf{2019 C7}} There are 60 empty boxes $B_1,\dots,B_{60}$ in a row on a table and an unlimited supply of pebbles. Given a positive integer $n$, Alice and Bob play the following game.
In the first round, Alice takes $n$ pebbles and distributes them into the 60 boxes as she wishes. Each subsequent round consists of two steps:
\begin{enumerate}[(a)]
    \item Bob chooses an integer $k$ with $1\leqslant k\leqslant 59$ and splits the boxes into the two groups $B_1,\dots,B_k$ and $B_{k+1},\dots,B_{60}$.
    \item Alice picks one of these two groups, adds one pebble to each box in that group, and removes one pebble from each box in the other group.
\end{enumerate}
Bob wins if, at the end of any round, some box contains no pebbles. Find the smallest $n$ such that Alice can prevent Bob from winning.
    \item \href{https://artofproblemsolving.com/community/c6h1876756p12752807}{\textbf{2018 C7}} Consider $2018$ pairwise crossing circles no three of which are concurrent. These circles subdivide the plane into regions bounded by circular $edges$ that meet at $vertices$. Notice that there are an even number of vertices on each circle. Given the circle, alternately colour the vertices on that circle red and blue. In doing so for each circle, every vertex is coloured twice- once for each of the two circle that cross at that point. If the two colours agree at a vertex, then it is assigned that colour; otherwise, it becomes yellow. Show that, if some circle contains at least $2061$ yellow points, then the vertices of some region are all yellow.
    \item \href{https://artofproblemsolving.com/community/c6h1671320p10632559}{\textbf{2017 C7}} For any finite sets $X$ and $Y$ of positive integers, denote by $f_X(k)$ the $k^{\text{th}}$ smallest positive integer not in $X$, and let$$X*Y=X\cup \{ f_X(y):y\in Y\}.$$Let $A$ be a set of $a>0$ positive integers and let $B$ be a set of $b>0$ positive integers. Prove that if $A*B=B*A$, then$$\underbrace{A*(A*\dots (A*(A*A))\dots )}_{\text{ A appears $b$ times}}=\underbrace{B*(B*\dots (B*(B*B))\dots )}_{\text{ B appears $a$ times}}.$$
    \item \boxed{\href{https://artofproblemsolving.com/community/c6h1271000p6642576}{\textbf{2016 C7}}} There are $n\geqslant 2$ line segments in the plane such that every two segments cross and no three segments meet at a point. Geoff has to choose an endpoint of each segment and place a frog on it facing the other endpoint. Then he will clap his hands $n-1$ times. Every time he claps,each frog will immediately jump forward to the next intersection point on its segment. Frogs never change the direction of their jumps. Geoff wishes to place the frogs in such a way that no two of them will ever occupy the same intersection point at the same time.
    \begin{enumerate}[(a)]
        \item  Prove that Geoff can always fulfill his wish if $n$ is odd.
        \item Prove that Geoff can never fulfill his wish if $n$ is even.
    \end{enumerate}
    \item \href{https://artofproblemsolving.com/community/c6h1268906p6622769}{\textbf{2015 C7}} In a company of people some pairs are enemies. A group of people is called unsociable if the number of members in the group is odd and at least $3$, and it is possible to arrange all its members around a round table so that every two neighbors are enemies. Given that there are at most $2015$ unsociable groups, prove that it is possible to partition the company into $11$ parts so that no two enemies are in the same part.
 
    \item \href{https://artofproblemsolving.com/community/c6h1113187p5083554}{\textbf{2014 C7}} Let $M$ be a set of $n \geqslant 4$ points in the plane, no three of which are collinear. Initially these points are connected with $n$ segments so that each point in $M$ is the endpoint of exactly two segments. Then, at each step, one may choose two segments $AB$ and $CD$ sharing a common interior point and replace them by the segments $AC$ and $BD$ if none of them is present at this moment. 
    
    Prove that it is impossible to perform $n^3 /4$ or more such moves.
    \item \boxed{\href{https://artofproblemsolving.com/community/c6h1181543p5720264}{\textbf{2013 C7}}} Let $n \geqslant 3$ be an integer, and consider a circle with $n + 1$ equally spaced points marked on it. Consider all labellings of these points with the numbers $0, 1, \dots , n$ such that each label is used exactly once; two such labellings are considered to be the same if one can be obtained from the other by a rotation of the circle. A labelling is called beautiful if, for any four labels $a < b < c < d$ with $a + d = b + c$, the chord joining the points labelled $a$ and $d$ does not intersect the chord joining the points labelled $b$ and $c$.

Let $M$ be the number of beautiful labelings, and let N be the number of ordered pairs $(x, y)$ of positive integers such that $x + y \leqslant n$ and $\gcd(x, y) = 1$. Prove that$$M = N + 1.$$
    \item \href{https://artofproblemsolving.com/community/c6h546174p3160574}{\textbf{2012 C7}} There are given $2^{500}$ points on a circle labeled $1,2,\dots ,2^{500}$ in some order. Prove that one can choose $100$ pairwise disjoint chords joining some of theses points so that the $100$ sums of the pairs of numbers at the endpoints of the chosen chord are equal.
    \item \href{https://artofproblemsolving.com/community/c6h488631p2738027}{\textbf{2011 C7}} On a square table of $2011$ by $2011$ cells we place a finite number of napkins that each cover a square of $52$ by $52$ cells. In each cell we write the number of napkins covering it, and we record the maximal number $k$ of cells that all contain the same nonzero number. 
    
    Considering all possible napkin configurations, what is the largest value of $k$?

    \newpage
    \item \href{https://artofproblemsolving.com/community/c6h418690p2362308}{\textbf{2010 C7}} Let $P_1, \dots , P_s$ be arithmetic progressions of integers, the following conditions being satisfied:
\begin{enumerate}[(a)]
    \item each integer belongs to at least one of them;
    \item  each progression contains a number which does not belong to other progressions.
\end{enumerate}
Denote by $n$ the least common multiple of the ratios of these progressions; let $n=p_1^{\alpha_1} \dots p_k^{\alpha_k}$ its prime factorization. Prove that\[s \geqslant  1 + \sum^k_{i=1} \alpha_i (p_i - 1).\]
    \item \boxed{\href{https://artofproblemsolving.com/community/c6h289051p1562840}{\textbf{2009 C7}}} 	Let $ a_1, a_2, \dots , a_n$ be distinct positive integers and let $ M$ be a set of $ n - 1$ positive integers not containing $ s = a_1 + a_2 + \dots + a_n.$ A grasshopper is to jump along the real axis, starting at the point $ 0$ and making $ n$ jumps to the right with lengths $ a_1, a_2, \dots , a_n$ in some order. Prove that the order can be chosen in such a way that the grasshopper never lands on any point in $ M.$
    \item \href{https://artofproblemsolving.com/community/c6h214710p1187185}{\textbf{2007 C7}} Let $ \alpha < \frac {3 - \sqrt {5}}{2}$ be a positive real number. Prove that there exist positive integers $ n$ and $ p > \alpha \cdot 2^n$ for which one can select $ 2 \cdot p$ pairwise distinct subsets $ S_1, \dots, S_p, T_1, \dots, T_p$ of the set $ \{1,2, \dots, n\}$ such that $ S_i \cap T_j \neq \emptyset$ for all $ 1 \leqslant i,j \leqslant p$.
 
    \item \href{https://artofproblemsolving.com/community/c6h155703p875009}{\textbf{2006 C7}} Consider a convex polyhedron without parallel edges and without an edge parallel to any face other than the two faces adjacent to it. Call a pair of points of the polyhedron antipodal if there exist two parallel planes passing through these points and such that the polyhedron is contained between these planes. Let $A$ be the number of antipodal pairs of vertices, and let $B$ be the number of antipodal pairs of midpoint edges. Determine the difference $A-B$ in terms of the numbers of vertices, edges, and faces.
\end{itemize}
\newpage

\section*{C8 Problems} \label{C8}
\begin{itemize}[leftmargin=*]

\item \href{https://www.imo-official.org/assets/documents/problems/2026/IMO2026SL.pdf}{\textbf{2026 C8}} TO BE REVEALED!

\item \boxed{\href{https://artofproblemsolving.com/community/c6h3610487p35341197}{\textbf{2025 C8}}} Consider a $2025\times2025$ grid of unit squares. Matilda wishes to place on the grid some rectangular tiles, possibly of different sizes, such that each side of every tile lies on a grid line and every unit square is covered by at most one tile.

Determine the minimum number of tiles Matilda needs to place so that each row and each column of the grid has exactly one unit square that is not covered by any tile. 
\item \href{https://artofproblemsolving.com/community/c6h3610448p35340921}{\textbf{2024 C8}} Let $n$ be a positive integer. Given an $n\times n$ board, the unit cell in the top left corner is initially coloured black, and the other cells are coloured white. We then apply a series of colouring operations to the board. In each operation, we choose a $2\times 2$ square with exactly one cell coloured black and we colour the remaining three cells of that $2\times 2$ square black.

Determine all values of $n$ such that we can colour the whole board black.
    \item \boxed{\href{https://artofproblemsolving.com/community/c6h2883218p25635163}{\textbf{2022 C8}}} Let $n$ be a positive integer. A Nordic square is an $n \times n$ board containing all the integers from $1$ to $n^2$ so that each cell contains exactly one number. Two different cells are considered adjacent if they share a common side. Every cell that is adjacent only to cells containing larger numbers is called a valley. An uphill path is a sequence of one or more cells such that:
    \begin{enumerate}[(a)]
        \item the first cell in the sequence is a valley,
        \item each subsequent cell in the sequence is adjacent to the previous cell, and
        \item the numbers written in the cells in the sequence are in increasing order.
    \end{enumerate}
Find, as a function of $n$, the smallest possible total number of uphill paths in a Nordic square.
    \item \href{https://artofproblemsolving.com/community/c6h2882558p25627579}{\textbf{2021 C8}} Determine the largest integer $N$ for which there exists a table $T$ of integers with $N$ rows and $100$ columns that has the following properties:
    \begin{enumerate}[(a)]
        \item  Every row contains the numbers $1$, $2$, $\dots$, $100$ in some order.
        \item For any two distinct rows $r$ and $s$, there is a column $c$ such that $|T(r,c) - T(s, c)|\geqslant  2$. (Here $T(r,c)$ is the entry in row $r$ and column $c$.)
    \end{enumerate}

    \item \href{https://artofproblemsolving.com/community/c6h2625868p22698110}{\textbf{2020 C8}} Players $A$ and $B$ play a game on a blackboard that initially contains 2020 copies of the number 1 . In every round, player $A$ erases two numbers $x$ and $y$ from the blackboard, and then player $B$ writes one of the numbers $x+y$ and $|x-y|$ on the blackboard. The game terminates as soon as, at the end of some round, one of the following holds:
    \begin{enumerate}[(a)]
        \item one of the numbers on the blackboard is larger than the sum of all other numbers;
\item there are only zeros on the blackboard
    \end{enumerate}
Player $B$ must then give as many cookies to player $A$ as there are numbers on the blackboard. Player $A$ wants to get as many cookies as possible, whereas player $B$ wants to give as few as possible.

Determine the number of cookies that $A$ receives if both players play optimally.
\newpage
    \item \href{https://artofproblemsolving.com/community/c6h2278997p17828730}{\textbf{2019 C8}} Alice has a map of Wonderland, a country consisting of $n \geqslant  2$ towns. For every pair of towns, there is a narrow road going from one town to the other. One day, all the roads are declared to be "one way" only. Alice has no information on the direction of the roads, but the King of Hearts has offered to help her. She is allowed to ask him a number of questions. For each question in turn, Alice chooses a pair of towns and the King of Hearts tells her the direction of the road connecting those two towns.
    
Alice wants to know whether there is at least one town in Wonderland with at most one outgoing road. Prove that she can always find out by asking at most $4n$ questions.
    \item \href{https://artofproblemsolving.com/community/c6h1671285p10632334}{\textbf{2017 C8}} 	Let $n$ be a given positive integer. In the Cartesian plane, each lattice point with nonnegative coordinates initially contains a butterfly, and there are no other butterflies. The neighborhood of a lattice point $c$ consists of all lattice points within the axis-aligned $(2n+1) \times (2n+1)$ square entered at $c$, apart from $c$ itself. We call a butterfly lonely, crowded, or comfortable, depending on whether the number of butterflies in its neighborhood $N$ is respectively less than, greater than, or equal to half of the number of lattice points in $N$. Every minute, all lonely butterflies fly away simultaneously. This process goes on for as long as there are any lonely butterflies. Assuming that the process eventually stops, determine the number of comfortable butterflies at the final state.
    \item \href{https://artofproblemsolving.com/community/c6h1480704p8639279}{\textbf{2016 C8}} Let $n$ be a positive integer. Determine the smallest positive integer $k$ with the following property: it is possible to mark $k$ cells on a $2n \times 2n$ board so that there exists a unique partition of the board into $1 \times 2$ and $2 \times 1$ dominoes, none of which contain two marked cells.
    \item \href{https://artofproblemsolving.com/community/c6h1113188p5083555}{\textbf{2014 C8}} 	A card deck consists of $1024$ cards. On each card, a set of distinct decimal digits is written in such a way that no two of these sets coincide (thus, one of the cards is empty). Two players alternately take cards from the deck, one card per turn. After the deck is empty, each player checks if he can throw out one of his cards so that each of the ten digits occurs on an even number of his remaining cards. If one player can do this but the other one cannot, the one who can is the winner; otherwise a draw is declared.
    
Determine all possible first moves of the first player after which he has a winning strategy.
    \item \href{https://artofproblemsolving.com/community/c6h597140p3543469}{\textbf{2013 C8}} Players $A$ and $B$ play a "paintful" game on the real line. Player $A$ has a pot of paint with four units of black ink. A quantity $p$ of this ink suffices to blacken a (closed) real interval of length $p$. In every round, player $A$ picks some positive integer $m$ and provides $1/2^m $ units of ink from the pot. Player $B$ then picks an integer $k$ and blackens the interval from $k/2^m$ to $(k+1)/2^m$ (some parts of this interval may have been blackened before). The goal of player $A$ is to reach a situation where the pot is empty and the interval $[0,1]$ is not completely blackened.
    
Decide whether there exists a strategy for player $A$ to win in a finite number of moves.
\newpage
    \item \href{https://artofproblemsolving.com/community/c6h355789p1932933}{\textbf{2009 C8}} For any integer $n\geqslant  2$, we compute the integer $h(n)$ by applying the following procedure to its decimal representation. Let $r$ be the rightmost digit of $n$.
\begin{enumerate}[(a)]
    \item If $r=0$, then the decimal representation of $h(n)$ results from the decimal representation of $n$ by removing this rightmost digit $0$.
    \item If $1\leqslant r \leqslant 9$ we split the decimal representation of $n$ into a maximal right part $R$ that solely consists of digits not less than $r$ and into a left part $L$ that either is empty or ends with a digit strictly smaller than $r$. Then the decimal representation of $h(n)$ consists of the decimal representation of $L$, followed by two copies of the decimal representation of $R-1$. For instance, for the number $17,151,345,543$, we will have $L=17,151$, $R=345,543$ and $h(n)=17,151,345,542,345,542$.
\end{enumerate}
Prove that, starting with an arbitrary integer $n\geqslant  2$, iterated application of $h$ produces the integer $1$ after finitely many steps.
    \item \href{https://artofproblemsolving.com/community/c6h207871p1143934}{\textbf{2007 C8}} Given is a convex polygon $ P$ with $ n$ vertices. Triangle whose vertices lie on vertices of $ P$ is called good if all its sides are unit length. Prove that there are at most $ \frac {2n}{3}$ good triangles.
\end{itemize}
\newpage

\section*{C9 Problems} \label{C9}
\begin{itemize}[leftmargin=*]


\item \href{https://www.imo-official.org/assets/documents/problems/2026/IMO2026SL.pdf}{\textbf{2026 C9}} TO BE REVEALED!

 
    \item \href{https://artofproblemsolving.com/community/c6h3107328p28104256}{\textbf{2022 C9}} Let $\mathbb Z_{\geqslant 0}$ be the set of non-negative integers, and let $f:\mathbb Z_{\geqslant 0}\times \mathbb Z_{\geqslant 0} \to \mathbb Z_{\geqslant 0}$ be a bijection such that whenever $f(x_1,y_1) > f(x_2, y_2)$, we have $f(x_1+1, y_1) > f(x_2 + 1, y_2)$ and $f(x_1, y_1+1) > f(x_2, y_2+1)$.

Let $N$ be the number of pairs of integers $(x,y)$ with $0\leqslant x,y<100$, such that $f(x,y)$ is odd. Find the smallest and largest possible values of $N$.
    \item \href{https://artofproblemsolving.com/community/c6h2279019p17828982}{\textbf{2019 C9}}  For any two different real numbers $x$ and $y$, we define $D(x,y)$ to be the unique integer $d$ satisfying $2^d\leqslant |x-y| < 2^{d+1}$. Given a set of reals $\mathcal F$, and an element $x\in \mathcal F$, we say that the scales of $x$ in $\mathcal F$ are the values of $D(x,y)$ for $y\in\mathcal F$ with $x\neq y$. Let $k$ be a given positive integer.

Suppose that each member $x$ of $\mathcal F$ has at most $k$ different scales in $\mathcal F$ (note that these scales may depend on $x$). What is the maximum possible size of $\mathcal F$?
    \item \href{https://artofproblemsolving.com/community/c6h1113189p5083556}{\textbf{2014 C9}} There are $n$ circles drawn on a piece of paper in such a way that any two circles intersect in two points, and no three circles pass through the same point. Turbo the snail slides along the circles in the following fashion. Initially he moves on one of the circles in clockwise direction. Turbo always keeps sliding along the current circle until he reaches an intersection with another circle. Then he continues his journey on this new circle and also changes the direction of moving, i.e. from clockwise to anticlockwise or $\textit{vice versa}$.
Suppose that Turbo's path entirely covers all circles. Prove that $n$ must be odd.
\end{itemize}
\newpage






\title{Geometry IMO Shortlist}
 
\section*{G1 Problems}  \label{G1}
\begin{itemize}[leftmargin=*]


\item \href{https://www.imo-official.org/assets/documents/problems/2026/IMO2026SL.pdf}{\textbf{2026 G1}} TO BE REVEALED!

\item \href{https://artofproblemsolving.com/community/c6h3866857p38676945}{\textbf{2025 G1}} Let $ABC$ be a triangle such that $\angle A$ is obtuse. Let $D$ and $E$ be points on sides $AB$ and $AC$, respectively, such that $BDEC$ is cyclic. Suppose the circumcircle of triangle $ADE$ intersects side $BC$ at two points $X$ and $Y$, where $B$, $X$, $Y$ and $C$ lie in that order. Let the circumcircles of triangles $BDX$ and $CEY$ be $\Omega$ and $\Gamma$, respectively. Let $P$ be a point such that $PD$ is tangent to $\Omega$ and $PE$ is tangent to $\Gamma$. Let $Q$ be a point such that $QB$ is tangent to $\Omega$ and $QC$ is tangent to $\Gamma$. 
    
    Prove that points $A$, $P$ and $Q$ are collinear.
\item \href{https://artofproblemsolving.com/community/c6h3610481p35341119}{\textbf{2024 G1}} Let $ABCD$ be a cyclic quadrilateral such that $AC<BD<AD$ and $\angle DBA<90^\circ$. Point $E$ lies on the line through $D$ parallel to $AB$ such that $E$ and $C$ lie on opposite sides of line $AD$, and $AC=DE$. Point $F$ lies on the line through $A$ parallel to $CD$ such that $F$ and $C$ lie on opposite sides of line $AD$, and $BD=AF$.

Prove that the perpendicular bisectors of segments $BC$ and $EF$ intersect on the circumcircle of $ABCD$.
\item \href{https://artofproblemsolving.com/community/c6h3359760p31218557}{\textbf{2023 G1}} Let $ABCDE$ be a convex pentagon such that $\angle ABC = \angle AED = 90^\circ$. Suppose that the midpoint of $CD$ is the circumcenter of triangle $ABE$. Let $O$ be the circumcenter of triangle $ACD$.

Prove that line $AO$ passes through the midpoint of segment $BE$.
    \item \boxed{\href{https://artofproblemsolving.com/community/c6h2883216p25635154}{\textbf{2022 G1}}}  Let $ABCDE$ be a convex pentagon such that $BC=DE$. Assume that there is a point $T$ inside $ABCDE$ with $TB=TD,TC=TE$ and $\angle ABT = \angle TEA$. Let line $AB$ intersect lines $CD$ and $CT$ at points $P$ and $Q$, respectively. Assume that the points $P,B,A,Q$ occur on their line in that order. Let line $AE$ intersect $CD$ and $DT$ at points $R$ and $S$, respectively. Assume that the points $R,E,A,S$ occur on their line in that order. 
    
    Prove that the points $P,S,Q,R$ lie on a circle.
\item \href{https://artofproblemsolving.com/community/c6h2882542p25627509}{\textbf{2021 G1}} Let $ABCD$ be a parallelogram with $AC=BC.$ A point $P$ is chosen on the extension of ray $AB$ past $B.$ The circumcircle of $ACD$ meets the segment $PD$ again at $Q.$ The circumcircle of triangle $APQ$ meets the segment $PC$ at $R.$ 

Prove that lines $CD,AQ,BR$ are concurrent.

\item \href{https://artofproblemsolving.com/community/c6h2625929p22698577}{\textbf{2020 G1}} Let $ABC$ be an isosceles triangle with $BC=CA$, and let $D$ be a point inside side $AB$ such that $AD< DB$. Let $P$ and $Q$ be two points inside sides $BC$ and $CA$, respectively, such that $\angle DPB = \angle DQA = 90^{\circ}$. Let the perpendicular bisector of $PQ$ meet line segment $CQ$ at $E$, and let the circumcircles of triangles $ABC$ and $CPQ$ meet again at point $F$, different from $C$.
Suppose that $P$, $E$, $F$ are collinear. 

Prove that $\angle ACB = 90^{\circ}$.
\item \href{https://artofproblemsolving.com/community/c6h2278984p17828603}{\textbf{2019 G1}} Let $ABC$ be a triangle. Circle $\Gamma$ passes through $A$, meets segments $AB$ and $AC$ again at points $D$ and $E$ respectively, and intersects segment $BC$ at $F$ and $G$ such that $F$ lies between $B$ and $G$. The tangent to circle $BDF$ at $F$ and the tangent to circle $CEG$ at $G$ meet at point $T$. Suppose that points $A$ and $T$ are distinct. 

Prove that line $AT$ is parallel to $BC$.
\item \boxed{\href{https://artofproblemsolving.com/community/c6h1670580p10626500}{\textbf{2018 G1}}}  Let $\Gamma$ be the circumcircle of acute triangle $ABC$. Points $D$ and $E$ are on segments $AB$ and $AC$ respectively such that $AD = AE$. The perpendicular bisectors of $BD$ and $CE$ intersect minor arcs $AB$ and $AC$ of $\Gamma$ at points $F$ and $G$ respectively.

Prove that lines $DE$ and $FG$ are either parallel or they are the same line.
\item \href{https://artofproblemsolving.com/community/c6h1671264p10632265}{\textbf{2017 G1}} Let $ABCDE$ be a convex pentagon such that $AB=BC=CD$, $\angle{EAB}=\angle{BCD}$, and $\angle{EDC}=\angle{CBA}$. Prove that the perpendicular line from $E$ to $BC$ and the line segments $AC$ and $BD$ are concurrent.
\item \boxed{\href{https://artofproblemsolving.com/community/c6h1270466p6637656}{\textbf{2016 G1}}}  Triangle $BCF$ has a right angle at $B$. Let $A$ be the point on line $CF$ such that $FA=FB$ and $F$ lies between $A$ and $C$. Point $D$ is chosen so that $DA=DC$ and $AC$ is the bisector of $\angle{DAB}$. Point $E$ is chosen so that $EA=ED$ and $AD$ is the bisector of $\angle{EAC}$. Let $M$ be the midpoint of $CF$. Let $X$ be the point such that $AMXE$ is a parallelogram. 

Prove that $BD,FX$ and $ME$ are concurrent.
\item \dbox{\href{https://artofproblemsolving.com/community/c6h1268782p6621351}{\textbf{2015 G1}}} 	Let $ABC$ be an acute triangle with orthocenter $H$. Let $G$ be the point such that the quadrilateral $ABGH$ is a parallelogram. Let $I$ be the point on the line $GH$ such that $AC$ bisects $HI$. Suppose that the line $AC$ intersects the circumcircle of the triangle $GCI$ at $C$ and $J$. 

Prove that $IJ = AH$.
\item \boxed{\href{https://artofproblemsolving.com/community/c6h597090p3543136}{\textbf{2014 G1}}}  Let $P$ and $Q$ be on segment $BC$ of an acute triangle $ABC$ such that $\angle PAB=\angle BCA$ and $\angle CAQ=\angle ABC$. Let $M$ and $N$ be the points on $AP$ and $AQ$, respectively, such that $P$ is the midpoint of $AM$ and $Q$ is the midpoint of $AN$. 

Prove that the intersection of $BM$ and $CN$ is on the circumference of triangle $ABC$.
\item \boxed{\href{https://artofproblemsolving.com/community/c6h1181533p5720174}{\textbf{2013 G1}}}  	Let $ABC$ be an acute triangle with orthocenter $H$, and let $W$ be a point on the side $BC$, lying strictly between $B$ and $C$. The points $M$ and $N$ are the feet of the altitudes from $B$ and $C$, respectively. Denote by $\omega_1$ is the circumcircle of $BWN$, and let $X$ be the point on $\omega_1$ such that $WX$ is a diameter of $\omega_1$. Analogously, denote by $\omega_2$ the circumcircle of triangle $CWM$, and let $Y$ be the point such that $WY$ is a diameter of $\omega_2$. 

Prove that $X,Y$ and $H$ are collinear.
\item \boxed{\href{https://artofproblemsolving.com/community/c6h488347p2736397}{\textbf{2012 G1}}}  	Given triangle $ABC$ the point $J$ is the centre of the excircle opposite the vertex $A$. This excircle is tangent to the side $BC$ at $M$, and to the lines $AB$ and $AC$ at $K$ and $L$, respectively. The lines $LM$ and $BJ$ meet at $F$, and the lines $KM$ and $CJ$ meet at $G$. Let $S$ be the point of intersection of the lines $AF$ and $BC$, and let $T$ be the point of intersection of the lines $AG$ and $BC$. 

Prove that $M$ is the midpoint of $ST.$ 
\item \href{https://artofproblemsolving.com/community/c6h488823p2739318}{\textbf{2011 G1}} Let $ABC$ be an acute triangle. Let $\omega$ be a circle whose centre $L$ lies on the side $BC$. Suppose that $\omega$ is tangent to $AB$ at $B'$ and $AC$ at $C'$. Suppose also that the circumcentre $O$ of triangle $ABC$ lies on the shorter arc $B'C'$ of $\omega$. 

Prove that the circumcircle of $ABC$ and $\omega$ meet at two points.

\newpage

\item \href{https://artofproblemsolving.com/community/c6h418633p2361970}{\textbf{2010 G1}} Let $ABC$ be an acute triangle with $D, E, F$ the feet of the altitudes lying on $BC, CA, AB$ respectively. One of the intersection points of the line $EF$ and the circumcircle is $P.$ The lines $BP$ and $DF$ meet at point $Q.$ 

Prove that $AP = AQ.$
\item \boxed{\href{https://artofproblemsolving.com/community/c6h289054p1562847}{\textbf{2009 G1}}}  Let $ ABC$ be a triangle with $ AB = AC$ . The angle bisectors of $ \angle C AB$ and $ \angle AB C$ meet the sides $ B C$ and $ C A$ at $ D$ and $ E$ , respectively. Let $ K$ be the incentre of triangle $ ADC$. Suppose that $ \angle B E K = 45^\circ$. 

Find all possible values of $ \angle CAB$.

\item \boxed{\href{https://artofproblemsolving.com/community/c6h215224p1190553}{\textbf{2008 G1}}}  Let $ H$ be the orthocenter of an acute-angled triangle $ ABC$. The circle $ \Gamma_{A}$ centered at the midpoint of $ BC$ and passing through $ H$ intersects the sideline $ BC$ at points $ A_{1}$ and $ A_{2}$. Similarly, define the points $ B_{1}$, $ B_{2}$, $ C_{1}$ and $ C_{2}$. 

Prove that the six points $ A_{1}$, $ A_{2}$, $ B_{1}$, $ B_{2}$, $ C_{1}$ and $ C_{2}$ are concyclic.
\item \boxed{\href{https://artofproblemsolving.com/community/c6h159986p894655}{\textbf{2007 G1}}}  In triangle $ ABC$ the bisector of angle $ BCA$ intersects the circumcircle again at $ R$, the perpendicular bisector of $ BC$ at $ P$, and the perpendicular bisector of $ AC$ at $ Q$. The midpoint of $ BC$ is $ K$ and the midpoint of $ AC$ is $ L$. 

Prove that the triangles $ RPK$ and $ RQL$ have the same area.
\item \boxed{\href{https://artofproblemsolving.com/community/c6h101303p571966}{\textbf{2006 G1}}}  Let $ABC$ be triangle with incenter $I$. A point $P$ in the interior of the triangle satisfies\[\angle PBA+\angle PCA = \angle PBC+\angle PCB.\]Show that $AP \geqslant  AI$, and that equality holds if and only if $P=I$.
\end{itemize}
\newpage
\section*{G2 Problems}  \label{G2}
\begin{itemize}[leftmargin=*]


\item \href{https://www.imo-official.org/assets/documents/problems/2026/IMO2026SL.pdf}{\textbf{2026 G2}} TO BE REVEALED!

\item \href{https://artofproblemsolving.com/community/c6h3866847p38676931}{\textbf{2025 G2}} Let $ABC$ be an acute-angled triangle. Let $I_B$ and $I_C$ be the excenters of triangle $ABC$ opposite $B$ and $C$, respectively. Let $E$ be a point on line $BA$ such that $\angle EI_BB=90^\circ$. Let $F$ be a point on line $CA$ such that $\angle CI_CF=90^\circ$. Let the circle with center $I_B$ and radius $I_BE$ intersect the segment $I_BB$ at $X$. Let the circle with center $I_C$ and radius $I_CF$ intersect the segment $I_CC$ at $Y$. 
    
Prove that $BY$ is perpendicular to $CX$.



\item \boxed{\href{https://artofproblemsolving.com/community/c6h3359767p31218657}{\textbf{2024 G2}}} Let $ABC$ be a triangle with $AB < AC < BC$. Let the incenter and incircle of triangle $ABC$ be $I$ and $\omega$, respectively. Let $X$ be the point on line $BC$ different from $C$ such that the line through $X$ parallel to $AC$ is tangent to $\omega$. Similarly, let $Y$ be the point on line $BC$ different from $B$ such that the line through $Y$ parallel to $AB$ is tangent to $\omega$. Let $AI$ intersect the circumcircle of triangle $ABC$ at $P \ne A$. Let $K$ and $L$ be the midpoints of $AC$ and $AB$, respectively.

Prove that $\angle KIL + \angle YPX = 180^{\circ}$.
\item \href{https://artofproblemsolving.com/community/c6h3359729p31218382}{\textbf{2023 G2}} Let $ABC$ be a triangle with $AC > BC,$ let $\omega$ be the circumcircle of $\triangle ABC,$ and let $r$ be its radius. Point $P$ is chosen on $\overline{AC}$ such taht $BC=CP,$ and point $S$ is the foot of the perpendicular from $P$ to $\overline{AB}$. Ray $BP$ mets $\omega$ again at $D$. Point $Q$ is chosen on line $SP$ such that $PQ = r$ and $S,P,Q$ lie on a line in that order. Finally, let $E$ be a point satisfying $\overline{AE} \perp \overline{CQ}$ and $\overline{BE} \perp \overline{DQ}$. 

Prove that $E$ lies on $\omega$.
\item  \href{https://artofproblemsolving.com/community/c6h3107329p28104258}{\textbf{2022 G2}} In the acute-angled triangle $ABC$, the point $F$ is the foot of the altitude from $A$, and $P$ is a point on the segment $AF$. The lines through $P$ parallel to $AC$ and $AB$ meet $BC$ at $D$ and $E$, respectively. Points $X \ne A$ and $Y \ne A$ lie on the circles $ABD$ and $ACE$, respectively, such that $DA = DX$ and $EA = EY$. 

Prove that $B, C, X,$ and $Y$ are concyclic.
\item \boxed{\href{https://artofproblemsolving.com/community/c6h2625855p22698001}{\textbf{2021 G2}}}  Let $\Gamma$ be a circle with centre $I$, and $A B C D$ a convex quadrilateral such that each of the segments $A B, B C, C D$ and $D A$ is tangent to $\Gamma$. Let $\Omega$ be the circumcircle of the triangle $A I C$. The extension of $B A$ beyond $A$ meets $\Omega$ at $X$, and the extension of $B C$ beyond $C$ meets $\Omega$ at $Z$. The extensions of $A D$ and $C D$ beyond $D$ meet $\Omega$ at $Y$ and $T$, respectively. 

Prove that\[A D+D T+T X+X A=C D+D Y+Y Z+Z C.\]
\item \boxed{\href{https://artofproblemsolving.com/community/c6h2278651p17821635}{\textbf{2020 G2}}}  Consider the convex quadrilateral $ABCD$. The point $P$ is in the interior of $ABCD$. The following ratio equalities hold: \[\angle PAD:\angle PBA:\angle DPA=1:2:3=\angle CBP:\angle BAP:\angle BPC\]Prove that the following three lines meet in a point: the internal bisectors of angles $\angle ADP$ and $\angle PCB$ and the perpendicular bisector of segment $AB$.

\newpage

\item \href{https://artofproblemsolving.com/community/c6h2278992p17828685}{\textbf{2019 G2}} Let $ABC$ be an acute-angled triangle and let $D, E$, and $F$ be the feet of altitudes from $A, B$, and $C$ to sides $BC, CA$, and $AB$, respectively. Denote by $\omega_B$ and $\omega_C$ the incircles of triangles $BDF$ and $CDE$, and let these circles be tangent to segments $DF$ and $DE$ at $M$ and $N$, respectively. Let line $MN$ meet circles $\omega_B$ and $\omega_C$ again at $P \ne M$ and $Q \ne N$, respectively. 

Prove that $MP = NQ$.




\item \href{https://artofproblemsolving.com/community/c6h1876755p12752806}{\textbf{2018 G2}} Let $ABC$ be a triangle with $AB=AC$, and let $M$ be the midpoint of $BC$. Let $P$ be a point such that $PB<PC$ and $PA$ is parallel to $BC$. Let $X$ and $Y$ be points on the lines $PB$ and $PC$, respectively, so that $B$ lies on the segment $PX$, $C$ lies on the segment $PY$, and $\angle PXM=\angle PYM$. 

Prove that the quadrilateral $APXY$ is cyclic.
\item \boxed{\href{https://artofproblemsolving.com/community/c6h1480682p8639236}{\textbf{2017 G2}}}  Let $R$ and $S$ be different points on a circle $\Omega$ such that $RS$ is not a diameter. Let $\ell$ be the tangent line to $\Omega$ at $R$. Point $T$ is such that $S$ is the midpoint of the line segment $RT$. Point $J$ is chosen on the shorter arc $RS$ of $\Omega$ so that the circumcircle $\Gamma$ of triangle $JST$ intersects $\ell$ at two distinct points. Let $A$ be the common point of $\Gamma$ and $\ell$ that is closer to $R$. Line $AJ$ meets $\Omega$ again at $K$. 

Prove that the line $KT$ is tangent to $\Gamma$.

\item \href{https://artofproblemsolving.com/community/c6h1480699p8639270}{\textbf{2016 G2}} Let $ABC$ be a triangle with circumcircle $\Gamma$ and incenter $I$ and let $M$ be the midpoint of $\overline{BC}$. The points $D$, $E$, $F$ are selected on sides $\overline{BC}$, $\overline{CA}$, $\overline{AB}$ such that $\overline{ID} \perp \overline{BC}$, $\overline{IE}\perp \overline{AI}$, and $\overline{IF}\perp \overline{AI}$. Suppose that the circumcircle of $\triangle AEF$ intersects $\Gamma$ at a point $X$ other than $A$. 

Prove that lines $XD$ and $AM$ meet on $\Gamma$.
\item \boxed{\href{https://artofproblemsolving.com/community/c6h1113163p5083464}{\textbf{2015 G2}}}  Triangle $ABC$ has circumcircle $\Omega$ and circumcenter $O$. A circle $\Gamma$ with center $A$ intersects the segment $BC$ at points $D$ and $E$, such that $B$, $D$, $E$, and $C$ are all different and lie on line $BC$ in this order. Let $F$ and $G$ be the points of intersection of $\Gamma$ and $\Omega$, such that $A$, $F$, $B$, $C$, and $G$ lie on $\Omega$ in this order. Let $K$ be the second point of intersection of the circumcircle of triangle $BDF$ and the segment $AB$. Let $L$ be the second point of intersection of the circumcircle of triangle $CGE$ and the segment $CA$. Suppose that the lines $FK$ and $GL$ are different and intersect at the point $X$. 

Prove that $X$ lies on the line $AO$.
\item \href{https://artofproblemsolving.com/community/c6h1113190p5083557}{\textbf{2014 G2}} 	Let $ABC$ be a triangle. The points $K, L,$ and $M$ lie on the segments $BC, CA,$ and $AB,$ respectively, such that the lines $AK, BL,$ and $CM$ intersect in a common point. 

Prove that it is possible to choose two of the triangles $ALM, BMK,$ and $CKL$ whose inradii sum up to at least the inradius of the triangle $ABC$.
\item \href{https://artofproblemsolving.com/community/c6h597143p3543487}{\textbf{2013 G2}} Let $\omega$ be the circumcircle of a triangle $ABC$. Denote by $M$ and $N$ the midpoints of the sides $AB$ and $AC$, respectively, and denote by $T$ the midpoint of the arc $BC$ of $\omega$ not containing $A$. The circumcircles of the triangles $AMT$ and $ANT$ intersect the perpendicular bisectors of $AC$ and $AB$ at points $X$ and $Y$, respectively; assume that $X$ and $Y$ lie inside the triangle $ABC$. The lines $MN$ and $XY$ intersect at $K$. 

Prove that $KA=KT$.
\newpage

\item \href{https://artofproblemsolving.com/community/c6h546176p3160578}{\textbf{2012 G2}} 	Let $ABCD$ be a cyclic quadrilateral whose diagonals $AC$ and $BD$ meet at $E$. The extensions of the sides $AD$ and $BC$ beyond $A$ and $B$ meet at $F$. Let $G$ be the point such that $ECGD$ is a parallelogram, and let $H$ be the image of $E$ under reflection in $AD$. 

Prove that $D,H,F,G$ are concyclic.

\item \href{https://artofproblemsolving.com/community/c6h488825p2739321}{\textbf{2011 G2}} 	Let $A_1A_2A_3A_4$ be a non-cyclic quadrilateral. Let $O_1$ and $r_1$ be the circumcentre and the circumradius of the triangle $A_2A_3A_4$. Define $O_2,O_3,O_4$ and $r_2,r_3,r_4$ in a similar way. Prove that \[\frac{1}{O_1A_1^2-r_1^2}+\frac{1}{O_2A_2^2-r_2^2}+\frac{1}{O_3A_3^2-r_3^2}+\frac{1}{O_4A_4^2-r_4^2}=0.\]
\item \boxed{\href{https://artofproblemsolving.com/community/c6h356195p1936916}{\textbf{2010 G2}}}  Let $P$ be a point interior to triangle $ABC$ (with $CA \neq CB$). The lines $AP$, $BP$ and $CP$ meet again its circumcircle $\Gamma$ at $K$, $L$, respectively $M$. The tangent line at $C$ to $\Gamma$ meets the line $AB$ at $S$. Show that from $SC = SP$ follows $MK = ML$.
\item \boxed{\href{https://artofproblemsolving.com/community/c6h288839p1561572}{\textbf{2009 G2}}}  	Let $ ABC$ be a triangle with circumcentre $ O$. The points $ P$ and $ Q$ are interior points of the sides $ CA$ and $ AB$ respectively. Let $ K,L$ and $ M$ be the midpoints of the segments $ BP,CQ$ and $ PQ$. respectively, and let $ \Gamma$ be the circle passing through $ K,L$ and $ M$. Suppose that the line $ PQ$ is tangent to the circle $ \Gamma$. 

Prove that $ OP = OQ$.
\item \href{https://artofproblemsolving.com/community/c6h287864p1555917}{\textbf{2008 G2}} Given trapezoid $ ABCD$ with parallel sides $ AB$ and $ CD$, assume that there exist points $ E$ on line $ BC$ outside segment $ BC$, and $ F$ inside segment $ AD$ such that $ \angle DAE = \angle CBF$. Denote by $ I$ the point of intersection of $ CD$ and $ EF$, and by $ J$ the point of intersection of $ AB$ and $ EF$. Let $ K$ be the midpoint of segment $ EF$, assume it does not lie on line $ AB$. 

Prove that $ I$ belongs to the circumcircle of $ ABK$ if and only if $ K$ belongs to the circumcircle of $ CDJ$.
\item \href{https://artofproblemsolving.com/community/c6h207870p1143932}{\textbf{2007 G2}} Denote by $ M$ midpoint of side $ BC$ in an isosceles triangle $ \triangle ABC$ with $ AC = AB$. Take a point $ X$ on a smaller arc $MA$ of circumcircle of triangle $ \triangle ABM$. Denote by $ T$ point inside of angle $ BMA$ such that $ \angle TMX = 90$ and $ TX = BX$. 

Prove that $ \angle MTB - \angle CTM$ does not depend on choice of $ X$.
\item \href{https://artofproblemsolving.com/community/c6h155705p875014}{\textbf{2006 G2}} Let $ ABCD$ be a trapezoid with parallel sides $ AB > CD$. Points $ K$ and $ L$ lie on the line segments $ AB$ and $ CD$, respectively, so that $AK/KB=DL/LC$. Suppose that there are points $ P$ and $ Q$ on the line segment $ KL$ satisfying\[\angle{APB} = \angle{BCD}\qquad\text{and}\qquad \angle{CQD} = \angle{ABC}.\]Prove that the points $ P$, $ Q$, $ B$ and $ C$ are concyclic.
\end{itemize}

\newpage
\section*{G3 Problems}  \label{G3}
\begin{itemize}[leftmargin=*]


\item \href{https://www.imo-official.org/assets/documents/problems/2026/IMO2026SL.pdf}{\textbf{2026 G3}} TO BE REVEALED!

\item \href{https://artofproblemsolving.com/community/c6h3866884p38677024}{\textbf{2025 G3}} Let $ABC$ be an acute-triangle with $AB<AC$. Let $\Gamma$ be the circumcircle of $ABC$. The perpendicular of $B$ to $AC$ intersects $\Gamma$ again at $D\ne B$. Let $Q$ be the point on $\Gamma$ such that $Q\ne B$ and $AB=AQ$. Lines $BQ$ and $AC$ intersect at $R$. Line $DR$ intersects $\Gamma$ again at $S\ne D$. Segment $BC$ intersects the circumcircle of triangle $QRS$ at $T$. Lines $TQ$ and $RS$ intersect at $X$.

    Prove that $CX$ is perpendicular to $AB$.


\item \href{https://artofproblemsolving.com/community/c6h3610478p35341061}{\textbf{2024 G3}} Let $ABCDE$ be a convex pentagon and let $M$ be the midpoint of $AB$. Suppose that segment $AB$ is tangent to the circumcircle of triangle $CME$ at $M$ and that $D$ lies on the circumircles of $AME$ and $BMC$. Lines $AD$ and $ME$ interesect at $K$, and lines $BD$ and $MC$ intersect at $L$. Points $P$ and $Q$ lie on line $EC$ so that $\angle PDC = \angle EDQ = \angle ADB$.

Prove that lines $KP, LQ,$ and $MD$ are concurrent.
\item \href{https://artofproblemsolving.com/community/c6h3359737p31218405}{\textbf{2023 G3}} Let $ABCD$ be a cyclic quadrilateral with $\angle BAD < \angle ADC$. Let $M$ be the midpoint of the arc $CD$ not containing $A$. Suppose there is a point $P$ inside $ABCD$ such that $\angle ADB = \angle CPD$ and $\angle ADP = \angle PCB$.

Prove that lines $AD, PM$, and $BC$ are concurrent.
    \item \href{https://artofproblemsolving.com/community/c6h3107322p28104239}{\textbf{2022 G3}} Let $ABCD$ be a cyclic quadrilateral. Assume that the points $Q, A, B, P$ are collinear in this order, in such a way that the line $AC$ is tangent to the circle $ADQ$, and the line $BD$ is tangent to the circle $BCP$. Let $M$ and $N$ be the midpoints of segments $BC$ and $AD$, respectively. 
    
    Prove that the following three lines are concurrent: line $CD$, the tangent of circle $ANQ$ at point $A$, and the tangent to circle $BMP$ at point $B$.
\item \href{https://artofproblemsolving.com/community/c6h2882549p25627540}{\textbf{2021 G3}} Consider a $100\times 100$ square unit lattice $\textbf{L}$ (hence $\textbf{L}$ has $10000$ points). Suppose $\mathcal{F}$ is a set of polygons such that all vertices of polygons in $\mathcal{F}$ lie in $\textbf{L}$ and every point in $\textbf{L}$ is the vertex of exactly one polygon in $\mathcal{F}.$ 

Find the maximum possible sum of the areas of the polygons in $\mathcal{F}.$
\item \href{https://artofproblemsolving.com/community/c6h2625881p22698213}{\textbf{2020 G3}} Let $ABCD$ be a convex quadrilateral with $\angle ABC>90$, $CDA>90$ and $\angle DAB=\angle BCD$. Denote by $E$ and $F$ the reflections of $A$ in lines $BC$ and $CD$, respectively. Suppose that the segments $AE$ and $AF$ meet the line $BD$ at $K$ and $L$, respectively. 

Prove that the circumcircles of triangles $BEK$ and $DFL$ are tangent to each other.
\item \boxed{\href{https://artofproblemsolving.com/community/c6h1876070p12744870}{\textbf{2019 G3}}}  	In triangle $ABC$, point $A_1$ lies on side $BC$ and point $B_1$ lies on side $AC$. Let $P$ and $Q$ be points on segments $AA_1$ and $BB_1$, respectively, such that $PQ$ is parallel to $AB$. Let $P_1$ be a point on line $PB_1$, such that $B_1$ lies strictly between $P$ and $P_1$, and $\angle PP_1C=\angle BAC$. Similarly, let $Q_1$ be the point on line $QA_1$, such that $A_1$ lies strictly between $Q$ and $Q_1$, and $\angle CQ_1Q=\angle CBA$. 

Prove that points $P,Q,P_1$, and $Q_1$ are concyclic.

\newpage

\item \href{https://artofproblemsolving.com/community/c6h1876760p12752814}{\textbf{2018 G3}} A circle $\omega$ with radius $1$ is given. A collection $T$ of triangles is called good, if the following conditions hold: each triangle from $T$ is inscribed in $\omega$; no two triangles from $T$ have a common interior point. 

Determine all positive real numbers $t$ such that, for each positive integer $n$, there exists a good collection of $n$ triangles, each of perimeter greater than $t$.
\item \href{https://artofproblemsolving.com/community/c6h1671271p10632285}{\textbf{2017 G3}} Let $O$ be the circumcenter of an acute triangle $ABC$. Line $OA$ intersects the altitudes of $ABC$ through $B$ and $C$ at $P$ and $Q$, respectively. The altitudes meet at $H$. 

Prove that the circumcenter of triangle $PQH$ lies on a median of triangle $ABC$.

\item \href{https://artofproblemsolving.com/community/c6h1480713p8639302}{\textbf{2016 G3}} Let $B = (-1, 0)$ and $C = (1, 0)$ be fixed points on the coordinate plane. A nonempty, bounded subset $S$ of the plane is said to be nice if
\begin{enumerate}[(i)]
    \item there is a point $T$ in $S$ such that for every point $Q$ in $S$, the segment $TQ$ lies entirely in $S$; and
    \item for any triangle $P_1P_2P_3$, there exists a unique point $A$ in $S$ and a permutation $\sigma$ of the indices $\{1, 2, 3\}$ for which triangles $ABC$ and $P_{\sigma(1)}P_{\sigma(2)}P_{\sigma(3)}$ are similar.
\end{enumerate}
Prove that there exist two distinct nice subsets $S$ and $S'$ of the set $\{(x, y) : x \geqslant  0, y \geqslant  0\}$ such that if $A \in S$ and $A' \in S'$ are the unique choices of points in $\text{(ii)}$, then the product $BA \cdot BA'$ is a constant independent of the triangle $P_1P_2P_3$.
\item \href{https://artofproblemsolving.com/community/c6h1268851p6622204}{\textbf{2015 G3}} Let $ABC$ be a triangle with $\angle{C} = 90^{\circ}$, and let $H$ be the foot of the altitude from $C$. A point $D$ is chosen inside the triangle $CBH$ so that $CH$ bisects $AD$. Let $P$ be the intersection point of the lines $BD$ and $CH$. Let $\omega$ be the semicircle with diameter $BD$ that meets the segment $CB$ at an interior point. A line through $P$ is tangent to $\omega$ at $Q$. 

Prove that the lines $CQ$ and $AD$ meet on $\omega$.

\item \href{https://artofproblemsolving.com/community/c6h1113191p5083558}{\textbf{2014 G3}} Let $\Omega$ and $O$ be the circumcircle and the circumcentre of an acute-angled triangle $ABC$ with $AB > BC$. The angle bisector of $\angle ABC$ intersects $\Omega$ at $M \ne B$. Let $\Gamma$ be the circle with diameter $BM$. The angle bisectors of $\angle AOB$ and $\angle BOC$ intersect $\Gamma$ at points $P$ and $Q,$ respectively. The point $R$ is chosen on the line $P Q$ so that $BR = MR$. 

Prove that $BR\parallel AC$.
(Here we always assume that an angle bisector is a ray.)

\item \href{https://artofproblemsolving.com/community/c6h597239p3544090}{\textbf{2013 G3}} In a triangle $ABC$, let $D$ and $E$ be the feet of the angle bisectors of angles $A$ and $B$, respectively. A rhombus is inscribed into the quadrilateral $AEDB$ (all vertices of the rhombus lie on different sides of $AEDB$). Let $\varphi$ be the non-obtuse angle of the rhombus. 

Prove that $\varphi \leqslant \max \{  \angle BAC, \angle ABC  \}$.

\item \href{https://artofproblemsolving.com/community/c6h546177p3160579}{\textbf{2012 G3}} In an acute triangle $ABC$ the points $D,E$ and $F$ are the feet of the altitudes through $A,B$ and $C$ respectively. The incenters of the triangles $AEF$ and $BDF$ are $I_1$ and $I_2$ respectively; the circumcenters of the triangles $ACI_1$ and $BCI_2$ are $O_1$ and $O_2$ respectively. 

Prove that $I_1I_2$ and $O_1O_2$ are parallel.
\newpage

\item \href{https://artofproblemsolving.com/community/c6h488827p2739325}{\textbf{2011 G3}} Let $ABCD$ be a convex quadrilateral whose sides $AD$ and $BC$ are not parallel. Suppose that the circles with diameters $AB$ and $CD$ meet at points $E$ and $F$ inside the quadrilateral. Let $\omega_E$ be the circle through the feet of the perpendiculars from $E$ to the lines $AB,BC$ and $CD$. Let $\omega_F$ be the circle through the feet of the perpendiculars from $F$ to the lines $CD,DA$ and $AB$. 

Prove that the midpoint of the segment $EF$ lies on the line through the two intersections of $\omega_E$ and $\omega_F$.
\item \href{https://artofproblemsolving.com/community/c6h418634p2361975}{\textbf{2010 G3}} Let $A_1A_2 \dots A_n$ be a convex polygon. Point $P$ inside this polygon is chosen so that its projections $P_1, \dots , P_n$ onto lines $A_1A_2, \dots , A_nA_1$ respectively lie on the sides of the polygon. 

Prove that for arbitrary points $X_1, \dots , X_n$ on sides $A_1A_2, \dots , A_nA_1$ respectively, \[\max \left\{ \frac{X_1X_2}{P_1P_2}, \dots, \frac{X_nX_1}{P_nP_1} \right\} \geqslant  1.\]
\item \href{https://artofproblemsolving.com/community/c6h355790p1932935}{\textbf{2009 G3}} Let $ABC$ be a triangle. The incircle of $ABC$ touches the sides $AB$ and $AC$ at the points $Z$ and $Y$, respectively. Let $G$ be the point where the lines $BY$ and $CZ$ meet, and let $R$ and $S$ be points such that the two quadrilaterals $BCYR$ and $BCSZ$ are parallelogram. 

Prove that $GR=GS$.
\item \href{https://artofproblemsolving.com/community/c6h287865p1555919}{\textbf{2008 G3}} 	Let $ ABCD$ be a convex quadrilateral and let $ P$ and $ Q$ be points in $ ABCD$ such that $ PQDA$ and $ QPBC$ are cyclic quadrilaterals. Suppose that there exists a point $ E$ on the line segment $ PQ$ such that $ \angle PAE = \angle QDE$ and $ \angle PBE = \angle QCE$. 

Show that the quadrilateral $ ABCD$ is cyclic.
\item \href{https://artofproblemsolving.com/community/c6h214644p1186774}{\textbf{2007 G3}} The diagonals of a trapezoid $ ABCD$ intersect at point $ P$. Point $ Q$ lies between the parallel lines $ BC$ and $ AD$ such that $ \angle AQD = \angle CQB$, and line $ CD$ separates points $ P$ and $ Q$. 

Prove that $ \angle BQP = \angle DAQ$.
\item \href{https://artofproblemsolving.com/community/c6h130813p741369}{\textbf{2006 G3}} Let $ ABCDE$ be a convex pentagon such that \[ \angle BAC = \angle CAD = \angle DAE\qquad \text{and}\qquad \angle ABC = \angle ACD = \angle ADE.\] The diagonals $BD$ and $CE$ meet at $P$. 

Prove that the line $AP$ bisects the side $CD$.
\end{itemize}
\newpage
\section*{G4 Problems}  \label{G4}
\begin{itemize}[leftmargin=*]


\item \boxed{\href{https://artofproblemsolving.com/community/c6h3866200p38664791}{\textbf{2026 G4}}} Let $ABC$ be a triangle and let points $M$ and $N$ be the midpoints of sides $AB$ and $AC$, respectively. Let points $K$ and $L$ be chosen inside triangles $BMC$ and $BNC$, respectively, such that $K$ lies inside the angle $LBA$ and $L$ lies inside the angle $ACK$. Suppose that
$$\angle KBA= \angle ACL, \quad \angle LBK= \angle LNC, \quad \text{and} \quad \angle LCK= \angle BMK.$$

Let $O$ be the circumcentre of triangle $AKL$. Prove that $OM = ON$.

\item \boxed{\href{https://artofproblemsolving.com/community/c6h3609790p35332018}{\textbf{2025 G4}}}  Let $\Omega$ and $\Gamma$ be circles with centres $M$ and $N$, respectively, such that the radius of $\Omega$ is less than the radius of $\Gamma$. Suppose $\Omega$ and $\Gamma$ intersect at two distinct points $A$ and $B$. Line $MN$ intersects $\Omega$ at $C$ and $\Gamma$ at $D$, so that $C, M, N, D$ lie on $MN$ in that order. Let $P$ be the circumcentre of triangle $ACD$. Line $AP$ meets $\Omega$ again at $E\neq A$ and meets $\Gamma$ again at $F\neq A$. Let $H$ be the orthocentre of triangle $PMN$.

Prove that the line through $H$ parallel to $AP$ is tangent to the circumcircle of triangle $BEF$.
\item \href{https://artofproblemsolving.com/community/c6h3610440p35340910}{\textbf{2024 G4}} Let $ABCD$ be a quadrilateral with $AB$ parallel to $CD$ and $AB<CD$. Lines $AD$ and $BC$ intersect at a point $P$. Point $X$ distinct from $C$ lies on the circumcircle of triangle $ABC$ such that $PC=PX$. Point $Y$ distinct from $D$ lies on the circumcircle of triangle $ABD$ such that $PD=PY$. Lines $AX$ and $BY$ intersect at $Q$.

Prove that $PQ$ is parallel to $AB$.
\item \boxed{\href{https://artofproblemsolving.com/community/c6h3106748p28097552}{\textbf{2023 G4}} } Let $ABC$ be an acute-angled triangle with $AB < AC$. Let $\Omega$ be the circumcircle of $ABC$. Let $S$ be the midpoint of the arc $CB$ of $\Omega$ containing $A$. The perpendicular from $A$ to $BC$ meets $BS$ at $D$ and meets $\Omega$ again at $E \neq A$. The line through $D$ parallel to $BC$ meets line $BE$ at $L$. Denote the circumcircle of triangle $BDL$ by $\omega$. Let $\omega$ meet $\Omega$ again at $P \neq B$. 

Prove that the line tangent to $\omega$ at $P$ meets line $BS$ on the internal angle bisector of $\angle BAC$.
    \item  \href{https://artofproblemsolving.com/community/c6h3107324p28104244}{\textbf{2022 G4}} Let $ABC$ be an acute-angled triangle with $AC > AB$, let $O$ be its circumcentre, and let $D$ be a point on the segment $BC$. The line through $D$ perpendicular to $BC$ intersects the lines $AO, AC,$ and $AB$ at $W, X,$ and $Y,$ respectively. The circumcircles of triangles $AXY$ and $ABC$ intersect again at $Z \ne A$. Prove that if $W \ne D$ and $OW = OD,$ then $DZ$ is tangent to the circle $AXY.$
\item \href{https://artofproblemsolving.com/community/c6h2882537p25627497}{\textbf{2021 G4}} Let $ABCD$ be a quadrilateral inscribed in a circle $\Omega.$ Let the tangent to $\Omega$ at $D$ meet rays $BA$ and $BC$ at $E$ and $F,$ respectively. A point $T$ is chosen inside $\triangle ABC$ so that $\overline{TE}\parallel\overline{CD}$ and $\overline{TF}\parallel\overline{AD}.$ Let $K\ne D$ be a point on segment $DF$ satisfying $TD=TK.$ 

Prove that lines $AC,DT,$ and $BK$ are concurrent.
\item \href{https://artofproblemsolving.com/community/c6h2625900p22698346}{\textbf{2020 G4}} In the plane, there are $n \geqslant  6$ pairwise disjoint disks $D_{1}, D_{2}, \dots, D_{n}$ with radii $R_{1} \geqslant  R_{2} \geqslant  \dots \geqslant  R_{n}$. For every $i=1,2, \dots, n$, a point $P_{i}$ is chosen in disk $D_{i}$. Let $O$ be an arbitrary point in the plane. Prove that\[O P_{1}+O P_{2}+\dots+O P_{n} \geqslant  R_{6}+R_{7}+\dots+R_{n}.\](A disk is assumed to contain its boundary.)
\item \href{https://artofproblemsolving.com/community/c6h2278998p17828733}{\textbf{2019 G4}} Let $P$ be a point inside triangle $ABC$. Let $AP$ meet $BC$ at $A_1$, let $BP$ meet $CA$ at $B_1$, and let $CP$ meet $AB$ at $C_1$. Let $A_2$ be the point such that $A_1$ is the midpoint of $PA_2$, let $B_2$ be the point such that $B_1$ is the midpoint of $PB_2$, and let $C_2$ be the point such that $C_1$ is the midpoint of $PC_2$. 

Prove that points $A_2, B_2$, and $C_2$ cannot all lie strictly inside the circumcircle of triangle $ABC$.
\item \href{https://artofproblemsolving.com/community/c6h1876789p12752914}{\textbf{2018 G4}} A point $T$ is chosen inside a triangle $ABC$. Let $A_1$, $B_1$, and $C_1$ be the reflections of $T$ in $BC$, $CA$, and $AB$, respectively. Let $\Omega$ be the circumcircle of the triangle $A_1B_1C_1$. The lines $A_1T$, $B_1T$, and $C_1T$ meet $\Omega$ again at $A_2$, $B_2$, and $C_2$, respectively. 

Prove that the lines $AA_2$, $BB_2$, and $CC_2$ are concurrent on $\Omega$.
\item \href{https://artofproblemsolving.com/community/c6h1671273p10632290}{\textbf{2017 G4}} 	In triangle $ABC$, let $\omega$ be the excircle opposite to $A$. Let $D, E$ and $F$ be the points where $\omega$ is tangent to $BC, CA$, and $AB$, respectively. The circle $AEF$ intersects line $BC$ at $P$ and $Q$. Let $M$ be the midpoint of $AD$. 

Prove that the circle $MPQ$ is tangent to $\omega$.
\item \href{https://artofproblemsolving.com/community/c6h1480714p8639304}{\textbf{2016 G4}} 	Let $ABC$ be a triangle with $AB = AC \neq BC$ and let $I$ be its incentre. The line $BI$ meets $AC$ at $D$, and the line through $D$ perpendicular to $AC$ meets $AI$ at $E$. 

Prove that the reflection of $I$ in $AC$ lies on the circumcircle of triangle $BDE$.
\item \href{https://artofproblemsolving.com/community/c6h1268843p6622159}{\textbf{2015 G4}} Let $ABC$ be an acute triangle and let $M$ be the midpoint of $AC$. A circle $\omega$ passing through $B$ and $M$ meets the sides $AB$ and $BC$ at points $P$ and $Q$ respectively. Let $T$ be the point such that $BPTQ$ is a parallelogram. Suppose that $T$ lies on the circumcircle of $ABC$. 

Determine all possible values of $\displaystyle frac{BT}{BM}$.
\item \href{https://artofproblemsolving.com/community/c6h1113192p5083561}{\textbf{2014 G4}} 	Consider a fixed circle $\Gamma$ with three fixed points $A, B,$ and $C$ on it. Also, let us fix a real number $\lambda \in(0,1)$. For a variable point $P \not\in\{A, B, C\}$ on $\Gamma$, let $M$ be the point on the segment $CP$ such that $CM =\lambda\cdot  CP$ . Let $Q$ be the second point of intersection of the circumcircles of the triangles $AMP$ and $BMC$. Prove that as $P$ varies, the point $Q$ lies on a fixed circle.
\item \href{https://artofproblemsolving.com/community/c6h597240p3544091}{\textbf{2013 G4}} Let $ABC$ be a triangle with $\angle B > \angle C$. Let $P$ and $Q$ be two different points on line $AC$ such that $\angle PBA = \angle QBA = \angle ACB $ and $A$ is located between $P$ and $C$. Suppose that there exists an interior point $D$ of segment $BQ$ for which $PD=PB$. Let the ray $AD$ intersect the circle $ABC$ at $R \neq A$. 

Prove that $QB = QR$.
\item \href{https://artofproblemsolving.com/community/c6h597240p3544091}{\textbf{2012 G4}} Let $ABC$ be a triangle with $AB \neq AC$ and circumcenter $O$. The bisector of $\angle BAC$ intersects $BC$ at $D$. Let $E$ be the reflection of $D$ with respect to the midpoint of $BC$. The lines through $D$ and $E$ perpendicular to $BC$ intersect the lines $AO$ and $AD$ at $X$ and $Y$ respectively. 

Prove that the quadrilateral $BXCY$ is cyclic.
\item \href{https://artofproblemsolving.com/community/c6h488829p2739327}{\textbf{2011 G4}} 	Let $ABC$ be an acute triangle with circumcircle $\Omega$. Let $B_0$ be the midpoint of $AC$ and let $C_0$ be the midpoint of $AB$. Let $D$ be the foot of the altitude from $A$ and let $G$ be the centroid of the triangle $ABC$. Let $\omega$ be a circle through $B_0$ and $C_0$ that is tangent to the circle $\Omega$ at a point $X\not= A$. 

Prove that the points $D,G$ and $X$ are collinear.
\item \boxed{\href{https://artofproblemsolving.com/community/c6h356083p1935927}{\textbf{2010 G4}}}  Given a triangle $ABC$, with $I$ as its incenter and $\Gamma$ as its circumcircle, $AI$ intersects $\Gamma$ again at $D$. Let $E$ be a point on the arc $BDC$, and $F$ a point on the segment $BC$, such that $\angle BAF=\angle CAE < \dfrac12\angle BAC$. If $G$ is the midpoint of $IF$, prove that the meeting point of the lines $EI$ and $DG$ lies on $\Gamma$.
\item \href{https://artofproblemsolving.com/community/c6h355791p1932936}{\textbf{2009 G4}} Given a cyclic quadrilateral $ABCD$, let the diagonals $AC$ and $BD$ meet at $E$ and the lines $AD$ and $BC$ meet at $F$. The midpoints of $AB$ and $CD$ are $G$ and $H$, respectively. 

Show that $EF$ is tangent at $E$ to the circle through the points $E$, $G$ and $H$.
\item \href{https://artofproblemsolving.com/community/c6h287866p1555921}{\textbf{2008 G4}} 	In an acute triangle $ ABC$ segments $ BE$ and $ CF$ are altitudes. Two circles passing through the point $ A$ and $ F$ and tangent to the line $ BC$ at the points $ P$ and $ Q$ so that $ B$ lies between $ C$ and $ Q$. 

Prove that lines $ PE$ and $ QF$ intersect on the circumcircle of triangle $AEF$.
\item \boxed{\href{https://artofproblemsolving.com/community/c6h159837p893744}{\textbf{2007 G4}}}  Consider five points $ A$, $ B$, $ C$, $ D$ and $ E$ such that $ ABCD$ is a parallelogram and $ BCED$ is a cyclic quadrilateral. Let $ \ell$ be a line passing through $ A$. Suppose that $ \ell$ intersects the interior of the segment $ DC$ at $ F$ and intersects line $ BC$ at $ G$. Suppose also that $ EF = EG = EC$. 

Prove that $ \ell$ is the bisector of angle $ DAB$.
\item \href{https://artofproblemsolving.com/community/c6h149173p842901}{\textbf{2006 G4}} A point $D$ is chosen on the side $AC$ of a triangle $ABC$ with $\angle C < \angle A < 90^\circ$ in such a way that $BD=BA$. The incircle of $ABC$ is tangent to $AB$ and $AC$ at points $K$ and $L$, respectively. Let $J$ be the incenter of triangle $BCD$. 

Prove that the line $KL$ intersects the line segment $AJ$ at its midpoint.
\end{itemize}
\newpage
\section*{G5 Problems}  \label{G5}
\begin{itemize}[leftmargin=*]


\item \href{https://www.imo-official.org/assets/documents/problems/2026/IMO2026SL.pdf}{\textbf{2026 G5}} TO BE REVEALED!

\item \href{https://artofproblemsolving.com/community/c6h3866849p38676933}{\textbf{2025 G5}} Let $ABCD$ be a convex quadrilateral with $\angle ADC > 90^\circ$. Suppose that points $X$ and $Y$ lie on diagonal $AC$ such that 
    \[ \angle ADX = \angle YDC = 90^\circ, \quad \angle CAD = \angle CBX, \quad \text{and} \quad \angle DCA = \angle YBA. \]
    Let $O$ be the circumcircle of triangle $BXY$. Suppose that $R\ne O$ lies on the perpendicular bisector of $AC$ such that $OR$ and $AC$ are parallel.

    Prove that line $RO$ bisects angle $\angle DRB$.


\item \href{https://artofproblemsolving.com/community/c6h3610468p35340974}{\textbf{2024 G5}} Let $ABC$ be a triangle with incentre $I$, and let $\Omega$ be the circumcircle of triangle $BIC$. Let $K$ be a point in the interior of segment $BC$ such that $\angle BAK < \angle KAC$. The angle bisector of $\angle BKA$ intersects $\Omega$ at points $W$ and $X$ such that $A$ and $W$ lie on the same side of $BC$, and the angle bisector of $\angle CKA$ intersects $\Omega$ at points $Y$ and $Z$ such that $A$ and $Y$ lie on the same side of $BC$.

Prove that $\angle WAY = \angle ZAX$.
\item \href{https://artofproblemsolving.com/community/c6h3359731p31218385}{\textbf{2023 G5}} Let $ABC$ be an acute-angled triangle with circumcircle $\omega$ and circumcentre $O$. Points $D\neq B$ and $E\neq C$ lie on $\omega$ such that $BD\perp AC$ and $CE\perp AB$. Let $CO$ meet $AB$ at $X$, and $BO$ meet $AC$ at $Y$.

Prove that the circumcircles of triangles $BXD$ and $CYE$ have an intersection lie on line $AO$.
    \item  \href{https://artofproblemsolving.com/community/c6h3107344p28104319}{\textbf{2022 G5}} Let $ABC$ be a triangle and $\ell_1,\ell_2$ be two parallel lines. Let $\ell_i$ intersects line $BC,CA,AB$ at $X_i,Y_i,Z_i$, respectively. Let $\Delta_i$ be the triangle formed by the line passed through $X_i$ and perpendicular to $BC$, the line passed through $Y_i$ and perpendicular to $CA$, and the line passed through $Z_i$ and perpendicular to $AB$. Prove that the circumcircles of $\Delta_1$ and $\Delta_2$ are tangent.
\item \href{https://artofproblemsolving.com/community/c6h2882556p25627574}{\textbf{2021 G5}} Let $ABCD$ be a cyclic quadrilateral whose sides have pairwise different lengths. Let $O$ be the circumcenter of $ABCD$. The internal angle bisectors of $\angle ABC$ and $\angle ADC$ meet $AC$ at $B_1$ and $D_1$, respectively. Let $O_B$ be the center of the circle which passes through $B$ and is tangent to $\overline{AC}$ at $D_1$. Similarly, let $O_D$ be the center of the circle which passes through $D$ and is tangent to $\overline{AC}$ at $B_1$.
Assume that $\overline{BD_1} \parallel \overline{DB_1}$. 

Prove that $O$ lies on the line $\overline{O_BO_D}$.
\item \href{https://artofproblemsolving.com/community/c6h2625889p22698285}{\textbf{2020 G5}} 	Let $ABCD$ be a cyclic quadrilateral. Points $K, L, M, N$ are chosen on $AB, BC, CD, DA$ such that $KLMN$ is a rhombus with $KL \parallel AC$ and $LM \parallel BD$. Let $\omega_A, \omega_B, \omega_C, \omega_D$ be the incircles of $\triangle ANK, \triangle BKL, \triangle CLM, \triangle DMN$. 

Prove that the common internal tangents to $\omega_A$, and $\omega_C$ and the common internal tangents to $\omega_B$ and $\omega_D$ are concurrent.
\item \href{https://artofproblemsolving.com/community/c6h2279008p17828826}{\textbf{2019 G5}} Let $ABCDE$ be a convex pentagon with $CD= DE$ and $\angle EDC \ne 2 \cdot \angle ADB$.
Suppose that a point $P$ is located in the interior of the pentagon such that $AP =AE$ and $BP= BC$.
Prove that $P$ lies on the diagonal $CE$ if and only if area $(BCD)$ + area $(ADE)$ = area $(ABD)$ + area $(ABP)$.
\item \href{https://artofproblemsolving.com/community/c6h1876775p12752854}{\textbf{2018 G5}} Let $ABC$ be a triangle with circumcircle $\Omega$ and incentre $I$. A line $\ell$ intersects the lines $AI$, $BI$, and $CI$ at points $D$, $E$, and $F$, respectively, distinct from the points $A$, $B$, $C$, and $I$. The perpendicular bisectors $x$, $y$, and $z$ of the segments $AD$, $BE$, and $CF$, respectively determine a triangle $\Theta$. 

Show that the circumcircle of the triangle $\Theta$ is tangent to $\Omega$.
\item \href{https://artofproblemsolving.com/community/c6h1671278p10632312}{\textbf{2017 G5}} Let $ABCC_1B_1A_1$ be a convex hexagon such that $AB=BC$, and suppose that the line segments $AA_1, BB_1$, and $CC_1$ have the same perpendicular bisector. Let the diagonals $AC_1$ and $A_1C$ meet at $D$, and denote by $\omega$ the circle $ABC$. Let $\omega$ intersect the circle $A_1BC_1$ again at $E \neq B$. 

Prove that the lines $BB_1$ and $DE$ intersect on $\omega$.


\item \href{https://artofproblemsolving.com/community/c6h1480717p8639313}{\textbf{2016 G5}} Let $D$ be the foot of perpendicular from $A$ to the Euler line (the line passing through the circumcentre and the orthocentre) of an acute scalene triangle $ABC$. A circle $\omega$ with centre $S$ passes through $A$ and $D$, and it intersects sides $AB$ and $AC$ at $X$ and $Y$ respectively. Let $P$ be the foot of altitude from $A$ to $BC$, and let $M$ be the midpoint of $BC$. 

Prove that the circumcentre of triangle $XSY$ is equidistant from $P$ and $M$.
\item \href{https://artofproblemsolving.com/community/c6h1268908p6622796}{\textbf{2015 G5}} 	Let $ABC$ be a triangle with $CA \neq CB$. Let $D$, $F$, and $G$ be the midpoints of the sides $AB$, $AC$, and $BC$ respectively. A circle $\Gamma$ passing through $C$ and tangent to $AB$ at $D$ meets segments $AF$ and $BG$ at $H$ and $I$, respectively. The points $H'$ and $I'$ are symmetric to $H$ and $I$ about $F$ and $G$, respectively. The line $H'I'$ meets $CD$ and $FG$ at $Q$ and $M$, respectively. The line $CM$ meets $\Gamma$ again at $P$.

Prove that $CQ = QP$.
\item \boxed{\href{https://artofproblemsolving.com/community/c6h596927p3542092}{\textbf{2014 G5}}}  Convex quadrilateral $ABCD$ has $\angle ABC = \angle CDA = 90^{\circ}$. Point $H$ is the foot of the perpendicular from $A$ to $BD$. Points $S$ and $T$ lie on sides $AB$ and $AD$, respectively, such that $H$ lies inside triangle $SCT$ and\[
\angle CHS - \angle CSB = 90^{\circ}, \quad \angle THC - \angle DTC = 90^{\circ}. \]Prove that line $BD$ is tangent to the circumcircle of triangle $TSH$.
\item \href{https://artofproblemsolving.com/community/c6h597241p3544093}{\textbf{2013 G5}} Let $ABCDEF$ be a convex hexagon with $AB=DE$, $BC=EF$, $CD=FA$, and $\angle A-\angle D = \angle C -\angle F = \angle E -\angle B$. Prove that the diagonals $AD$, $BE$, and $CF$ are concurrent.

\item \boxed{\href{https://artofproblemsolving.com/community/c6h488511p2737425}{\textbf{2012 G5}}}  Let $ABC$ be a triangle with $\angle BCA=90^{\circ}$, and let $D$ be the foot of the altitude from $C$. Let $X$ be a point in the interior of the segment $CD$. Let $K$ be the point on the segment $AX$ such that $BK=BC$. Similarly, let $L$ be the point on the segment $BX$ such that $AL=AC$. Let $M$ be the point of intersection of $AL$ and $BK$. 

Show that $MK=ML$.
\item \href{https://artofproblemsolving.com/community/c6h488830p2739330}{\textbf{2011 G5}} Let $ABC$ be a triangle with incentre $I$ and circumcircle $\omega$. Let $D$ and $E$ be the second intersection points of $\omega$ with $AI$ and $BI$, respectively. The chord $DE$ meets $AC$ at a point $F$, and $BC$ at a point $G$. Let $P$ be the intersection point of the line through $F$ parallel to $AD$ and the line through $G$ parallel to $BE$. Suppose that the tangents to $\omega$ at $A$ and $B$ meet at a point $K$. 

Prove that the three lines $AE,BD$ and $KP$ are either parallel or concurrent.
\item \href{https://artofproblemsolving.com/community/c6h418635p2361976}{\textbf{2010 G5}} 	Let $ABCDE$ be a convex pentagon such that $BC \parallel AE,$ $AB = BC +  AE,$ and $\angle ABC = \angle CDE.$ Let $M$ be the midpoint of $CE,$ and let $O$ be the circumcenter of triangle $BCD.$ 

Given that $\angle DMO = 90^{\circ},$ prove that $2 \angle BDA = \angle CDE.$
\item \href{https://artofproblemsolving.com/community/c6h355792p1932937}{\textbf{2009 G5}} 	Let $P$ be a polygon that is convex and symmetric to some point $O$. Prove that for some parallelogram $R$ satisfying $P\subset R$ we have\[\frac{|R|}{|P|}\leqslant \sqrt 2\]
where $|R|$ and $|P|$ denote the area of the sets $R$ and $P$, respectively.

\item \href{https://artofproblemsolving.com/community/c6h287867p1555923}{\textbf{2008 G5}} Let $ k$ and $ n$ be integers with $ 0\leqslant k\leqslant n - 2$. Consider a set $ L$ of $ n$ lines in the plane such that no two of them are parallel and no three have a common point. Denote by $ I$ the set of intersections of lines in $ L$. Let $ O$ be a point in the plane not lying on any line of $ L$. A point $ X\in I$ is colored red if the open line segment $ OX$ intersects at most $ k$ lines in $ L$. 

Prove that $ I$ contains at least $ \dfrac{1}{2}(k + 1)(k + 2)$ red points.
\item \href{https://artofproblemsolving.com/community/c6h214646p1186790}{\textbf{2007 G5}} Let $ ABC$ be a fixed triangle, and let $ A_1$, $ B_1$, $ C_1$ be the midpoints of sides $ BC$, $ CA$, $ AB$, respectively. Let $ P$ be a variable point on the circumcircle. Let lines $ PA_1$, $ PB_1$, $ PC_1$ meet the circumcircle again at $ A'$, $ B'$, $ C'$, respectively. Assume that the points $ A$, $ B$, $ C$, $ A'$, $ B'$, $ C'$ are distinct, and lines $ AA'$, $ BB'$, $ CC'$ form a triangle.

Prove that the area of this triangle does not depend on $ P$.
\item \href{https://artofproblemsolving.com/community/c6h155708p875023}{\textbf{2006 G5}} 	In triangle $ABC$, let $J$ be the center of the excircle tangent to side $BC$ at $A_{1}$ and to the extensions of the sides $AC$ and $AB$ at $B_{1}$ and $C_{1}$ respectively. Suppose that the lines $A_{1}B_{1}$ and $AB$ are perpendicular and intersect at $D$. Let $E$ be the foot of the perpendicular from $C_{1}$ to line $DJ$. 

Determine the angles $\angle{BEA_{1}}$ and $\angle{AEB_{1}}$.
\end{itemize}
\newpage
\section*{G6 Problems}  \label{G6}
\begin{itemize}[leftmargin=*]

\item \href{https://www.imo-official.org/assets/documents/problems/2026/IMO2026SL.pdf}{\textbf{2026 G6}} TO BE REVEALED!

\item \href{https://artofproblemsolving.com/community/c6h3866848p38676932}{\textbf{2025 G6}} Let $ABC$ be a triangle with $\angle ABC > \angle CAB \geqslant 60^\circ$. Let $\Omega$ be the circumcircle of triangle $ABC$. The lines tangent to $\Omega$ at points $B$ and $C$ intersect at $P$. Lines $BP$ and $AC$ intersect at $E$. Lines $CP$ and $AB$ intersect at $F$. Let $D$ be the reflection of point $A$ in line $EF$. The circumcircles of triangles $DEB$ and $DFC$ intersect again at $Q \ne D$. Let $Z$ be the intersection of lines $DP$ and $BC$.

Prove that points $A$, $Z$, and $Q$ are collinear.


\item \href{https://artofproblemsolving.com/community/c6h3610439p35340907}{\textbf{2024 G6}} Let $ABC$ be an acute triangle with $AB < AC$, and let $\Gamma$ be the circumcircle of $ABC$. Points $X$ and $Y$ lie on $\Gamma$ so that $XY$ and $BC$ meet on the external angle bisector of $\angle BAC$. Suppose that the tangents to $\Gamma$ at $X$ and $Y$ intersect at a point $T$ on the same side of $BC$ as $A$, and that $TX$ and $TY$ intersect $BC$ at $U$ and $V$, respectively. Let $J$ be the centre of the excircle of triangle $TUV$ opposite the vertex $T$.

Prove that $AJ$ bisects $\angle BAC$.
\item \href{https://artofproblemsolving.com/community/c6h3359733p31218391}{\textbf{2023 G6}} Let $ABC$ be an acute-angled triangle with circumcircle $\omega$. A circle $\Gamma$ is internally tangent to $\omega$ at $A$ and also tangent to $BC$ at $D$. Let $AB$ and $AC$ intersect $\Gamma$ at $P$ and $Q$ respectively. Let $M$ and $N$ be points on line $BC$ such that $B$ is the midpoint of $DM$ and $C$ is the midpoint of $DN$. Lines $MP$ and $NQ$ meet at $K$ and intersect $\Gamma$ again at $I$ and $J$ respectively. The ray $KA$ meets the circumcircle of triangle $IJK$ again at $X\neq K$.

Prove that $\angle BXP = \angle CXQ$.

\item  \href{https://artofproblemsolving.com/community/c6h3107320p28104233}{\textbf{2022 G6}} Let $ABC$ be an acute triangle with altitude $\overline{AH}$, and let $P$ be a variable point such that the angle bisectors $k$ and $\ell$ of $\angle PBC$ and $\angle PCB$, respectively, meet on $\overline{AH}$. Let $k$ meet $\overline{AC}$ at $E$, $\ell$ meet $\overline{AB}$ at $F$, and $\overline{EF}$ meet $\overline{AH}$ at $Q$. 
    
Prove that as $P$ varies, line $PQ$ passes through a fixed point.
\item \href{https://artofproblemsolving.com/community/c6h2882534p25627484}{\textbf{2021 G6}} Find all integers $n\geqslant  3$ for which every convex equilateral $n$-gon of side length $1$ contains an equilateral triangle of side length $1$. (Here, polygons contain their boundaries.)
\item \href{https://artofproblemsolving.com/community/c6h2625869p22698132}{\textbf{2020 G6}} Let $ABC$ be a triangle with $AB < AC$, incenter $I$, and $A$ excenter $I_{A}$. The incircle meets $BC$ at $D$. Define $E = AD\cap BI_{A}$, $F = AD\cap CI_{A}$. 

Show that the circumcircle of $\triangle AID$ and $\triangle I_{A}EF$ are tangent to each other
\item \href{https://artofproblemsolving.com/community/c6h2279013p17828908}{\textbf{2019 G6}} Let $I$ be the incentre of acute-angled triangle $ABC$. Let the incircle meet $BC, CA$, and $AB$ at $D, E$, and $F,$ respectively. Let line $EF$ intersect the circumcircle of the triangle at $P$ and $Q$, such that $F$ lies between $E$ and $P$. 

Prove that $\angle DPA + \angle AQD =\angle QIP$.
\item \boxed{\href{https://artofproblemsolving.com/community/c6h1671293p10632360}{\textbf{2018 G6}}}  A convex quadrilateral $ABCD$ satisfies $AB\cdot CD = BC\cdot DA$. Point $X$ lies inside $ABCD$ so that\[\angle{XAB} = \angle{XCD}\quad\,\,\text{and}\quad\,\,\angle{XBC} = \angle{XDA}.\]Prove that $\angle{BXA} + \angle{DXC} = 180^\circ$.
\item \href{https://artofproblemsolving.com/community/c6h1671314p10632474}{\textbf{2017 G6}} Let $n\geqslant 3$ be an integer. Two regular $n$-gons $\mathcal{A}$ and $\mathcal{B}$ are given in the plane. Prove that the vertices of $\mathcal{A}$ that lie inside $\mathcal{B}$ or on its boundary are consecutive. (That is, prove that there exists a line separating those vertices of $\mathcal{A}$ that lie inside $\mathcal{B}$ or on its boundary from the other vertices of $\mathcal{A}$.)



\item \href{https://artofproblemsolving.com/community/c6h1480719p8639316}{\textbf{2016 G6}} 	Let $ABCD$ be a convex quadrilateral with $\angle ABC = \angle ADC < 90^{\circ}$. The internal angle bisectors of $\angle ABC$ and $\angle ADC$ meet $AC$ at $E$ and $F$ respectively, and meet each other at point $P$. Let $M$ be the midpoint of $AC$ and let $\omega$ be the circumcircle of triangle $BPD$. Segments $BM$ and $DM$ intersect $\omega$ again at $X$ and $Y$ respectively. Denote by $Q$ the intersection point of lines $XE$ and $YF$. 

Prove that $PQ \perp AC$.
\item \boxed{\href{https://artofproblemsolving.com/community/c6h1112748p5079655}{\textbf{2015 G6}}}  Let $ABC$ be an acute triangle with $AB > AC$. Let $\Gamma $ be its circumcircle, $H$ its orthocenter, and $F$ the foot of the altitude from $A$. Let $M$ be the midpoint of $BC$. Let $Q$ be the point on $\Gamma$ such that $\angle HQA = 90^{\circ}$ and let $K$ be the point on $\Gamma$ such that $\angle HKQ = 90^{\circ}$. Assume that the points $A$, $B$, $C$, $K$ and $Q$ are all different and lie on $\Gamma$ in this order. 

Prove that the circumcircles of triangles $KQH$ and $FKM$ are tangent to each other.
\item \href{https://artofproblemsolving.com/community/c6h1113193p5083563}{\textbf{2014 G6}} Let $ABC$ be a fixed acute-angled triangle. Consider some points $E$ and $F$ lying on the sides $AC$ and $AB$, respectively, and let $M$ be the midpoint of $EF$ . Let the perpendicular bisector of $EF$ intersect the line $BC$ at $K$, and let the perpendicular bisector of $MK$ intersect the lines $AC$ and $AB$ at $S$ and $T$ , respectively. We call the pair $(E, F )$ $\textit{interesting}$, if the quadrilateral $KSAT$ is cyclic. Suppose that the pairs $(E_1 , F_1 )$ and $(E_2 , F_2 )$ are interesting. 

Prove that $\displaystyle \frac{E_1 E_2}{AB}=\frac{F_1 F_2}{AC}$.
\item \boxed{\href{https://artofproblemsolving.com/community/c6h1181536p5720184}{\textbf{2013 G6}}}  Let the excircle of triangle $ABC$ opposite the vertex $A$ be tangent to the side $BC$ at the point $A_1$. Define the points $B_1$ on $CA$ and $C_1$ on $AB$ analogously, using the excircles opposite $B$ and $C$, respectively. Suppose that the circumcentre of triangle $A_1B_1C_1$ lies on the circumcircle of triangle $ABC$. Prove that triangle $ABC$ is right-angled.
\item \href{https://artofproblemsolving.com/community/c6h546182p3160589}{\textbf{2012 G6}} Let $ABC$ be a triangle with circumcenter $O$ and incenter $I$. The points $D,E$ and $F$ on the sides $BC,CA$ and $AB$ respectively are such that $BD+BF=CA$ and $CD+CE=AB$. The circumcircles of the triangles $BFD$ and $CDE$ intersect at $P \neq D$. 

Prove that $OP=OI$.
\item \href{https://artofproblemsolving.com/community/c6h488831p2739334}{\textbf{2011 G6}} Let $ABC$ be a triangle with $AB=AC$ and let $D$ be the midpoint of $AC$. The angle bisector of $\angle BAC$ intersects the circle through $D,B$ and $C$ at the point $E$ inside the triangle $ABC$. The line $BD$ intersects the circle through $A,E$ and $B$ in two points $B$ and $F$. The lines $AF$ and $BE$ meet at a point $I$, and the lines $CI$ and $BD$ meet at a point $K$. 

Show that $I$ is the incentre of triangle $KAB$.
\item \href{https://artofproblemsolving.com/community/c6h418636p2361979}{\textbf{2010 G6}} 	The vertices $X, Y , Z$ of an equilateral triangle $XYZ$ lie respectively on the sides $BC, CA, AB$ of an acute-angled triangle $ABC.$ 

Prove that the incenter of triangle $ABC$ lies inside triangle $XYZ$.

\newpage

\item \href{https://artofproblemsolving.com/community/c6h355793p1932938}{\textbf{2009 G6}} 	Let the sides $AD$ and $BC$ of the quadrilateral $ABCD$ (such that $AB$ is not parallel to $CD$) intersect at point $P$. Points $O_1$ and $O_2$ are circumcenters and points $H_1$ and $H_2$ are orthocenters of triangles $ABP$ and $CDP$, respectively. Denote the midpoints of segments $O_1H_1$ and $O_2H_2$ by $E_1$ and $E_2$, respectively. 

Prove that the perpendicular from $E_1$ on $CD$, the perpendicular from $E_2$ on $AB$ and the lines $H_1H_2$ are concurrent.
\item \href{https://artofproblemsolving.com/community/c6h287868p1555924}{\textbf{2008 G6}} There is given a convex quadrilateral $ ABCD$. Prove that there exists a point $ P$ inside the quadrilateral such that \[ \angle PAB + \angle PDC = \angle PBC + \angle PAD = \angle PCD + \angle PBA = \angle PDA + \angle PCB = 90^{\circ} \] if and only if the diagonals $ AC$ and $ BD$ are perpendicular.
\item \href{https://artofproblemsolving.com/community/c6h214651p1186813}{\textbf{2007 G6}} Determine the smallest positive real number $ k$ with the following property. Let $ ABCD$ be a convex quadrilateral, and let points $ A_1$, $ B_1$, $ C_1$, and $ D_1$ lie on sides $ AB$, $ BC$, $ CD$, and $ DA$, respectively. Consider the areas of triangles $ AA_1D_1$, $ BB_1A_1$, $ CC_1B_1$ and $ DD_1C_1$; let $ S$ be the sum of the two smallest ones, and let $ S_1$ be the area of quadrilateral $ A_1B_1C_1D_1$. Then we always have $ kS_1\geqslant S$.
\item \href{https://artofproblemsolving.com/community/c6h155710p875026}{\textbf{2006 G6}} Circles $ w_{1}$ and $ w_{2}$ with centres $ O_{1}$ and $ O_{2}$ are externally tangent at point $ D$ and internally tangent to a circle $ w$ at points $ E$ and $ F$ respectively. Line $ t$ is the common tangent of $ w_{1}$ and $ w_{2}$ at $ D$. Let $ AB$ be the diameter of $ w$ perpendicular to $ t$, so that $ A, E, O_{1}$ are on the same side of $ t$. 

Prove that lines $ AO_{1}$, $ BO_{2}$, $ EF$ and $ t$ are concurrent.
\end{itemize}
\newpage
\section*{G7 Problems}  \label{G7}
\begin{itemize}[leftmargin=*]


\item \href{https://www.imo-official.org/assets/documents/problems/2026/IMO2026SL.pdf}{\textbf{2026 G7}} TO BE REVEALED!

\item \href{https://artofproblemsolving.com/community/c6h3866846p38676930}{\textbf{2025 G7}} Let $ABCD$ be a cyclic quadrilateral with circumcircle $\Omega$ such that $CD > BD > BC$. The tangents to $\Omega$ at $B$ and $C$ meet at a point $S$. A line $\ell$ is called coastal if:
\begin{itemize}[leftmargin=*]
\item[$\bullet$] $\ell$ does not pass through $A$,
\item[$\bullet$] $\ell$ intersects ray $SD$ at a point $X_\ell$ beyond $D$, and
\item[$\bullet$] there is a circle tangent to line $\ell$, line $X_\ell A$, segment $SB$ and segment $SC$.
\end{itemize}
Prove that there is a circle that is tangent to $\Omega$ and all coastal lines. 


\item \href{https://artofproblemsolving.com/community/c6h3610452p35340925}{\textbf{2024 G7}} Let \(ABC\) be a triangle with incenter \(I\) such that \(AB<AC<BC\). The second intersections of \(AI\), \(BI\), and \(CI\) with the circumcircle of triangle \(ABC\) are \(M_{A}\), \(M_{B}\), and \(M_{C}\), respectively. Lines \(AI\) and \(BC\) intersect at \(D\) and lines \(BM_{C}\) and \(CM_{B}\) intersect at \(X\). Suppose the circumcircle of triangles \(XM_{B}M_{C}\) and \(XBC\) intersect again at \(S\neq X\). Lines \(BX\) and \(CX\) intersect the circumcircle of triangle \(SXM_{A}\) again at \(P\neq X\) and \(Q\neq X\), respectively.

Prove that the circumcenter of triangle \(SID\) lies on \(PQ\).
\item \href{https://artofproblemsolving.com/community/c6h3359736p31218400}{\textbf{2023 G7}} Let $ABC$ be an acute, scalene triangle with orthocentre $H$. Let $\ell_a$ be the line through the reflection of $B$ with respect to $CH$ and the reflection of $C$ with respect to $BH$. Lines $\ell_b$ and $\ell_c$ are defined similarly. Suppose lines $\ell_a$, $\ell_b$, and $\ell_c$ determine a triangle $\mathcal T$.

Prove that the orthocentre of $\mathcal T$, the circumcentre of $\mathcal T$, and $H$ are collinear.
    \item  \href{https://artofproblemsolving.com/community/c6h3107352p28104378}{\textbf{2022 G7}} Two triangles $ABC, A'B'C'$ have the same orthocenter $H$ and the same circumcircle with center $O$. Letting $PQR$ be the triangle formed by $AA', BB', CC'$, prove that the circumcenter of $PQR$ lies on $OH$.
\item \boxed{\href{https://artofproblemsolving.com/community/c6h2625862p22698068}{\textbf{2021 G7}}}  Let $D$ be an interior point of the acute triangle $ABC$ with $AB > AC$ so that $\angle DAB = \angle CAD.$ The point $E$ on the segment $AC$ satisfies $\angle ADE =\angle BCD,$ the point $F$ on the segment $AB$ satisfies $\angle FDA =\angle DBC,$ and the point $X$ on the line $AC$ satisfies $CX = BX.$ Let $O_1$ and $O_2$ be the circumcenters of the triangles $ADC$ and $EXD,$ respectively. 

Prove that the lines $BC, EF,$ and $O_1O_2$ are concurrent.
\item \href{https://artofproblemsolving.com/community/c6h2625883p22698237}{\textbf{2020 G7}} Let $P$ be a point on the circumcircle of acute triangle $ABC$. Let $D,E,F$ be the reflections of $P$ in the $A$-midline, $B$-midline, and $C$-midline. Let $\omega$ be the circumcircle of the triangle formed by the perpendicular bisectors of $AD, BE, CF$. 

Show that the circumcircles of $\triangle ADP, \triangle BEP, \triangle CFP,$ and $\omega$ share a common point.
\item \boxed{\href{https://artofproblemsolving.com/community/c6h1876745p12752769}{\textbf{2019 G7}}}  Let $I$ be the incentre of acute triangle $ABC$ with $AB\neq AC$. The incircle $\omega$ of $ABC$ is tangent to sides $BC, CA$, and $AB$ at $D, E,$ and $F$, respectively. The line through $D$ perpendicular to $EF$ meets $\omega$ at $R$. Line $AR$ meets $\omega$ again at $P$. The circumcircles of triangle $PCE$ and $PBF$ meet again at $Q$. 

Prove that lines $DI$ and $PQ$ meet on the line through $A$ perpendicular to $AI$.
\newpage

\item \href{https://artofproblemsolving.com/community/c6h1876778p12752860}{\textbf{2018 G7}} 	Let $O$ be the circumcentre, and $\Omega$ be the circumcircle of an acute-angled triangle $ABC$. Let $P$ be an arbitrary point on $\Omega$, distinct from $A$, $B$, $C$, and their antipodes in $\Omega$. Denote the circumcentres of the triangles $AOP$, $BOP$, and $COP$ by $O_A$, $O_B$, and $O_C$, respectively. The lines $\ell_A$, $\ell_B$, $\ell_C$ perpendicular to $BC$, $CA$, and $AB$ pass through $O_A$, $O_B$, and $O_C$, respectively. 

Prove that the circumcircle of triangle formed by $\ell_A$, $\ell_B$, and $\ell_C$ is tangent to the line $OP$.
\item \href{https://artofproblemsolving.com/community/c6h1671266p10632269}{\textbf{2017 G7}} A convex quadrilateral $ABCD$ has an inscribed circle with center $I$. Let $I_a, I_b, I_c$ and $I_d$ be the incenters of the triangles $DAB, ABC, BCD$ and $CDA$, respectively. Suppose that the common external tangents of the circles $AI_bI_d$ and $CI_bI_d$ meet at $X$, and the common external tangents of the circles $BI_aI_c$ and $DI_aI_c$ meet at $Y$. 

Prove that $\angle{XIY}=90^{\circ}$.



\item \href{https://artofproblemsolving.com/community/c6h1480720p8639324}{\textbf{2016 G7}} 	Let $I$ be the incentre of a non-equilateral triangle $ABC$, $I_A$ be the $A$-excentre, $I'_A$ be the reflection of $I_A$ in $BC$, and $l_A$ be the reflection of line $AI'_A$ in $AI$. Define points $I_B$, $I'_B$ and line $l_B$ analogously. Let $P$ be the intersection point of $l_A$ and $l_B$. Prove that $P$ lies on line $OI$ where $O$ is the circumcentre of triangle $ABC$. Let one of the tangents from $P$ to the incircle of triangle $ABC$ meet the circumcircle at points $X$ and $Y$. 

Show that $\angle XIY = 120^{\circ}$.
\item \href{https://artofproblemsolving.com/community/c6h1268837p6622082}{\textbf{2015 G7}} 	Let $ABCD$ be a convex quadrilateral, and let $P$, $Q$, $R$, and $S$ be points on the sides $AB$, $BC$, $CD$, and $DA$, respectively. Let the line segment $PR$ and $QS$ meet at $O$. Suppose that each of the quadrilaterals $APOS$, $BQOP$, $CROQ$, and $DSOR$ has an incircle. 

Prove that the lines $AC$, $PQ$, and $RS$ are either concurrent or parallel to each other.
\item \href{https://artofproblemsolving.com/community/c6h1113194p5083564}{\textbf{2014 G7}} Let $ABC$ be a triangle with circumcircle $\Omega$ and incentre $I$. Let the line passing through $I$ and perpendicular to $CI$ intersect the segment $BC$ and the arc $BC$ (not containing $A$) of $\Omega$ at points $U$ and $V$ , respectively. Let the line passing through $U$ and parallel to $AI$ intersect $AV$ at $X$, and let the line passing through $V$ and parallel to $AI$ intersect $AB$ at $Y$ . Let $W$ and $Z$ be the midpoints of $AX$ and $BC$, respectively. 

Prove that if the points $I, X,$ and $Y$ are collinear, then the points $I, W ,$ and $Z$ are also collinear.
\item \href{https://artofproblemsolving.com/community/c6h546183p3160591}{\textbf{2012 G7}} Let $ABCD$ be a convex quadrilateral with non-parallel sides $BC$ and $AD$. Assume that there is a point $E$ on the side $BC$ such that the quadrilaterals $ABED$ and $AECD$ are circumscribed. 

Prove that there is a point $F$ on the side $AD$ such that the quadrilaterals $ABCF$ and $BCDF$ are circumscribed if and only if $AB$ is parallel to $CD$.
\item \href{https://artofproblemsolving.com/community/c6h488832p2739339}{\textbf{2011 G7}} Let $ABCDEF$ be a convex hexagon all of whose sides are tangent to a circle $\omega$ with centre $O$. Suppose that the circumcircle of triangle $ACE$ is concentric with $\omega$. Let $J$ be the foot of the perpendicular from $B$ to $CD$. Suppose that the perpendicular from $B$ to $DF$ intersects the line $EO$ at a point $K$. Let $L$ be the foot of the perpendicular from $K$ to $DE$. 

Prove that $DJ=DL$.
\newpage

\item \href{https://artofproblemsolving.com/community/c6h418638p2361996}{\textbf{2010 G7}} Three circular arcs $\gamma_1, \gamma_2,$ and $\gamma_3$ connect the points $A$ and $C.$ These arcs lie in the same half-plane defined by line $AC$ in such a way that arc $\gamma_2$ lies between the arcs $\gamma_1$ and $\gamma_3.$ Point $B$ lies on the segment $AC.$ Let $h_1, h_2$, and $h_3$ be three rays starting at $B,$ lying in the same half-plane, $h_2$ being between $h_1$ and $h_3.$ For $i, j = 1, 2, 3,$ denote by $V_{ij}$ the point of intersection of $h_i$ and $\gamma_j$ (see the Figure below). Denote by $\widehat{V_{ij}V_{kj}}\widehat{V_{kl}V_{il}}$ the curved quadrilateral, whose sides are the segments $V_{ij}V_{il},$ $V_{kj}V_{kl}$ and arcs $V_{ij}V_{kj}$ and $V_{il}V_{kl}.$ We say that this quadrilateral is $circumscribed$ if there exists a circle touching these two segments and two arcs. 

Prove that if the curved quadrilaterals $\widehat{V_{11}V_{21}}\widehat{V_{22}V_{12}}, \widehat{V_{12}V_{22}}\widehat{V_{23}V_{13}},\widehat{V_{21}V_{31}}\widehat{V_{32}V_{22}}$ are circumscribed, then the curved quadrilateral $\widehat{V_{22}V_{32}}\widehat{V_{33}V_{23}}$ is circumscribed, too.

\item \href{https://artofproblemsolving.com/community/c6h355794p1932939}{\textbf{2009 G7}} Let $ABC$ be a triangle with incenter $I$ and let $X$, $Y$ and $Z$ be the incenters of the triangles $BIC$, $CIA$ and $AIB$, respectively. Let the triangle $XYZ$ be equilateral. 

Prove that $ABC$ is equilateral too.




\item \boxed{\href{https://artofproblemsolving.com/community/c6h215427p1191671}{\textbf{2008 G7}}}  Let $ ABCD$ be a convex quadrilateral with $ BA\neq BC$. Denote the incircles of triangles $ ABC$ and $ ADC$ by $ \omega_{1}$ and $ \omega_{2}$ respectively. Suppose that there exists a circle $ \omega$ tangent to ray $ BA$ beyond $ A$ and to the ray $ BC$ beyond $ C$, which is also tangent to the lines $ AD$ and $ CD$. 

Prove that the common external tangents to $ \omega_{1}$ and $\omega_{2}$ intersect on $ \omega$.


\item \href{https://artofproblemsolving.com/community/c6h213519p1178972}{\textbf{2007 G7}} Given an acute triangle $ ABC$ with $ \angle B > \angle C$. Point $ I$ is the incenter, and $ R$ the circumradius. Point $ D$ is the foot of the altitude from vertex $ A$. Point $ K$ lies on line $ AD$ such that $ AK = 2R$, and $ D$ separates $ A$ and $ K$. Lines $ DI$ and $ KI$ meet sides $ AC$ and $ BC$ at $ E,F$ respectively. Let $ IE = IF$. 

Prove that $ \angle B\leqslant 3\angle C$.
\item \href{https://artofproblemsolving.com/community/c6h155711p875028}{\textbf{2006 G7}} 	In a triangle $ ABC$, let $ M_{a}$, $ M_{b}$, $ M_{c}$ be the midpoints of the sides $ BC$, $ CA$, $ AB$, respectively, and $ T_{a}$, $ T_{b}$, $ T_{c}$ be the midpoints of the arcs $ BC$, $ CA$, $ AB$ of the circumcircle of $ ABC$, not containing the vertices $ A$, $ B$, $ C$, respectively. For $ i \in \left\{a, b, c\right\}$, let $ w_{i}$ be the circle with $ M_{i}T_{i}$ as diameter. Let $ p_{i}$ be the common external common tangent to the circles $ w_{j}$ and $ w_{k}$ (for all $ \left\{i, j, k\right\}= \left\{a, b, c\right\}$) such that $ w_{i}$ lies on the opposite side of $ p_{i}$ than $ w_{j}$ and $ w_{k}$ do.

Prove that the lines $ p_{a}$, $ p_{b}$, $ p_{c}$ form a triangle similar to $ ABC$ and find the ratio of similitude.
\end{itemize}
\newpage
\section*{G8 Problems}  \label{G8}
\begin{itemize}[leftmargin=*]


\item \href{https://www.imo-official.org/assets/documents/problems/2026/IMO2026SL.pdf}{\textbf{2026 G8}} TO BE REVEALED!

\item \href{https://artofproblemsolving.com/community/c6h3610449p35340922}{\textbf{2024 G8}} Let $ABC$ be a triangle with $AB<AC<BC$, and let $D$ be a point in the interior of segment $BC$. Let $E$ be a point on the circumcircle of triangle $ABC$ such that $A$ and $E$ lie on opposite sides of the line $BC$ and $\angle{BAD}=\angle{EAC}$. Let $I,I_B,I_C,J_B$ and $J_C$ be the incenters of triangles $ABC,ABD,ADC,ABE$, and $AEC$, respectively. 

Prove that $I_B,I_C,J_B$, and $J_C$ are concyclic if and only if $AI,I_BJ_C$, and $J_BI_C$ concur.
\item \boxed{\href{https://artofproblemsolving.com/community/c6h3107345p28104331}{\textbf{2023 G8}}} Let $ABC$ be an equilateral triangle. Let $A_1,B_1,C_1$ be interior points of $ABC$ such that $BA_1=A_1C$, $CB_1=B_1A$, $AC_1=C_1B$, and
$$\angle BA_1C+\angle CB_1A+\angle AC_1B=480^\circ$$Let $BC_1$ and $CB_1$ meet at $A_2,$ let $CA_1$ and $AC_1$ meet at $B_2,$ and let $AB_1$ and $BA_1$ meet at $C_2.$
Prove that if triangle $A_1B_1C_1$ is scalene, then the three circumcircles of triangles $AA_1A_2, BB_1B_2$
and $CC_1C_2$ all pass through two common points.

(A scalene triangle is one where no two sides have equal length.)
    \item  \href{https://artofproblemsolving.com/community/c6h3107351p28104372}{\textbf{2022 G8}} Let $AA'BCC'B'$ be a convex cyclic hexagon such that $AC$ is tangent to the incircle of the triangle $A'B'C'$, and $A'C'$ is tangent to the incircle of the triangle $ABC$. Let the lines $AB$ and $A'B'$ meet at $X$ and let the lines $BC$ and $B'C'$ meet at $Y$. 
    
    Prove that if $XBYB'$ is a convex quadrilateral, then it has an incircle.
\item \href{https://artofproblemsolving.com/community/c6h2882551p25627547}{\textbf{2021 G8}} Let $ABC$ be a triangle with circumcircle $\omega$ and let $\Omega_A$ be the $A$-excircle. Let $X$ and $Y$ be the intersection points of $\omega$ and $\Omega_A$. Let $P$ and $Q$ be the projections of $A$ onto the tangent lines to $\Omega_A$ at $X$ and $Y$ respectively. The tangent line at $P$ to the circumcircle of the triangle $APX$ intersects the tangent line at $Q$ to the circumcircle of the triangle $AQY$ at a point $R$. 

Prove that $\overline{AR} \perp \overline{BC}$.
\item \href{https://artofproblemsolving.com/community/c6h2625877p22698191}{\textbf{2020 G8}} Let $ABC$ be a triangle with incenter $I$ and circumcircle $\Gamma$. Circles $\omega_{B}$ passing through $B$ and $\omega_{C}$ passing through $C$ are tangent at $I$. Let $\omega_{B}$ meet minor arc $AB$ of $\Gamma$ at $P$ and $AB$ at $M\neq B$, and let $\omega_{C}$ meet minor arc $AC$ of $\Gamma$ at $Q$ and $AC$ at $N\neq C$. Rays $PM$ and $QN$ meet at $X$. Let $Y$ be a point such that $YB$ is tangent to $\omega_{B}$ and $YC$ is tangent to $\omega_{C}$. 

Show that $A,X,Y$ are collinear.
\item \href{https://artofproblemsolving.com/community/c6h2278995p17828714}{\textbf{2019 G8}} 	Let $\mathcal L$ be the set of all lines in the plane and let $f$ be a function that assigns to each line $\ell\in\mathcal L$ a point $f(\ell)$ on $f(\ell)$. Suppose that for any point $X$, and for any three lines $\ell_1,\ell_2,\ell_3$ passing through $X$, the points $f(\ell_1),f(\ell_2),f(\ell_3)$, and $X$ lie on a circle.

Prove that there is a unique point $P$ such that $f(\ell)=P$ for any line $\ell$ passing through $P$.

\newpage

\item \href{https://artofproblemsolving.com/community/c6h1671302p10632384}{\textbf{2017 G8}} There are $2017$ mutually external circles drawn on a blackboard, such that no two are tangent and no three share a common tangent. A tangent segment is a line segment that is a common tangent to two circles, starting at one tangent point and ending at the other one. Luciano is drawing tangent segments on the blackboard, one at a time, so that no tangent segment intersects any other circles or previously drawn tangent segments. Luciano keeps drawing tangent segments until no more can be drawn. Find all possible numbers of tangent segments when Luciano stops drawing.
\item \href{https://artofproblemsolving.com/community/c6h1480722p8639326}{\textbf{2016 G8}} Let $A_1, B_1$ and $C_1$ be points on sides $BC$, $CA$ and $AB$ of an acute triangle $ABC$ respectively, such that $AA_1$, $BB_1$ and $CC_1$ are the internal angle bisectors of triangle $ABC$. Let $I$ be the incentre of triangle $ABC$, and $H$ be the orthocentre of triangle $A_1B_1C_1$. 

Show that \[AH + BH + CH \geqslant  AI + BI + CI.\]
\item \href{https://artofproblemsolving.com/community/c6h1268911p6622814}{\textbf{2015 G8}} 	A triangulation of a convex polygon $\Pi$ is a partitioning of $\Pi$ into triangles by diagonals having no common points other than the vertices of the polygon. We say that a triangulation is a Thaiangulation if all triangles in it have the same area. Prove that any two different Thaiangulations of a convex polygon $\Pi$ differ by exactly two triangles. (In other words, prove that it is possible to replace one pair of triangles in the first Thaiangulation with a different pair of triangles so as to obtain the second Thaiangulation.)
\item \href{https://artofproblemsolving.com/community/c6h546185p3160596}{\textbf{2012 G8}} Let $ABC$ be a triangle with circumcircle $\omega$ and $\ell$ a line without common points with $\omega$. Denote by $P$ the foot of the perpendicular from the center of $\omega$ to $\ell$. The side-lines $BC,CA,AB$ intersect $\ell$ at the points $X,Y,Z$ different from $P$. 

Prove that the circumcircles of the triangles $AXP$, $BYP$ and $CZP$ have a common point different from $P$ or are mutually tangent at $P$.
\item \boxed{\href{https://artofproblemsolving.com/community/c6h418983p2365045}{\textbf{2011 G8}}}  Let $ABC$ be an acute triangle with circumcircle $\Gamma$. Let $\ell$ be a tangent line to $\Gamma$, and let $\ell_a, \ell_b$ and $\ell_c$ be the lines obtained by reflecting $\ell$ in the lines $BC$, $CA$ and $AB$, respectively. 

Show that the circumcircle of the triangle determined by the lines $\ell_a, \ell_b$ and $\ell_c$ is tangent to the circle $\Gamma$.

\item \href{https://artofproblemsolving.com/community/c6h355795p1932940}{\textbf{2009 G8}} Let $ABCD$ be a circumscribed quadrilateral. Let $g$ be a line through $A$ which meets the segment $BC$ in $M$ and the line $CD$ in $N$. Denote by $I_1$, $I_2$ and $I_3$ the incenters of $\triangle ABM$, $\triangle MNC$ and $\triangle NDA$, respectively. 

Prove that the orthocenter of $\triangle I_1I_2I_3$ lies on $g$.

\item \href{https://artofproblemsolving.com/community/c6h214649p1186805}{\textbf{2007 G8}} Point $ P$ lies on side $ AB$ of a convex quadrilateral $ ABCD$. Let $ \omega$ be the incircle of triangle $ CPD$, and let $ I$ be its incenter. Suppose that $ \omega$ is tangent to the incircles of triangles $ APD$ and $ BPC$ at points $ K$ and $ L$, respectively. Let lines $ AC$ and $ BD$ meet at $ E$, and let lines $ AK$ and $ BL$ meet at $ F$. 

Prove that points $ E$, $ I$, and $ F$ are collinear.


\item \href{https://artofproblemsolving.com/community/c6h155712p875030}{\textbf{2006 G8}} Let $ABCD$ be a convex quadrilateral. A circle passing through the points $A$ and $D$ and a circle passing through the points $B$ and $C$ are externally tangent at a point $P$ inside the quadrilateral. Suppose that\[\angle{PAB}+\angle{PDC}\leqslant  90^\circ\qquad\text{and}\qquad\angle{PBA}+\angle{PCD}\leqslant  90^\circ.\]Prove that $AB+CD \geqslant BC+AD$.
\end{itemize}
\newpage
\section*{G9 Problems}  \label{G9}
\begin{itemize}[leftmargin=*]


\item \href{https://www.imo-official.org/assets/documents/problems/2026/IMO2026SL.pdf}{\textbf{2026 G9}} TO BE REVEALED!

 
    \item \boxed{\href{https://artofproblemsolving.com/community/c6h2278656p17821732}{\textbf{2020 G9}}}  Prove that there exists a positive constant $c$ such that the following statement is true:
Consider an integer $n > 1$, and a set $\mathcal S$ of $n$ points in the plane such that the distance between any two different points in $\mathcal S$ is at least 1. It follows that there is a line $\ell$ separating $\mathcal S$ such that the distance from any point of $\mathcal S$ to $\ell$ is at least $cn^{-1/3}$. (A line $\ell$ separates a set of points S if some segment joining two points in $\mathcal S$ crosses $\ell$.)

\textbf{Note.} Weaker results with $cn^{-1/3}$ replaced by $cn^{-\alpha}$ may be awarded points depending on the value of the constant $\alpha > 1/3$.
\item \href{https://artofproblemsolving.com/community/c6h155713p875036}{\textbf{2006 G9}} Points $ A_{1}$, $ B_{1}$, $ C_{1}$ are chosen on the sides $ BC$, $ CA$, $ AB$ of a triangle $ ABC$ respectively. The circumcircles of triangles $ AB_{1}C_{1}$, $ BC_{1}A_{1}$, $ CA_{1}B_{1}$ intersect the circumcircle of triangle $ ABC$ again at points $ A_{2}$, $ B_{2}$, $ C_{2}$ respectively ($ A_{2}\neq A, B_{2}\neq B, C_{2}\neq C$). Points $ A_{3}$, $ B_{3}$, $ C_{3}$ are symmetric to $ A_{1}$, $ B_{1}$, $ C_{1}$ with respect to the midpoints of the sides $ BC$, $ CA$, $ AB$ respectively. 

Prove that the triangles $ A_{2}B_{2}C_{2}$ and $ A_{3}B_{3}C_{3}$ are similar.
\end{itemize}

\newpage
\section*{G10 Problems}  \label{G10}
\begin{itemize}[leftmargin=*]
 
    \item \boxed{\href{https://artofproblemsolving.com/community/c6h101488p572824}{\textbf{2006 G10}}}  Assign to each side $b$ of a convex polygon $P$ the maximum area of a triangle that has $b$ as a side and is contained in $P$. Show that the sum of the areas assigned to the sides of $P$ is at least twice the area of $P$.
\end{itemize}



\newpage
 
\title{Number Theory IMO Shortlist}



\section*{N1 Problems} \label{N1}

\begin{itemize}[leftmargin=*]

\item \boxed{\href{https://artofproblemsolving.com/community/c6h3866213p38664896}{\textbf{2026 N1}}} There are $2026$ integers greater than $1$ written on a blackboard, not necessarily different. In a move, Confucius chooses two integers $m>1$ and $n>1$ from different places on the blackboard and replaces these two integers with
\[\gcd(m,n) \quad \text{ and } \quad  \frac{\mathrm{lcm}(m,n)}{\gcd(m,n)}.\]
He continues to make moves while it is possible to do so.
\begin{enumerate}[leftmargin=*]
\item[(a)] Prove that, regardless of the choices of Confucius, after finitely many moves, exactly one integer $M$ on the blackboard is greater than $1$.

\item[(b)] Prove that the value of $M$ does not depend on the choices of Confucius.
\end{enumerate}
%(Note that $\gcd(x,y)$ denotes the greatest common divisor of positive integers $x$ and $y$, and $\mathrm{lcm}(x,y)$ denotes the least common multiple of $x$ and $y$.)
% I think everyone knows this bruh^

\item \href{https://artofproblemsolving.com/community/c6h3866859p38676947}{\textbf{2025 N1}} Let $n$ and $k$ be positive integers such that $n<k<2n$. Suppose $a_1, a_2, a_3,\dots, a_k$ are positive integers such that   $2^{a_1}+2^{a_2}+\dots+2^{a_k}$ 
    is divisible by $2^n-1$. 
    
    Show that at least three of $a_1,a_2,\ldots, a_k$ have the same remainder when divided by $n$.


\item \href{https://artofproblemsolving.com/community/c6h3610447p35340920}{\textbf{2024 N1}} Find all positive integers $n$ with the following property: for all positive divisors $d$ of $n$, we have $d+1\mid n$ or $d+1$ is prime.	

\item \boxed{\href{https://artofproblemsolving.com/community/c6h3106752p28097575}{\textbf{2023 N1}}} Determine all composite integers $n>1$ that satisfy the following property: if $d_1$, $d_2$, $\dots$, $d_k$ are all the positive divisors of $n$ with $1 = d_1 < d_2 < \dots < d_k = n$, then $d_i$ divides $d_{i+1} + d_{i+2}$ for every $1 \leqslant i \leqslant k - 2$.	

    \item \href{https://artofproblemsolving.com/community/c6h3107332p28104270}{\textbf{2022 N1}} A number is called Norwegian if it has three distinct positive divisors whose sum is equal to $2022$. Determine the smallest Norwegian number. 
    
    (The total number of positive divisors of a Norwegian number is allowed to be larger than $3$.) 
    
    \item \href{https://artofproblemsolving.com/community/c6h2882583p25627819}{\textbf{2021 N1}} 	Find all positive integers $n\geqslant  1$ such that there exists a pair $(a,b)$ of positive integers, such that $a^2+b+3$ is not divisible by the cube of any prime, and$$n=\frac{ab+3b+8}{a^2+b+3}.$$
    
    \item \href{https://artofproblemsolving.com/community/c6h2625898p22698342}{\textbf{2020 N1}} Given a positive integer $k$. Show that there exists a prime $p$ such that one can choose distinct integers $a_1,a_2\dots, a_{k+3} \in \{1, 2, \dots ,p-1\}$ such that $p$ divides $a_ia_{i+1}a_{i+2}a_{i+3}-i$ for all $i= 1, 2, \dots, k$.
    
    \item \boxed{\href{https://artofproblemsolving.com/community/c6h1876742p12752761}{\textbf{2019 N1}}} Find all pairs $(k,n)$ of positive integers such that\[ k!=(2^n-1)(2^n-2)(2^n-4)\dots(2^n-2^{n-1}). \]
    
    \item \href{https://artofproblemsolving.com/community/c6h1876768p12752840}{\textbf{2018 N1}}  Determine all pairs $(n, k)$ of distinct positive integers such that there exists a positive integer $s$ for which the number of divisors of $sn$ and of $sk$ are equal.
    \newpage
    
    \item \boxed{\href{https://artofproblemsolving.com/community/c6h1480154p8633268}{\textbf{2017 N1}}}  For each integer $a_0 > 1$, define the sequence $a_0, a_1, a_2, \dots$ for $n \geqslant  0$ as
$$a_{n+1} = 
\begin{cases}
\sqrt{a_n} & \text{if } \sqrt{a_n} \text{ is an integer,} \\
a_n + 3 & \text{otherwise.}
\end{cases}
$$Determine all values of $a_0$ such that there exists a number $A$ such that $a_n = A$ for infinitely many values of $n$.
    \item \href{https://artofproblemsolving.com/community/c6h1480723p8639330}{\textbf{2016 N1}}  For any positive integer $k$, denote the sum of digits of $k$ in its decimal representation by $S(k)$. Find all polynomials $P(x)$ with integer coefficients such that for any positive integer $n \geqslant  2016$, the integer $P(n)$ is positive and$$S(P(n)) = P(S(n)).$$
    \item \href{https://artofproblemsolving.com/community/c6h1268859p6622268}{\textbf{2015 N1}}	Determine all positive integers $M$ such that the sequence $a_0, a_1, a_2, \dots$ defined by\[ a_0 = M + \frac{1}{2}   \qquad  \textrm{and} \qquad    a_{k+1} = a_k\lfloor a_k \rfloor   \quad \textrm{for} \, k = 0, 1, 2, \dots \]contains at least one integer term.
    \item \href{https://artofproblemsolving.com/community/c6h1113195p5083566}{\textbf{2014 N1}} Let $n \geqslant  2$ be an integer, and let $A_n$ be the set\[A_n = \{2^n  - 2^k\mid k \in \mathbb{Z},\, 0 \leqslant k < n\}.\]Determine the largest positive integer that cannot be written as the sum of one or more (not necessarily distinct) elements of $A_n$.
    \item \href{https://artofproblemsolving.com/community/c6h597243p3544096}{\textbf{2013 N1}} Let $\mathbb{Z} _{>0}$ be the set of positive integers. Find all functions $f: \mathbb{Z} _{>0}\rightarrow \mathbb{Z} _{>0}$ such that
\[ m^2 + f(n) \mid mf(m) +n \]
for all positive integers $m$ and $n$.
    \item \href{https://artofproblemsolving.com/community/c6h546186p3160599}{\textbf{2012 N1}}  Call admissible a set $A$ of integers that has the following property:
If $x,y \in A$ (possibly $x=y$) then $x^2+kxy+y^2 \in A$ for every integer $k$.
Determine all pairs $m,n$ of nonzero integers such that the only admissible set containing both $m$ and $n$ is the set of all integers.
    \item \href{https://artofproblemsolving.com/community/c6h488539p2737648}{\textbf{2011 N1}} For any integer $d > 0,$ let $f(d)$ be the smallest possible integer that has exactly $d$ positive divisors (so for example we have $f(1)=1, f(5)=16,$ and $f(6)=12$). Prove that for every integer $k \geqslant  0$ the number $f\left(2^k\right)$ divides $f\left(2^{k+1}\right).$
    \item \href{https://artofproblemsolving.com/community/c6h418639p2361998}{\textbf{2010 N1}} Find the least positive integer $n$ for which there exists a set $\{s_1, s_2, \dots , s_n\}$ consisting of $n$ distinct positive integers such that
\[ \left( 1 - \frac{1}{s_1} \right) \left( 1 - \frac{1}{s_2} \right) \dots \left( 1 - \frac{1}{s_n} \right) = \frac{51}{2010}.\]
    \item \boxed{\href{https://artofproblemsolving.com/community/c6h288838p1561571}{\textbf{2009 N1}}}  Let $ n$ be a positive integer and let $ a_1,a_2,a_3,\dots,a_k$ $ ( k\geqslant  2)$ be distinct integers in the set $ { 1,2,\dots,n}$ such that $ n$ divides $ a_i(a_{i + 1} - 1)$ for $ i = 1,2,\dots,k - 1$. Prove that $ n$ does not divide $ a_k(a_1 - 1).$
    \item \href{https://artofproblemsolving.com/community/c6h287869p1555927}{\textbf{2008 N1}} Let $n$ be a positive integer and let $p$ be a prime number. Prove that if $a$, $b$, $c$ are integers (not necessarily positive) satisfying the equations\[ a^n + pb = b^n + pc = c^n + pa\]then $a = b = c$.
    \item \href{https://artofproblemsolving.com/community/c6h201181p1106459}{\textbf{2007 N1}} 	Find all pairs of natural numbers $ (a, b)$ such that $ 7^a - 3^b$ divides $ a^4 + b^2$.
    \item \boxed{\href{https://artofproblemsolving.com/community/c6h101485p572815}{\textbf{2006 N1}}}  Determine all pairs $(x, y)$ of integers such that\[1+2^{x}+2^{2x+1}= y^{2}.\]
\end{itemize}
\newpage
 
\section*{N2 Problems} \label{N2}
\begin{itemize}[leftmargin=*]


\item \href{https://www.imo-official.org/assets/documents/problems/2026/IMO2026SL.pdf}{\textbf{2026 N2}} TO BE REVEALED!

\item \href{https://artofproblemsolving.com/community/c6h3866850p38676934}{\textbf{2025 N2}} In an $n \times n$ board, for each $1 \leqslant i \leqslant n$ and $1 \leqslant j \leqslant n$, the cell in the $i^{\text{th}}$ row from the bottom and $j^{\text{th}}$ column from the left contains $\gcd(i,j)$ leaves. A koala travels from the bottom left corner to the top right corner. At each step, the koala moves up or to the right by a single cell. The koala eats the leaves in each cell it visits, including the bottom left and top right cells.

Determine the maximum number of leaves that the koala can eat.

\item \href{https://artofproblemsolving.com/community/c6h3610444p35340915}{\textbf{2024 N2}} Determine all finite, nonempty sets $\mathcal{S}$ of positive integers such that for every $a,b\in\mathcal{S}$ there exists $c\in\mathcal{S}$ with $a\mid b+2c$.	

\item \href{https://artofproblemsolving.com/community/c6h3359734p31218394}{\textbf{2023 N2}} Determine all ordered pairs $(a,p)$ of positive integers, with $p$ prime, such that $p^a+a^4$ is a perfect square.

\item \href{https://artofproblemsolving.com/community/c6h3107356p28104399}{\textbf{2022 N2}} 	Find all positive integers $n>2$ such that
$$ n! \mid \prod_{ \substack{p<q\leqslant n \\ p,q \text{ primes }} } (p+q).$$ 
\item \boxed{\href{https://artofproblemsolving.com/community/c6h2625909p22698392}{\textbf{2021 N2}}}  Let $n \geqslant  100$ be an integer. Ivan writes the numbers $n, n+1, \dots, 2 n$ each on different cards. He then shuffles these $n+1$ cards, and divides them into two piles. Prove that at least one of the piles contains two cards such that the sum of their numbers is a perfect square.
\item \href{https://artofproblemsolving.com/community/c6h2625867p22698104}{\textbf{2020 N2}} For each prime $p$, there is a kingdom of $p$-Landia consisting of $p$ islands numbered $1,2,\dots, p$. Two distinct islands numbered $n$ and $m$ are connected by a bridge if and only if  $p$ divides $(m^{2} -n + 1)(n^{2} -m + 1)$. The bridges may pass over each other, but cannot cross. Prove that for infinitely many $p$ there are two islands in $p$-Landia not connected by a chain of bridges.

\item \href{https://artofproblemsolving.com/community/c6h2278999p17828751}{\textbf{2019 N2}} Find all triples $(a, b, c)$ of positive integers such that $a^3 + b^3 + c^3 = (abc)^2$.
\item \href{https://artofproblemsolving.com/community/c6h1876750p12752787}{\textbf{2018 N2}} 	Let $n>1$ be a positive integer. Each cell of an $n\times n$ table contains an integer. Suppose that the following conditions are satisfied:
\begin{enumerate}[(i)]
\item Each number in the table is congruent to $1$ modulo $n$.
\item The sum of numbers in any row, as well as the sum of numbers in any column, is congruent to $n$ modulo $n^2$.
 
\end{enumerate}
Let $R_i$ be the product of the numbers in the $i^{\text{th}}$ row, and $C_j$ be the product of the number in the $j^{\text{th}}$ column. Prove that the sums $R_1+\dots+ R_n$ and $C_1+\dots +C_n$ are congruent modulo $n^4$.
\newpage

\item \href{https://artofproblemsolving.com/community/c6h1671275p10632294}{\textbf{2017 N2}} Let $ p \geqslant  2$ be a prime number. Eduardo and Fernando play the following game making moves alternately: in each move, the current player chooses an index $i$ in the set $\{0,1,2,\dots, p-1 \}$ that was not chosen before by either of the two players and then chooses an element $a_i$ from the set $\{0,1,2,3,4,5,6,7,8,9\}$. Eduardo has the first move. The game ends after all the indices have been chosen .Then the following number is computed:
$$M=a_0+a_110+a_210^2+\dots+a_{p-1}10^{p-1}= \sum_{i=0}^{p-1}a_i.10^i.$$
The goal of Eduardo is to make $M$ divisible by $p$, and the goal of Fernando is to prevent this.

Prove that Eduardo has a winning strategy.
\item \href{https://artofproblemsolving.com/community/c6h1480726p8639339}{\textbf{2016 N2}} Let $\tau(n)$ be the number of positive divisors of $n$. Let $\tau_1(n)$ be the number of positive divisors of $n$ which have remainders $1$ when divided by $3$. Find all positive integral values of the fraction $\frac{\tau(10n)}{\tau_1(10n)}$.
\item \href{https://artofproblemsolving.com/community/c6h1268852p6622214}{\textbf{2015 N2}} Let $a$ and $b$ be positive integers such that $a! + b!$ divides $a!b!$. Prove that $3a \geqslant  2b + 2$.
\item \href{https://artofproblemsolving.com/community/c6h1113196p5083567}{\textbf{2014 N2}} Determine all pairs $(x, y)$ of positive integers such that\[\sqrt[3]{7x^2-13xy+7y^2}=|x-y|+1.\]
\item \boxed{\href{https://artofproblemsolving.com/community/c6h1181540p5720240}{\textbf{2013 N2}}}	Assume that $k$ and $n$ are two positive integers. Prove that there exist positive integers $m_1 , \dots , m_k$ such that\[1+\frac{2^k-1}{n}=\left(1+\frac1{m_1}\right)\dots \left(1+\frac1{m_k}\right).\]
\item \href{https://artofproblemsolving.com/community/c6h546188p3160603}{\textbf{2012 N2}} Find all triples $(x,y,z)$ of positive integers such that $x \leqslant y \leqslant z$ and
\[x^3(y^3+z^3)=2012(xyz+2).\]

\item \href{https://artofproblemsolving.com/community/c6h488540p2737650}{\textbf{2011 N2}} Consider a polynomial $P(x) =  (x+d_1)(x+d_2)\dots (x+d_9)$ where $d_1, d_2, \dots , d_9$ are nine distinct integers. Prove that there exists an integer $N,$ such that for all integers $x \geqslant  N$ the number $P(x)$ is divisible by a prime number greater than $20$.

\item \href{https://artofproblemsolving.com/community/c6h418640p2361999}{\textbf{2010 N2}} Find all pairs $(m,n)$ of nonnegative integers for which\[m^2 + 2 \cdot 3^n = m\left(2^{n+1} - 1\right).\]
\item \href{https://artofproblemsolving.com/community/c6h355796p1932941}{\textbf{2009 N2}} A positive integer $N$ is called balanced, if $N=1$ or if $N$ can be written as a product of an even number of not necessarily distinct primes. Given positive integers $a$ and $b$, consider the polynomial $P$ defined by $P(x)=(x+a)(x+b)$.

\begin{enumerate}[(a)]
    \item Prove that there exist distinct positive integers $a$ and $b$ such that all the number $P(1)$, $P(2)$,$\dots$, $P(50)$ are balanced.
    \item Prove that if $P(n)$ is balanced for all positive integers $n$, then $a=b$.
\end{enumerate}

\item \href{https://artofproblemsolving.com/community/c6h287870p1555929}{\textbf{2008 N2}} Let $ a_1$, $ a_2$, $ \dots$, $ a_n$ be distinct positive integers, $ n\geqslant  3$. Prove that there exist distinct indices $ i$ and $ j$ such that $ a_i + a_j$ does not divide any of the numbers $ 3a_1$, $ 3a_2$, $ \dots$, $ 3a_n$.
\item \href{https://artofproblemsolving.com/community/c6h214712p1187198}{\textbf{2007 N2}} 	Let $b,n > 1$ be integers. Suppose that for each $k > 1$ there exists an integer $a_k$ such that $b - a^n_k$ is divisible by $k$. Prove that $b = A^n$ for some integer $A$.
\item \href{https://artofproblemsolving.com/community/c6h155715p875055}{\textbf{2006 N2}} For $ x \in (0, 1)$ let $ y \in (0, 1)$ be the number whose $ n$-th digit after the decimal point is the $ 2^{n}$-th digit after the decimal point of $ x$. Show that if $ x$ is rational then so is $ y$.
\end{itemize}
\newpage

\section*{N3 Problems} \label{N3}
\begin{itemize}[leftmargin=*]


\item \href{https://www.imo-official.org/assets/documents/problems/2026/IMO2026SL.pdf}{\textbf{2026 N3}} TO BE REVEALED!

\item \boxed{\href{https://artofproblemsolving.com/community/c6h3611033p35347364}{\textbf{2025 N3}}} A proper divisor of a positive integer $N$ is a positive divisor of $N$ other than $N$ itself.

The infinite sequence $a_1, a_2, \dots$ consists of positive integers, each of which has at least three proper divisors. For each $n\geqslant  1$, the integer $a_{n+1}$ is the sum of the three largest proper divisors of $a_n$.

Determine all possible values of $a_1$.

\item \href{https://artofproblemsolving.com/community/c6h3610437p35340905}{\textbf{2024 N3}} Determine all sequences $a_1, a_2, \dots$ of positive integers such that for any pair of positive integers $m\leqslant n$, the arithmetic and geometric means
\[ \frac{a_m + a_{m+1} + \dots + a_n}{n-m+1}\quad\text{and}\quad (a_ma_{m+1}\dots a_n)^{\frac{1}{n-m+1}}\]are both integers.	

\item \href{https://artofproblemsolving.com/community/c6h3359721p31218370}{\textbf{2023 N3}}  For positive integers $n$ and $k \geqslant  2$, define $E_k(n)$ as the greatest exponent $r$ such that $k^r$ divides $n!$. Prove that there are infinitely many $n$ such that $E_{10}(n) > E_9(n)$ and infinitely many $m$ such that $E_{10}(m) < E_9(m)$.	

\item \href{https://artofproblemsolving.com/community/c6h3107336p28104282}{\textbf{2022 N3}} Let $a > 1$ be a positive integer and $d > 1$ be a positive integer coprime to $a$. Let $x_1=1$, and for $k\geqslant  1$, define
$$x_{k+1} = \begin{cases}
x_k + d &\text{if } a \text{ does not divide } x_k \\
x_k/a & \text{if } a \text{ divides } x_k
\end{cases}$$Find, in terms of $a$ and $d$, the greatest positive integer $n$ for which there exists an index $k$ such that $x_k$ is divisible by $a^n$.
\item \href{https://artofproblemsolving.com/community/c6h2882565p25627623}{\textbf{2021 N3}} Find all positive integers $n$ with the following property: the $k$ positive divisors of $n$ have a permutation $(d_1,d_2,\dots,d_k)$ such that for $i=1,2,\dots,k$, the number $d_1+d_2+\dots+d_i$ is a perfect square.
\item \boxed{\href{https://artofproblemsolving.com/community/c6h2278645p17821528}{\textbf{2020 N3}}} A deck of $n > 1$ cards is given. A positive integer is written on each card. The deck has the property that the arithmetic mean of the numbers on each pair of cards is also the geometric mean of the numbers on some collection of one or more cards.
For which $n$ does it follow that the numbers on the cards are all equal?
\item \href{https://artofproblemsolving.com/community/c6h2279036p17829189}{\textbf{2019 N3}} We say that a set $S$ of integers is rootiful if, for any positive integer $n$ and any $a_0, a_1, \dots, a_n \in S$, all integer roots of the polynomial $a_0+a_1x+\dots+a_nx^n$ are also in $S$. Find all rootiful sets of integers that contain all numbers of the form $2^a - 2^b$ for positive integers $a$ and $b$.
\item \href{https://artofproblemsolving.com/community/c6h1876764p12752823}{\textbf{2018 N3}} 	Define the sequence $a_0,a_1,a_2,\hdots$ by $a_n=2^n+2^{\lfloor n/2\rfloor}$. Prove that there are infinitely many terms of the sequence which can be expressed as a sum of (two or more) distinct terms of the sequence, as well as infinitely many of those which cannot be expressed in such a way.
\newpage
\item \href{https://artofproblemsolving.com/community/c6h1671289p10632346}{\textbf{2017 N3}} 	Determine all integers $ n\geqslant  2$ having the following property: for any integers $a_1,a_2,\dots, a_n$ whose sum is not divisible by $n$, there exists an index $1 \leqslant i \leqslant n$ such that none of the numbers$$a_i,a_i+a_{i+1},\dots,a_i+a_{i+1}+\dots+a_{i+n-1}$$is divisible by $n$. Here, we let $a_i=a_{i-n}$ when $i >n$.

\item \boxed{\href{https://artofproblemsolving.com/community/c6h1270992p6642559}{\textbf{2016 N3}}} A set of positive integers is called fragrant if it contains at least two elements and each of its elements has a prime factor in common with at least one of the other elements. Let $P(n)=n^2+n+1$. What is the least possible positive integer value of $b$ such that there exists a non-negative integer $a$ for which the set$$\{P(a+1),P(a+2),\dots,P(a+b)\}$$is fragrant?
\item \href{https://artofproblemsolving.com/community/c6h1268874p6622379}{\textbf{2015 N3}} Let $m$ and $n$ be positive integers such that $m>n$. Define $x_k=\frac{m+k}{n+k}$ for $k=1,2,\dots,n+1$. Prove that if all the numbers $x_1,x_2,\dots,x_{n+1}$ are integers, then $x_1x_2\dots x_{n+1}-1$ is divisible by an odd prime.
\item \boxed{\href{https://artofproblemsolving.com/community/c6h597093p3543144}{\textbf{2014 N3}}} For each positive integer $n$, the Bank of Cape Town issues coins of denomination $\frac1n$. Given a finite collection of such coins (of not necessarily different denominations) with total value at most most $99+\frac12$, prove that it is possible to split this collection into $100$ or fewer groups, such that each group has total value at most $1$.
\item \href{https://artofproblemsolving.com/community/c6h597245p3544101}{\textbf{2013 N3}} Prove that there exist infinitely many positive integers $n$ such that the largest prime divisor of $n^4 + n^2 + 1$ is equal to the largest prime divisor of $(n+1)^4 + (n+1)^2 +1$.
\item \href{https://artofproblemsolving.com/community/c6h545764p3156840}{\textbf{2012 N3}} Determine all integers $m \geqslant  2$ such that every $n$ with $\frac{m}{3} \leqslant n \leqslant \frac{m}{2}$ divides the binomial coefficient $\binom{n}{m-2n}$.
\item \href{https://artofproblemsolving.com/community/c6h488541p2737651}{\textbf{2011 N3}} 	Let $n \geqslant  1$ be an odd integer. Determine all functions $f$ from the set of integers to itself, such that for all integers $x$ and $y$ the difference $f(x)-f(y)$ divides $x^n-y^n.$
\item \href{https://artofproblemsolving.com/community/c6h418641p2362006}{\textbf{2010 N3}} 	Find the smallest number $n$ such that there exist polynomials $f_1, f_2, \dots , f_n$ with rational coefficients satisfying\[x^2+7 = f_1\left(x\right)^2 + f_2\left(x\right)^2 + \dots + f_n\left(x\right)^2.\]
\item \href{https://artofproblemsolving.com/community/c6h355797p1932942}{\textbf{2009 N3}} 	Let $f$ be a non-constant function from the set of positive integers into the set of positive integer, such that $a-b$ divides $f(a)-f(b)$ for all distinct positive integers $a$, $b$. Prove that there exist infinitely many primes $p$ such that $p$ divides $f(c)$ for some positive integer $c$.
\item \href{https://artofproblemsolving.com/community/c6h287871p1555931}{\textbf{2008 N3}} 	Let $ a_0$, $ a_1$, $ a_2$, $ \dots$ be a sequence of positive integers such that the greatest common divisor of any two consecutive terms is greater than the preceding term; in symbols, $ \gcd (a_i, a_{i + 1}) > a_{i - 1}$. 

Prove that $ a_n\geqslant  2^n$ for all $ n\geqslant  0$.
\item \href{https://artofproblemsolving.com/community/c6h214713p1187204}{\textbf{2007 N3}} Let $ X$ be a set of 10,000 integers, none of them is divisible by 47. Prove that there exists a 2007-element subset $ Y$ of $ X$ such that $ a - b + c - d + e$ is not divisible by 47 for any $ a,b,c,d,e \in Y.$
\item \href{https://artofproblemsolving.com/community/c6h130811p741367}{\textbf{2006 N3}} 	We define a sequence $ \left(a_{1},a_{2},a_{3},\dots \right)$ by
\[ a_{n} = \frac {1}{n}\left(\left\lfloor\frac {n}{1}\right\rfloor + \left\lfloor\frac {n}{2}\right\rfloor + \dots + \left\lfloor\frac {n}{n}\right\rfloor\right),
\]where $\lfloor x\rfloor$ denotes the integer part of $x$.
\begin{enumerate}[(a)]
    \item Prove that $a_{n+1}>a_n$ infinitely often. 
    \item Prove that $a_{n+1}<a_n$ infinitely often.
\end{enumerate}

\end{itemize}


















\newpage






\newpage
\newpage
\section*{N4 Problems} \label{N4}
\begin{itemize}[leftmargin=*]


\item \href{https://www.imo-official.org/assets/documents/problems/2026/IMO2026SL.pdf}{\textbf{2026 N4}} TO BE REVEALED!


\item \href{https://artofproblemsolving.com/community/c6h3866861p38676953}{\textbf{2025 N4}} Let $\mathbb{Z}_{>0}$ be a set of positive integers. Determine all functions $f: \mathbb{Z}_{>0} \rightarrow \mathbb{Z}$ such that for every $n\in \mathbb{Z}_{>0}$, and every positive divisor $d$ of $n$, there exists a positive divisor $e$ of $n$ such that $n$ divides $d+f(e)$.

\item \boxed{\href{https://artofproblemsolving.com/community/c6h3358926p31205957}{\textbf{2024 N4}}} Determine all pairs $(a,b)$ of positive integers for which there exist positive integers $g$ and $N$ such that
$$\gcd (a^n+b,b^n+a)=g$$holds for all integers $n\geqslant  N.$ (Note that $\gcd(x, y)$ denotes the greatest common divisor of integers $x$ and $y.$)

\item \href{https://artofproblemsolving.com/community/c6h3359730p31218384}{\textbf{2023 N4}} Let $a_1, \dots, a_n, b_1, \dots, b_n$ be $2n$ positive integers such that the $n+1$ products
\[a_1 a_2 a_3 \dots a_n \; , \; b_1 a_2 a_3 \dots a_n \; , \; b_1 b_2 a_3 \dots a_n \; , \; \dots \; , \; b_1 b_2 b_3 \dots b_n\]form a strictly increasing arithmetic progression in that order. Determine the smallest possible integer that could be the common difference of such an arithmetic progression.	

    \item  \boxed{\href{https://artofproblemsolving.com/community/c6h2883217p25635158}{\textbf{2022 N4}}} Find all triples $(a,b,p)$ of positive integers with $p$ prime and\[ a^p=b!+p.\] 
\item  \href{https://artofproblemsolving.com/community/c6h2882562p25627605}{\textbf{2021 N4}} Let $r>1$ be a rational number. Alice plays a solitaire game on a number line. Initially there is a red bead at $0$ and a blue bead at $1$. In a move, Alice chooses one of the beads and an integer $k \in \mathbb{Z}$. If the chosen bead is at $x$, and the other bead is at $y$, then the bead at $x$ is moved to the point $x'$ satisfying $x'-y=r^k(x-y)$. Find all $r$ for which Alice can move the red bead to $1$ in at most $2021$ moves.
\item  \href{https://artofproblemsolving.com/community/c6h2625856p22698019}{\textbf{2020 N4}} For any odd prime $p$ and any integer $n,$ let $d_p (n) \in \{ 0,1, \dots, p-1 \}$ denote the remainder when $n$ is divided by $p.$ We say that $(a_0, a_1, a_2, \dots)$ is a p-sequence, if $a_0$ is a positive integer coprime to $p,$ and $a_{n+1} =a_n + d_p (a_n)$ for $n \geqslant  0.$ 
\begin{enumerate}[(a)]
    \item Do there exist infinitely many primes $p$ for which there exist $p$-sequences $(a_0, a_1, a_2, \dots)$ and $(b_0, b_1, b_2, \dots)$ such that $a_n >b_n$ for infinitely many $n,$ and $b_n > a_n$ for infinitely many $n?$ 
    \item Do there exist infinitely many primes $p$ for which there exist $p$-sequences $(a_0, a_1, a_2, \dots)$ and $(b_0, b_1, b_2, \dots)$ such that $a_0 <b_0,$ but $a_n >b_n$ for all $n \geqslant  1?$
\end{enumerate}

\item  \href{https://artofproblemsolving.com/community/c6h2279056p17829445}{\textbf{2019 N4}} Find all functions $f:\mathbb Z_{>0}\to \mathbb Z_{>0}$ such that $a+f(b)$ divides $a^2+bf(a)$ for all positive integers $a$ and $b$ with $a+b>2019$.
\item  \boxed{\href{https://artofproblemsolving.com/community/c6h1671291p10632353}{\textbf{2018 N4}}} Let $a_1$, $a_2$, $\dots$ be an infinite sequence of positive integers. Suppose that there is an integer $N > 1$ such that, for each $n \geqslant  N$, the number $$\frac{a_1}{a_2} + \frac{a_2}{a_3} + \dots + \frac{a_{n-1}}{a_n} + \frac{a_n}{a_1}$$is an integer. Prove that there is a positive integer $M$ such that $a_m = a_{m+1}$ for all $m \geqslant  M$.
\newpage
\item  \href{https://artofproblemsolving.com/community/c6h1671331p10632689}{\textbf{2017 N4}} Call a rational number short if it has finitely many digits in its decimal expansion. For a positive integer $m$, we say that a positive integer $t$ is $m-$tastic if there exists a number $c\in \{1,2,3,\dots ,2017\}$ such that $\dfrac{10^t-1}{c\cdot m}$ is short, and such that $\dfrac{10^k-1}{c\cdot m}$ is not short for any $1\leqslant k<t$. Let $S(m)$ be the set of $m-$tastic numbers. Consider $S(m)$ for $m=1,2,\dots{}.$ What is the maximum number of elements in $S(m)$?
\item  \href{https://artofproblemsolving.com/community/c6h1480731p8639358}{\textbf{2016 N4}} Let $n, m, k$ and $l$ be positive integers with $n \neq 1$ such that $n^k + mn^l + 1$ divides $n^{k+l} - 1$. Prove that 
\begin{enumerate}[(a)]
    \item $m = 1$ and $l = 2k$; or 
    \item $l|k$ and $m = \dfrac{n^{k-l}-1}{n^l-1}$.
\end{enumerate}

\item \href{https://artofproblemsolving.com/community/c6h1268845p6622169}{\textbf{2015 N4}} Suppose that $a_0, a_1, \dots $ and $b_0, b_1, \dots$ are two sequences of positive integers such that $a_0, b_0 \geqslant  2$ and\[ a_{n+1} = \gcd{(a_n, b_n)} + 1, \qquad b_{n+1} = \operatorname{lcm}{(a_n, b_n)} - 1. \]Show that the sequence $a_n$ is eventually periodic; in other words, there exist integers $N \geqslant  0$ and $t > 0$ such that $a_{n+t} = a_n$ for all $n \geqslant  N$.

\item \href{https://artofproblemsolving.com/community/c6h1113198p5083569}{\textbf{2014 N4}} 	Let $n > 1$ be a given integer. Prove that infinitely many terms of the sequence $(a_k )_{k\geqslant  1}$, defined by\[a_k=\left\lfloor\frac{n^k}{k}\right\rfloor,\]are odd.
\item \href{https://artofproblemsolving.com/community/c6h597246p3544103}{\textbf{2013 N4}} Determine whether there exists an infinite sequence of nonzero digits $a_1 , a_2 , a_3 , \dots $ and a positive integer $N$ such that for every integer $k > N$, the number $\overline{a_k a_{k-1}\dots a_1 }$ is a perfect square.
\item \href{https://artofproblemsolving.com/community/c6h546190p3160606}{\textbf{2012 N4}} An integer $a$ is called friendly if the equation $(m^2+n)(n^2+m)=a(m-n)^3$ has a solution over the positive integers.
\begin{enumerate}[(a)]
    \item Prove that there are at least $500$ friendly integers in the set $\{ 1,2,\dots ,2012\}$.
    \item Decide whether $a=2$ is friendly.
\end{enumerate}

\item \href{https://artofproblemsolving.com/community/c6h488542p2737653}{\textbf{2011 N4}} For each positive integer $k,$ let $t(k)$ be the largest odd divisor of $k.$ Determine all positive integers $a$ for which there exists a positive integer $n,$ such that all the differences

\[t(n+a)-t(n), t(n+a+1)-t(n+1), \dots, t(n+2a-1)-t(n+a-1)\]are divisible by 4.
\item \href{https://artofproblemsolving.com/community/c6h418642p2362008}{\textbf{2010 N4}} Let $a, b$ be integers, and let $P(x) = ax^3+bx.$ For any positive integer $n$ we say that the pair $(a,b)$ is $n$-good if $n | P(m)-P(k)$ implies $n | m - k$ for all integers $m, k.$ We say that $(a,b)$ is $very \ good$ if $(a,b)$ is $n$-good for infinitely many positive integers $n$.

\begin{enumerate}[(a)]
    \item Find a pair $(a,b)$ which is 51-good, but not very good.
    \item Show that all 2010-good pairs are very good.
\end{enumerate}

\newpage

\item \href{https://artofproblemsolving.com/community/c6h355798p1932944}{\textbf{2009 N4}} Find all positive integers $n$ such that there exists a sequence of positive integers $a_1$, $a_2$,$\dots$, $a_n$ satisfying:\[a_{k+1}=\frac{a_k^2+1}{a_{k-1}+1}-1\]for every $k$ with $2\leqslant k\leqslant n-1$.
\item \href{https://artofproblemsolving.com/community/c6h287872p1555932}{\textbf{2008 N4}} 	Let $ n$ be a positive integer. Show that the numbers
\[ \binom{2^n - 1}{0},\; \binom{2^n - 1}{1},\; \binom{2^n - 1}{2},\; \dots,\; \binom{2^n - 1}{2^{n - 1} - 1}\]
are congruent modulo $ 2^n$ to $ 1$, $ 3$, $ 5$, $ \dots$, $ 2^n - 1$ in some order.
\item \href{https://artofproblemsolving.com/community/c6h214716p1187214}{\textbf{2007 N4}} For every integer $ k \geqslant  2,$ prove that $ 2^{3k}$ divides the number
\[ \binom{2^{k + 1}}{2^{k}} - \binom{2^{k}}{2^{k - 1}}
\]
but $ 2^{3k + 1}$ does not.
\item \boxed{\href{https://artofproblemsolving.com/community/c6h101487p572821}{\textbf{2006 N4}}} Let $P(x)$ be a polynomial of degree $n > 1$ with integer coefficients and let $k$ be a positive integer. Consider the polynomial $Q(x) = P(P(\dots P(P(x)) \dots ))$, where $P$ occurs $k$ times. Prove that there are at most $n$ integers $t$ such that $Q(t) = t$.
\end{itemize}
\newpage

\newpage
\section*{N5 Problems}\label{N5}
\begin{itemize}[leftmargin=*]

\item \href{https://www.imo-official.org/assets/documents/problems/2026/IMO2026SL.pdf}{\textbf{2026 N5}} TO BE REVEALED!


\item \href{https://artofproblemsolving.com/community/c6h3866855p38676941}{\textbf{2025 N5}} The sequence of triples of positive integers $(a_0, b_0, c_0) , (a_1, b_1, c_1), \dots, (a_{2025}, b_{2025}, c_{2025})$ satisfies
    \[ (a_{i+1},b_{i+1},c_{i+1}) = (a_i+\gcd(b_i,c_i),b_i+\gcd(c_i,a_i) ,c_i+\gcd(a_i,b_i)  )    \]
    for $0\leqslant i \leqslant 2024$.

    Determine the smallest possible value of $a_{2025}$. 

\item \href{https://artofproblemsolving.com/community/c6h3610445p35340917}{\textbf{2024 N5}} Let $\mathcal{S}$ be a finite nonempty set of prime numbers. Let $1 = b_1 < b_2 < \dots$ be the sequence of all positive integers whose prime divisors all belong to $\mathcal{S}$. Prove that, for all but finitely many positive integers $n$, there exist positive integers $a_1, a_2, \dots, a_n$ such that
\[
    \frac{a_1}{b_1} + \frac{a_2}{b_2} + \dots + \frac{a_n}{b_n} = \left\lceil \frac{1}{b_1} + \frac{1}{b_2} + \dots + \frac{1}{b_n} \right\rceil.
\]	

\item \href{https://artofproblemsolving.com/community/c6h3359746p31218469}{\textbf{2023 N5}} Let $a_1<a_2<a_3<\dots$ be positive integers such that $a_{k+1}$ divides $2(a_1+a_2+\dots+a_k)$ for every $k\geqslant  1$. Suppose that for infinitely many primes $p$, there exists $k$ such that $p$ divides $a_k$. Prove that for every positive integer $n$, there exists $k$ such that $n$ divides $a_k$.	

    \item \href{https://artofproblemsolving.com/community/c6h3107337p28104286}{\textbf{2022 N5}} For each $1\leqslant i\leqslant 9$ and $T\in\mathbb N$, define $d_i(T)$ to be the total number of times the digit $i$ appears when all the multiples of $1829$ between $1$ and $T$ inclusive are written out in base $10$. Show that there are infinitely many $T\in\mathbb N$ such that there are precisely two distinct values among $d_1(T)$, $d_2(T)$, $\dots$, $d_9(T)$.
\item \href{https://artofproblemsolving.com/community/c6h2882536p25627492}{\textbf{2021 N5}} Show that $n!=a^{n-1}+b^{n-1}+c^{n-1}$ has only finitely many solutions in positive integers.
\item \href{https://artofproblemsolving.com/community/c6h2625925p22698553}{\textbf{2020 N5}} Determine all functions $f$ defined on the set of all positive integers and taking non-negative integer values, satisfying the three conditions:
\begin{enumerate}[(a)]
\item $f(n) \neq 0$ for at least one $n$;
\item $f(x y)=f(x)+f(y)$ for every positive integers $x$ and $y$;
\item there are infinitely many positive integers $n$ such that $f(k)=f(n-k)$ for all $k<n$
\end{enumerate}
\item \href{https://artofproblemsolving.com/community/c6h2279015p17828930}{\textbf{2019 N5}} 	Let $a$ be a positive integer. We say that a positive integer $b$ is $a$-good if $\tbinom{an}{b}-1$ is divisible by $an+1$ for all positive integers $n$ with $an \geqslant  b$. Suppose $b$ is a positive integer such that $b$ is $a$-good, but $b+2$ is not $a$-good. Prove that $b+1$ is prime.
\item \href{https://artofproblemsolving.com/community/c6h1876761p12752815}{\textbf{2018 N5}} Four positive integers $x,y,z$ and $t$ satisfy the relations
\[ xy - zt = x + y = z + t\]Is it possible that both $xy$ and $zt$ are perfect squares?
\item \href{https://artofproblemsolving.com/community/c6h1671301p10632380}{\textbf{2017 N5}} Find all pairs $(p,q)$ of prime numbers which $p>q$ and
$$\frac{(p+q)^{p+q}(p-q)^{p-q}-1}{(p+q)^{p-q}(p-q)^{p+q}-1}$$is an integer.
\newpage
\item \href{https://artofproblemsolving.com/community/c6h1480706p8639282}{\textbf{2016 N5}} Let $a$ be a positive integer which is not a perfect square, and consider the equation\[k = \frac{x^2-a}{x^2-y^2}.\]Let $A$ be the set of positive integers $k$ for which the equation admits a solution in $\mathbb Z^2$ with $x>\sqrt{a}$, and let $B$ be the set of positive integers for which the equation admits a solution in $\mathbb Z^2$ with $0\leqslant x<\sqrt{a}$. Show that $A=B$.
\item \boxed{\href{https://artofproblemsolving.com/community/c6h1112743p5079630}{\textbf{2015 N5}}} Find all positive integers $(a,b,c)$ such that
$$ab-c,\quad bc-a,\quad ca-b$$are all powers of $2$.
\item \href{https://artofproblemsolving.com/community/c6h1113199p5083571}{\textbf{2014 N5}} Find all triples $(p, x, y)$ consisting of a prime number $p$ and two positive integers $x$ and $y$ such that $x^{p -1} + y$ and $x + y^ {p -1}$ are both powers of $p$.
\item \href{https://artofproblemsolving.com/community/c6h597247p3544105}{\textbf{2013 N5}} 	Fix an integer $k>2$. Two players, called Ana and Banana, play the following game of numbers. Initially, some integer $n \geqslant  k$ gets written on the blackboard. Then they take moves in turn, with Ana beginning. A player making a move erases the number $m$ just written on the blackboard and replaces it by some number $m'$ with $k \leqslant m' < m$ that is coprime to $m$. The first player who cannot move anymore loses. An integer $n \geqslant  k $ is called good if Banana has a winning strategy when the initial number is $n$, and bad otherwise. Consider two integers $n,n' \geqslant  k$ with the property that each prime number $p \leqslant k$ divides $n$ if and only if it divides $n'$. Prove that either both $n$ and $n'$ are good or both are bad.
\item \href{https://artofproblemsolving.com/community/c6h546191p3160608}{\textbf{2012 N5}} For a nonnegative integer $n$ define $\operatorname{rad}(n)=1$ if $n=0$ or $n=1$, and $\operatorname{rad}(n)=p_1p_2\dots p_k$ where $p_1<p_2<\dots <p_k$ are all prime factors of $n$. Find all polynomials $f(x)$ with nonnegative integer coefficients such that $\operatorname{rad}(f(n))$ divides $\operatorname{rad}(f(n^{\operatorname{rad}(n)}))$ for every nonnegative integer $n$.

\item \boxed{\href{https://artofproblemsolving.com/community/c6h418981p2365041}{\textbf{2011 N5}}} Let $f$ be a function from the set of integers to the set of positive integers. Suppose that, for any two integers $m$ and $n$, the difference $f(m) - f(n)$ is divisible by $f(m- n)$. Prove that, for all integers $m$ and $n$ with $f(m) \leqslant f(n)$, the number $f(n)$ is divisible by $f(m)$.

\item \boxed{\href{https://artofproblemsolving.com/community/c6h356076p1935854}{\textbf{2010 N5}}} Find all functions $g:\mathbb{N}\rightarrow\mathbb{N}$ such that\[\left(g(m)+n\right)\left(g(n)+m\right)\]is a perfect square for all $m,n\in\mathbb{N}$.

\item \href{https://artofproblemsolving.com/community/c6h355799p1932945}{\textbf{2009 N5}} Let $P(x)$ be a non-constant polynomial with integer coefficients. Prove that there is no function $T$ from the set of integers into the set of integers such that the number of integers $x$ with $T^n(x)=x$ is equal to $P(n)$ for every $n\geqslant  1$, where $T^n$ denotes the $n$-fold application of $T$.

\item \href{https://artofproblemsolving.com/community/c6h287873p1555934}{\textbf{2008 N5}} For every $ n\in\mathbb{N}$ let $ d(n)$ denote the number of positive divisors of $ n$. Find all functions $ f: \mathbb{N}\to\mathbb{N}$ with the following properties: 
\begin{enumerate}[(a)]
\item $d\left(f(x)\right) = x$ for all $ x\in\mathbb{N}$.
\item $ f(xy)$ divides $ (x - 1)y^{xy - 1}f(x)$ for all $ x$, $ y\in\mathbb{N}$.
\end{enumerate}

\item \href{https://artofproblemsolving.com/community/c6h214717p1187222}{\textbf{2007 N5}} Find all surjective functions $ f: \mathbb{N} \to \mathbb{N}$ such that for every $ m,n \in \mathbb{N}$ and every prime $ p,$ the number $ f(m + n)$ is divisible by $ p$ if and only if $ f(m) + f(n)$ is divisible by $ p$.
\item \href{https://artofproblemsolving.com/community/c6h138053p780855}{\textbf{2006 N5}} 	Find all integer solutions of the equation\[\frac {x^{7} - 1}{x - 1} = y^{5} - 1.\]
\end{itemize}


\newpage
\section*{N6 Problems} \label{N6}
\begin{itemize}[leftmargin=*]

\item \href{https://www.imo-official.org/assets/documents/problems/2026/IMO2026SL.pdf}{\textbf{2026 N6}} TO BE REVEALED!


\item \href{https://artofproblemsolving.com/community/c6h3866880p38677010}{\textbf{2025 N6}} For positive integers $n$ and $k$, let $f_{k}(n)$ denote the remainder when $n$ is divided by $2^{k}-1.$ A positive integer $n$ is \textit{grouse} if
\[ f_{1}(n)\leqslant f_{2}(n)\leqslant f_{3}(n)\leqslant f_{4}(n)\leqslant \dots \]
\begin{enumerate}[leftmargin=*]
    \item[(a)]  Prove that there are infinitely many grouse numbers.
\item[(b)] Prove that there exists a positive integer $M$ such that there are at most $\dfrac{M}{2025^{2025}}$ grouse numbers less than $M$.
\end{enumerate}

\item \href{https://artofproblemsolving.com/community/c6h3610457p35340935}{\textbf{2024 N6}} Let $n$ be a positive integer. We say a polynomial $P$ with integer coefficients is $\emph{n-good}$ if there exists a polynomial $Q$ of degree $2$ with integer coefficients such that $Q(k)(P(k) + Q(k))$ is never divisible by $n$ for any integer $k$.

Determine all integers $n$ such that every polynomial with integer coefficients is an $n$-good polynomial.

\item \href{https://artofproblemsolving.com/community/c6h3359725p31218376}{\textbf{2023 N6}} A sequence of integers $a_0, a_1, \dots $ is called kawaii if $a_0 =0, a_1=1,$ and$$(a_{n+2}-3a_{n+1}+2a_n)(a_{n+2}-4a_{n+1}+3a_n)=0$$for all integers $n \geqslant  0$. An integer is called kawaii if it belongs to some kawaii sequence.
Suppose that two consecutive integers $m$ and $m+1$ are both kawaii (not necessarily belonging to the same kawaii sequence). Prove that $m$ is divisible by $3,$ and that $m/3$ is also kawaii.	

    \item \href{https://artofproblemsolving.com/community/c6h3107366p28104497}{\textbf{2022 N6}} Let $Q$ be a set of prime numbers, not necessarily finite. For a positive integer $n$ consider its prime factorization: define $p(n)$ to be the sum of all the exponents and $q(n)$ to be the sum of the exponents corresponding only to primes in $Q$. A positive integer $n$ is called special if $p(n)+p(n+1)$ and $q(n)+q(n+1)$ are both even integers. Prove that there is a constant $c>0$ independent of the set $Q$ such that for any positive integer $N>100$, the number of special integers in $[1,N]$ is at least $cN$.

(For example, if $Q=\{3,7\}$, then $p(42)=3$, $q(42)=2$, $p(63)=3$, $q(63)=3$, $p(2022)=3$, $q(2022)=1$.)
\item \href{https://artofproblemsolving.com/community/c6h2882554p25627555}{\textbf{2021 N6}} Determine all integers $n\geqslant  2$ with the following property: every $n$ pairwise distinct integers whose sum is not divisible by $n$ can be arranged in some order $a_1,a_2,\dots, a_n$ so that $n$ divides $1\cdot a_1+2\cdot a_2+\dots+n\cdot a_n.$
\item \href{https://artofproblemsolving.com/community/c6h2625861p22698064}{\textbf{2020 N6}} For a positive integer $n$, let $d(n)$ be the number of positive divisors of $n$, and let $\varphi(n)$ be the number of positive integers not exceeding $n$ which are coprime to $n$. Does there exist a constant $C$ such that
$$ \frac {\varphi ( d(n))}{d(\varphi(n))}\leqslant C$$for all $n\geqslant  1$.
\item \href{https://artofproblemsolving.com/community/c6h2279051p17829382}{\textbf{2019 N6}} 	Let $H = \{ \lfloor i\sqrt{2}\rfloor : i \in \mathbb Z_{>0}\} = \{1,2,4,5,7,\dots \}$ and let $n$ be a positive integer. Prove that there exists a constant $C$ such that, if $A\subseteq \{1,2,\dots, n\}$ satisfies $|A| \geqslant  C\sqrt{n}$, then there exist $a,b\in A$ such that $a-b\in H$. ($\mathbb Z_{>0}$ is the set of positive integers, and $\lfloor z\rfloor$ denotes the greatest integer less than or equal to $z$.)
\item \href{https://artofproblemsolving.com/community/c6h1876765p12752830}{\textbf{2018 N6}} Let $f : \{ 1, 2, 3, \dots \} \to \{ 2, 3, \dots \}$ be a function such that $f(m + n) | f(m) + f(n) $ for all pairs $m,n$ of positive integers. Prove that there exists a positive integer $c > 1$ which divides all values of $f$.
\item \href{https://artofproblemsolving.com/community/c6h1671292p10632354}{\textbf{2017 N6}} Find the smallest positive integer $n$ or show no such $n$ exists, with the following property: there are infinitely many distinct $n$-tuples of positive rational numbers $(a_1, a_2, \dots, a_n)$ such that both
$$a_1+a_2+\dots +a_n \quad \text{and} \quad \frac{1}{a_1} + \frac{1}{a_2} + \dots + \frac{1}{a_n}$$are integers.
\item \href{https://artofproblemsolving.com/community/c6h1480691p8639255}{\textbf{2016 N6}} Denote by $\mathbb{N}$ the set of all positive integers. Find all functions $f:\mathbb{N}\rightarrow \mathbb{N}$ such that for all positive integers $m$ and $n$, the integer $f(m)+f(n)-mn$ is nonzero and divides $mf(m)+nf(n)$.
\item \href{https://artofproblemsolving.com/community/c6h1268916p6622847}{\textbf{2015 N6}} Let $\mathbb{Z}_{>0}$ denote the set of positive integers. Consider a function $f: \mathbb{Z}_{>0} \to \mathbb{Z}_{>0}$. For any $m, n \in \mathbb{Z}_{>0}$ we write $f^n(m) = \underbrace{f(f(\dots f}_{n}(m)\dots))$. Suppose that $f$ has the following two properties:
\begin{enumerate}[(a)]
\item if $m, n \in \mathbb{Z}_{>0}$, then $\frac{f^n(m) - m}{n} \in \mathbb{Z}_{>0}$;
\item The set $\mathbb{Z}_{>0} \setminus \{f(n) \mid n\in \mathbb{Z}_{>0}\}$ is finite.
\end{enumerate}
Prove that the sequence $f(1) - 1, f(2) - 2, f(3) - 3, \dots$ is periodic.
\item \href{https://artofproblemsolving.com/community/c6h1113200p5083572}{\textbf{2014 N6}} Let $a_1 < a_2 <  \dots <a_n$ be pairwise coprime positive integers with $a_1$ being prime and $a_1 \geqslant  n + 2$. On the segment $I = [0, a_1 a_2  \dots a_n ]$ of the real line, mark all integers that are divisible by at least one of the numbers $a_1 ,   \dots , a_n$ . These points split $I$ into a number of smaller segments. Prove that the sum of the squares of the lengths of these segments is divisible by $a_1$.
\item \href{https://artofproblemsolving.com/community/c6h597248p3544109}{\textbf{2013 N6}} Determine all functions $f: \mathbb{Q} \rightarrow \mathbb{Z} $ satisfying
\[ f \left( \frac{f(x)+a} {b}\right) = f \left( \frac{x+a}{b} \right) \]
for all $x \in \mathbb{Q}$, $a \in \mathbb{Z}$, and $b \in \mathbb{Z}_{>0}$. (Here, $\mathbb{Z}_{>0}$ denotes the set of positive integers.)
\item \href{https://artofproblemsolving.com/community/c6h545765p3156844}{\textbf{2012 N6}} Let $x$ and $y$ be positive integers. If ${x^{2^n}}-1$ is divisible by $2^ny+1$ for every positive integer $n$, prove that $x=1$.
\item \href{https://artofproblemsolving.com/community/c6h488543p2737654}{\textbf{2011 N6}} Let $P(x)$ and $Q(x)$ be two polynomials with integer coefficients, such that no nonconstant polynomial with rational coefficients divides both $P(x)$ and $Q(x).$ Suppose that for every positive integer $n$ the integers $P(n)$ and $Q(n)$ are positive, and $2^{Q(n)}-1$ divides $3^{P(n)}-1.$ Prove that $Q(x)$ is a constant polynomial.
\item \href{https://artofproblemsolving.com/community/c6h418643p2362013}{\textbf{2010 N6}} The rows and columns of a $2^n \times 2^n$ table are numbered from $0$ to $2^{n}-1.$ The cells of the table have been coloured with the following property being satisfied: for each $0 \leqslant i,j \leqslant 2^n - 1,$ the $j$-th cell in the $i$-th row and the $(i+j)$-th cell in the $j$-th row have the same colour. (The indices of the cells in a row are considered modulo $2^n$.) Prove that the maximal possible number of colours is $2^n$.
\item \href{https://artofproblemsolving.com/community/c6h355800p1932946}{\textbf{2009 N6}} Let $k$ be a positive integer. Show that if there exists a sequence $a_0,a_1,\dots$ of integers satisfying the condition\[a_n=\frac{a_{n-1}+n^k}{n}\text{ for all } n\geqslant  1,\]then $k-2$ is divisible by $3$.
\item \boxed{\href{https://artofproblemsolving.com/community/c6h215221p1190546}{\textbf{2008 N6}}} Prove that there are infinitely many positive integers $ n$ such that $ n^{2} + 1$ has a prime divisor greater than $ 2n + \sqrt {2n}$.
\item \boxed{\href{https://artofproblemsolving.com/community/c6h214725p1187273}{\textbf{2007 N6}}} Let $ k$ be a positive integer. Prove that the number $ (4k^2 - 1)^2$ has a positive divisor of the form $ 8kn - 1$ if and only if $ k$ is even.
\item \href{https://artofproblemsolving.com/community/c6h155721p875072}{\textbf{2006 N6}} Let $ a > b > 1$ be relatively prime positive integers. Define the weight of an integer $ c$, denoted by $ w(c)$ to be the minimal possible value of $ |x| + |y|$ taken over all pairs of integers $ x$ and $ y$ such that\[ax + by = c.\]An integer $ c$ is called a local champion if $ w(c) \geqslant  w(c \pm a)$ and $ w(c) \geqslant  w(c \pm b)$. Find all local champions and determine their number.
\end{itemize}


\newpage
\section*{N7 Problems}  \label{N7}
\begin{itemize}[leftmargin=*]


\item \href{https://www.imo-official.org/assets/documents/problems/2026/IMO2026SL.pdf}{\textbf{2026 N7}} TO BE REVEALED!

\item \boxed{\href{https://artofproblemsolving.com/community/c6h3609789p35332016}{\textbf{2025 N7}}} Let $\mathbb{N}$ denote the set of positive integers. A function $f\colon\mathbb{N}\to\mathbb{N}$ is said to be bonza if
\[
f(a)~~\text{divides}~~b^a-f(b)^{f(a)}
\]for all positive integers $a$ and $b$.

Determine the smallest real constant $c$ such that $f(n)\leqslant cn$ for all bonza functions $f$ and all positive integers $n$.
\item \href{https://artofproblemsolving.com/community/c6h3610443p35340914}{\textbf{2024 N7}} Let $\mathbb{Z}_{>0}$ denote the set of positive integers. Let $f : \mathbb{Z}_{>0} \rightarrow \mathbb{Z}_{>0}$ be a function satisfying the following property: for $m,n \in \mathbb{Z}_{>0}$, the equation
\[
f(mn)^2 = f(m^2)f(f(n))f(mf(n))
\]holds if and only if $m$ and $n$ are coprime.

For each positive integer $n$, determine all the possible values of $f(n)$.	

\item \href{https://artofproblemsolving.com/community/c6h3359727p31218380}{\textbf{2023 N7}} Let $a,b,c,d$ be positive integers satisfying\[\frac{ab}{a+b}+\frac{cd}{c+d}=\frac{(a+b)(c+d)}{a+b+c+d}.\]Determine all possible values of $a+b+c+d$.	

    \item  \boxed{\href{https://artofproblemsolving.com/community/c6h2883213p25635143}{\textbf{2022 N7}}} Let $k$ be a positive integer and let $S$ be a finite set of odd prime numbers. Prove that there is at most one way (up to rotation and reflection) to place the elements of $S$ around the circle such that the product of any two neighbors is of the form $x^2+x+k$ for some positive integer $x$.
\item  \href{https://artofproblemsolving.com/community/c6h2882548p25627533}{\textbf{2021 N7}} 	Let $a_1,a_2,a_3,\dots$ be an infinite sequence of positive integers such that $a_{n+2m}$ divides $a_{n}+a_{n+m}$ for all positive integers $n$ and $m.$ Prove that this sequence is eventually periodic, i.e. there exist positive integers $N$ and $d$ such that $a_n=a_{n+d}$ for all $n>N.$
\item  \href{https://artofproblemsolving.com/community/c6h2625921p22698513}{\textbf{2020 N7}} Let $\mathcal{S}$ be a set consisting of $n \geqslant  3$ positive integers, none of which is a sum of two other distinct members of $\mathcal{S}$. Prove that the elements of $\mathcal{S}$ may be ordered as $a_1, a_2, \dots, a_n$ so that $a_i$ does not divide $a_{i - 1} + a_{i + 1}$ for all $i = 2, 3, \dots, n - 1$.
\item  \href{https://artofproblemsolving.com/community/c6h2279001p17828757}{\textbf{2019 N7}} Prove that there is a constant $c>0$ and infinitely many positive integers $n$ with the following property: there are infinitely many positive integers that cannot be expressed as the sum of fewer than $cn\log(n)$ pairwise coprime $n$th powers.
\item  \href{https://artofproblemsolving.com/community/c6h1876776p12752855}{\textbf{2018 N7}} Let $n \geqslant  2018$ be an integer, and let $a_1, a_2, \dots, a_n, b_1, b_2, \dots, b_n$ be pairwise distinct positive integers not exceeding $5n$. Suppose that the sequence
\[ \frac{a_1}{b_1}, \frac{a_2}{b_2}, \dots, \frac{a_n}{b_n} \]forms an arithmetic progression. Prove that the terms of the sequence are equal.
\newpage
\item  \boxed{\href{https://artofproblemsolving.com/community/c6h1480686p8639242}{\textbf{2017 N7}}} An ordered pair $(x, y)$ of integers is a primitive point if the greatest common divisor of $x$ and $y$ is $1$. Given a finite set $S$ of primitive points, prove that there exist a positive integer $n$ and integers $a_0, a_1, \dots , a_n$ such that, for each $(x, y)$ in $S$, we have:
$$a_0x^n + a_1x^{n-1} y + a_2x^{n-2}y^2 + \dots + a_{n-1}xy^{n-1} + a_ny^n = 1.$$
 
\item \boxed{\href{https://artofproblemsolving.com/community/c6h1270467p6637660}{\textbf{2016 N7}}}  Let $P=A_1A_2\dots A_k$ be a convex polygon in the plane. The vertices $A_1, A_2, \dots, A_k$ have integral coordinates and lie on a circle. Let $S$ be the area of $P$. An odd positive integer $n$ is given such that the squares of the side lengths of $P$ are integers divisible by $n$. Prove that $2S$ is an integer divisible by $n$.
\item  \href{https://artofproblemsolving.com/community/c6h1268919p6622862}{\textbf{2015 N7}} Let $\mathbb{Z}_{>0}$ denote the set of positive integers. For any positive integer $k$, a function $f: \mathbb{Z}_{>0} \to \mathbb{Z}_{>0}$ is called $k$-good if $\gcd(f(m) + n, f(n) + m) \leqslant k$ for all $m \neq n$. Find all $k$ such that there exists a $k$-good function.
\item  \href{https://artofproblemsolving.com/community/c6h1113202p5083576}{\textbf{2014 N7}} Let $c \geqslant  1$ be an integer. Define a sequence of positive integers by $a_1 = c$ and\[a_{n+1}=a_n^3-4c\cdot a_n^2+5c^2\cdot a_n+c\]for all $n\geqslant  1$. Prove that for each integer $n \geqslant  2$ there exists a prime number $p$ dividing $a_n$ but none of the numbers $a_1 , \dots , a_{n -1}$ .
\item  \href{https://artofproblemsolving.com/community/c6h597250p3544112}{\textbf{2013 N7}} Let $\nu$ be an irrational positive number, and let $m$ be a positive integer. A pair of $(a,b)$ of positive integers is called good if
\[a \left \lceil b\nu \right \rceil - b \left \lfloor a \nu \right \rfloor = m.\]A good pair $(a,b)$ is called excellent if neither of the pair $(a-b,b)$ and $(a,b-a)$ is good. Prove that the number of excellent pairs is equal to the sum of the positive divisors of $m$.
\item  \boxed{\href{https://artofproblemsolving.com/community/c6h488512p2737435}{\textbf{2012 N7}}} Find all positive integers $n$ for which there exist non-negative integers $a_1, a_2, \dots, a_n$ such that
\[
\frac{1}{2^{a_1}} + \frac{1}{2^{a_2}} + \dots + \frac{1}{2^{a_n}} = 
\frac{1}{3^{a_1}} + \frac{2}{3^{a_2}} + \dots + \frac{n}{3^{a_n}} = 1.
\]
\item  \href{https://artofproblemsolving.com/community/c6h488544p2737656}{\textbf{2011 N7}} Let $p$ be an odd prime number. For every integer $a,$ define the number $S_a = \sum^{p-1}_{j=1} \frac{a^j}{j}.$ Let $m,n \in \mathbb{Z},$ such that $S_3 + S_4 - 3S_2 = \frac{m}{n}.$ Prove that $p$ divides $m.$
\item  \href{https://artofproblemsolving.com/community/c6h355801p1932947}{\textbf{2009 N7}} Let $a$ and $b$ be distinct integers greater than $1$. Prove that there exists a positive integer $n$ such that $(a^n-1)(b^n-1)$ is not a perfect square.
\item  \href{https://artofproblemsolving.com/community/c6h214654p1186847}{\textbf{2007 N7}} For a prime $ p$ and a given integer $ n$ let $ \nu_p(n)$ denote the exponent of $ p$ in the prime factorisation of $ n!$. Given $ d \in \mathbb{N}$ and $ \{p_1,p_2,\dots,p_k\}$ a set of $ k$ primes, show that there are infinitely many positive integers $ n$ such that $ d\mid \nu_{p_i}(n)$ for all $ 1 \leqslant i \leqslant k$.
\item  \href{https://artofproblemsolving.com/community/c6h154177p867486}{\textbf{2006 N7}} For all positive integers $n$, show that there exists a positive integer $m$ such that $n$ divides $2^{m} + m$.
\end{itemize}

\newpage
\section*{N8 Problems}  \label{N8}
\begin{itemize}[leftmargin=*]


\item \boxed{\href{https://www.imo-official.org/assets/documents/problems/2026/IMO2026SL.pdf}{\textbf{2026 N8}}} Let $a_1, a_2, a_3, \ldots$ be an infinite sequence of positive integers greater than $1$. Suppose that for all positive integers $n$, the number $a_{n+1}$ is the smallest positive integer greater than $a_n$ such that $\text{gcd}(a_{n+1}, a_i)>1$ for every $i=1,2,\ldots, n$. Prove that there exist positive integers $T$ and $L$ such that
\[ a_{n+T}=a_n+L\]for every positive integer $n$.

(Note that $\text{gcd}(x,y)$ denotes the greatest common divisor of positive integers $x$ and $y$.)

 \item \href{https://artofproblemsolving.com/community/c6h3866854p38676940}{\textbf{2025 N8}} Prove that there are finitely many prime numbers $p$ such that
\begin{itemize}[leftmargin=*]
    \item[$\bullet$] $p-8$ is a perfect square, and

\item[$\bullet$] for all odd positive integers $k<\sqrt{p}$, the number $p-k^2$ has at most two distinct prime factors.
\end{itemize}


\item \href{https://artofproblemsolving.com/community/c6h3359735p31218397}{\textbf{2023 N8}} Determine all functions $f\colon\mathbb{Z}_{>0}\to\mathbb{Z}_{>0}$ such that, for all positive integers $a$ and $b$,
\[
f^{bf(a)}(a+1)=(a+1)f(b).
\]	

    \item  \href{https://artofproblemsolving.com/community/c6h3107338p28104291}{\textbf{2022 N8}} Prove that $5^n-3^n$ is not divisible by $2^n+65$  for any positive integer $n$.
\item  \href{https://artofproblemsolving.com/community/c6h2882533p25627467}{\textbf{2021 N8}} Find all positive integers $n$ for which there exists a polynomial $P(x) \in \mathbb{Z}[x]$ such that for every positive integer $m\geqslant  1$, the numbers $P^m(1), \dots, P^m(n)$ leave exactly $\lceil n/2^m\rceil$ distinct remainders when divided by $n$. (Here, $P^m$ means $P$ applied $m$ times.)
\item  \href{https://artofproblemsolving.com/community/c6h2278989p17828642}{\textbf{2019 N8}} Let $a$ and $b$ be two positive integers. Prove that the integer
\[a^2+\left\lceil\frac{4a^2}b\right\rceil\]is not a square.
\item  \href{https://artofproblemsolving.com/community/c6h1671295p10632363}{\textbf{2017 N8}} Let $p$ be an odd prime number and $\mathbb{Z}_{>0}$ be the set of positive integers. Suppose that a function $f:\mathbb{Z}_{>0}\times\mathbb{Z}_{>0}\to\{0,1\}$ satisfies the following properties:
\begin{enumerate}[(a)]
\item $f(1,1)=0$.
\item $f(a,b)+f(b,a)=1$ for any pair of relatively prime positive integers $(a,b)$ not both equal to 1;
\item $f(a+b,b)=f(a,b)$ for any pair of relatively prime positive integers $(a,b)$.
\end{enumerate}
Prove that
$$\sum_{n=1}^{p-1}f(n^2,p) \geqslant  \sqrt{2p}-2.$$
\item  \href{https://artofproblemsolving.com/community/c6h1480732p8639362}{\textbf{2016 N8}} 	Find all polynomials $P(x)$ of odd degree $d$ and with integer coefficients satisfying the following property: for each positive integer $n$, there exists $n$ positive integers $x_1, x_2, \dots, x_n$ such that $\dfrac12 < \dfrac{P(x_i)}{P(x_j)} < 2$ and $\dfrac{P(x_i)}{P(x_j)}$ is the $d$-th power of a rational number for every pair of indices $i$ and $j$ with $1 \leqslant i, j \leqslant n$.

\newpage

\item  \href{https://artofproblemsolving.com/community/c6h1268921p6622883}{\textbf{2015 N8}} 	For every positive integer $n$ with prime factorization $n = \prod_{i = 1}^{k} p_i^{\alpha_i}$, define
\[\mho(n) = \sum_{i: \; p_i > 10^{100}} \alpha_i.\]That is, $\mho(n)$ is the number of prime factors of $n$ greater than $10^{100}$, counted with multiplicity.

Find all strictly increasing functions $f: \mathbb{Z} \to \mathbb{Z}$ such that
\[\mho(f(a) - f(b)) \leqslant \mho(a - b) \quad \text{for all integers } a \text{ and } b \text{ with } a > b.\]

\item  \href{https://artofproblemsolving.com/community/c6h1113204p5083578}{\textbf{2014 N8}} For every real number $x$, let $||x||$ denote the distance between $x$ and the nearest integer.
Prove that for every pair $(a, b)$ of positive integers there exist an odd prime $p$ and a positive integer $k$ satisfying\[\displaystyle\left|\left|\frac{a}{p^k}\right|\right|+\left|\left|\frac{b}{p^k}\right|\right|+\left|\left|\frac{a+b}{p^k}\right|\right|=1.\]
\item  \href{https://artofproblemsolving.com/community/c6h545767p3156853}{\textbf{2012 N8}} Prove that for every prime $p>100$ and every integer $r$, there exist two integers $a$ and $b$ such that $p$ divides $a^2+b^5-r$.
\item  \href{https://artofproblemsolving.com/community/c6h488545p2737657}{\textbf{2011 N8}} Let $k \in \mathbb{Z}^+$ and set $n=2^k+1.$ Prove that $n$ is a prime number if and only if the following holds: there is a permutation $a_{1},\dots,a_{n-1}$ of the numbers $1,2, \dots, n-1$ and a sequence of integers $g_{1},\dots,g_{n-1},$ such that $n$ divides $g^{a_i}_i - a_{i+1}$ for every $i \in \{1,2,\dots,n-1\},$ where we set $a_n = a_1.$
\end{itemize}

\newpage
\section*{N9 Problems}  \label{N9}
\begin{itemize}[leftmargin=*]

\item \href{https://www.imo-official.org/assets/documents/problems/2026/IMO2026SL.pdf}{\textbf{2026 N9}} TO BE REVEALED!
\end{itemize}



\end{document}
